% Recovered from the previously completed PDF; no new mathematical content.
% Prose is editable; formula boxes retain original symbol positions.
% Build from repository root: xelatex -no-shell-escape -halt-on-error -output-directory=... print/fudan/NAME.tex
\documentclass[11pt,a4paper]{article}
\usepackage[margin=23mm,headheight=14pt,headsep=8.5mm,footskip=10mm]{geometry}
\usepackage{amsmath,amssymb}
\DeclareMathAlphabet{\mathsf}{OT1}{cmss}{m}{n}
\usepackage[no-math]{fontspec}
\usepackage{xeCJK}
\setmainfont{Latin Modern Roman}
\setCJKmainfont{FandolSong-Regular.otf}[Path=assets/fonts/fandol/,BoldFont=FandolSong-Bold.otf,ItalicFont=FandolSong-Regular.otf,Script=CustomDefault]
\xeCJKsetup{PunctStyle=quanjiao,CheckSingle=true}
\usepackage{setspace}\setstretch{1.12}
\setlength{\parindent}{2em}\setlength{\parskip}{3pt}
\emergencystretch=6em
\hfuzz=9pt
\usepackage{titlesec}
\titleformat{\section}{\centering\normalfont\Large\bfseries}{}{0pt}{}
\titlespacing*{\section}{0pt}{1.8\baselineskip}{1.0\baselineskip}
\usepackage{fancyhdr}\pagestyle{fancy}\fancyhf{}
\fancyhead[L]{\small 数学手记·高等代数}\fancyhead[R]{\small 15S01–15S07}
\fancyfoot[C]{\thepage}\renewcommand{\headrulewidth}{0.4pt}
\usepackage[hidelinks,unicode]{hyperref}
\hypersetup{pdftitle={15S01–15S07：题目、解答与小结},pdfauthor={Vincent Arqain}}
\newcommand{\RB}[4]{\raisebox{0pt}[#2bp][#3bp]{\hbox to #1bp{#4\hss}}}
\newcommand{\RG}[5]{\hbox to 0pt{\kern#1bp\raisebox{#2bp}[0pt][0pt]{\fontsize{#3bp}{#3bp}\csname font#4\endcsname\char#5}\hss}}
\newcommand{\RT}[4]{\hbox to 0pt{\kern#1bp\raisebox{#2bp}[0pt][0pt]{\fontsize{#3bp}{#3bp}\selectfont#4}\hss}}
\newcommand{\RS}[4]{\hbox to 0pt{\kern#1bp\raisebox{#2bp}[0pt][0pt]{\fontsize{#3bp}{#3bp}\selectfont\ensuremath{#4}}\hss}}
\newcommand{\RR}[4]{\hbox to 0pt{\kern#1bp\raisebox{#2bp}[0pt][0pt]{\rule{#3bp}{#4bp}}\hss}}
\newcommand{\fontmi}{\usefont{OML}{cmm}{m}{it}}
\newcommand{\fontsy}{\usefont{OMS}{cmsy}{m}{n}}
\newcommand{\fontex}{\usefont{OMX}{cmex}{m}{n}}
\newcommand{\fontrm}{\usefont{OT1}{cmr}{m}{n}}
\newcommand{\fontbf}{\usefont{OT1}{cmr}{bx}{n}}
\newcommand{\fontam}{\usefont{U}{msa}{m}{n}}
\newcommand{\fontbb}{\usefont{U}{msb}{m}{n}}
\widowpenalty=10000\clubpenalty=10000\raggedbottom
\begin{document}
\begin{center}{\LARGE\bfseries 15S01–15S07\par}\vspace{8pt}复旦大学高等代数习题集·题目与解答\end{center}
% Original page 1.
复旦大学高等代数习题集题目与解答

题目来源：谢启鸿编著，2023版，§1.3，第15–16页。

\section*{2015S01从逐点二选一到整体二选一}
题目设\allowbreak{}\hskip.1em \RB{31.994}{8.683}{3.226}{%
\RG{0.000}{0.001}{10.9091}{mi}{77}\RG{10.584}{-1.635}{7.9701}{mi}{110}\RG{16.220}{0.001}{10.9091}{rm}{40}\RG{20.462}{0.001}{10.9091}{bb}{82}\RG{28.341}{0.001}{10.9091}{rm}{41}%
}\hskip.1em\allowbreak{} 是\allowbreak{}\hskip.1em \RB{6.729}{5.323}{0.619}{%
\RG{0.000}{0.001}{10.9091}{mi}{110}%
}\hskip.1em\allowbreak{} 阶实方阵全体构成的实线性空间，\allowbreak{}\hskip.1em \RB{7.253}{5.323}{2.877}{%
\RG{0.000}{0.001}{10.9091}{mi}{39}%
}\hskip.1em\allowbreak{} 是其上的线性变换。对任意\allowbreak{}\hskip.1em \RB{32.571}{8.312}{2.604}{%
\RG{0.000}{0.001}{10.9091}{mi}{65}\RG{8.182}{0.001}{10.9091}{mi}{59}\RG{13.015}{0.001}{10.9091}{mi}{66}\RG{25.711}{0.001}{10.9091}{sy}{50}%
}\hskip.1em\allowbreak{}\allowbreak{}\hskip.1em \RB{31.993}{8.683}{3.226}{%
\RG{0.000}{0.001}{10.9091}{mi}{77}\RG{10.583}{-1.636}{7.9701}{mi}{110}\RG{16.219}{0.001}{10.9091}{rm}{40}\RG{20.462}{0.001}{10.9091}{bb}{82}\RG{28.340}{0.001}{10.9091}{rm}{41}%
}\hskip.1em\allowbreak{} ，或者

\allowbreak{}\hskip.1em \RB{98.131}{8.682}{3.227}{%
\RG{0.000}{0.000}{10.9091}{mi}{39}\RG{7.136}{0.000}{10.9091}{rm}{40}\RG{11.378}{0.000}{10.9091}{mi}{65}\RG{19.560}{0.000}{10.9091}{mi}{66}\RG{28.382}{0.000}{10.9091}{rm}{41}\RG{35.637}{0.000}{10.9091}{rm}{61}\RG{47.170}{0.000}{10.9091}{mi}{39}\RG{54.306}{0.000}{10.9091}{rm}{40}\RG{58.549}{0.000}{10.9091}{mi}{65}\RG{66.732}{0.000}{10.9091}{rm}{41}\RG{70.974}{0.000}{10.9091}{mi}{39}\RG{78.110}{0.000}{10.9091}{rm}{40}\RG{82.353}{0.000}{10.9091}{mi}{66}\RG{91.175}{0.000}{10.9091}{rm}{41}\RG{95.417}{0.000}{10.9091}{mi}{59}%
}\hskip.1em\allowbreak{} 或者\allowbreak{}\hskip.1em \RB{98.012}{8.682}{3.227}{%
\RG{0.000}{0.000}{10.9091}{mi}{39}\RG{7.137}{0.000}{10.9091}{rm}{40}\RG{11.379}{0.000}{10.9091}{mi}{65}\RG{19.561}{0.000}{10.9091}{mi}{66}\RG{28.383}{0.000}{10.9091}{rm}{41}\RG{35.638}{0.000}{10.9091}{rm}{61}\RG{47.171}{0.000}{10.9091}{mi}{39}\RG{54.307}{0.000}{10.9091}{rm}{40}\RG{58.550}{0.000}{10.9091}{mi}{66}\RG{67.372}{0.000}{10.9091}{rm}{41}\RG{71.615}{0.000}{10.9091}{mi}{39}\RG{78.751}{0.000}{10.9091}{rm}{40}\RG{82.994}{0.000}{10.9091}{mi}{65}\RG{91.176}{0.000}{10.9091}{rm}{41}\RG{95.418}{0.000}{10.9091}{mi}{58}%
}\hskip.1em\allowbreak{}

证明：或者第一个等式对任意\allowbreak{}\hskip.1em \RB{21.762}{8.315}{2.601}{%
\RG{0.000}{0.004}{10.9091}{mi}{65}\RG{8.182}{0.004}{10.9091}{mi}{59}\RG{13.015}{0.004}{10.9091}{mi}{66}%
}\hskip.1em\allowbreak{} 都成立，或者第二个等式对任意\allowbreak{}\hskip.1em \RB{21.762}{8.315}{2.601}{%
\RG{0.000}{0.004}{10.9091}{mi}{65}\RG{8.182}{0.004}{10.9091}{mi}{59}\RG{13.015}{0.004}{10.9091}{mi}{66}%
}\hskip.1em\allowbreak{} 都成立。

解：先用一个基本事实：线性空间不能是两个真子空间的并。事实上，若\allowbreak{}\hskip.1em \RB{54.780}{7.954}{0.737}{%
\RG{0.000}{0.003}{10.9091}{mi}{69}\RG{11.712}{0.003}{10.9091}{rm}{61}\RG{23.227}{0.003}{10.9091}{mi}{85}\RG{33.720}{0.003}{10.9091}{sy}{91}\RG{42.847}{0.003}{10.9091}{mi}{87}%
}\hskip.1em\allowbreak{} 且\allowbreak{}\hskip.1em \RB{24.184}{7.954}{2.602}{%
\RG{0.000}{0.003}{10.9091}{mi}{85}\RG{7.418}{0.003}{10.9091}{mi}{59}\RG{12.251}{0.003}{10.9091}{mi}{87}%
}\hskip.1em\allowbreak{}都是真子空间，取\allowbreak{}\hskip.1em \RB{49.220}{8.684}{3.225}{%
\RG{0.000}{0.002}{10.9091}{mi}{117}\RG{9.275}{0.002}{10.9091}{sy}{50}\RG{19.578}{0.002}{10.9091}{mi}{85}\RG{30.024}{0.002}{10.9091}{sy}{110}\RG{37.287}{0.002}{10.9091}{mi}{87}%
}\hskip.1em\allowbreak{} 、\allowbreak{}\hskip.1em \RB{51.117}{8.684}{3.225}{%
\RG{0.000}{0.002}{10.9091}{mi}{119}\RG{11.133}{0.002}{10.9091}{sy}{50}\RG{21.436}{0.002}{10.9091}{mi}{87}\RG{35.063}{0.002}{10.9091}{sy}{110}\RG{42.326}{0.002}{10.9091}{mi}{85}%
}\hskip.1em\allowbreak{} 。若\allowbreak{}\hskip.1em \RB{48.574}{7.953}{1.403}{%
\RG{0.000}{0.002}{10.9091}{mi}{117}\RG{8.053}{0.002}{10.9091}{rm}{43}\RG{18.347}{0.002}{10.9091}{mi}{119}\RG{29.480}{0.002}{10.9091}{sy}{50}\RG{39.783}{0.002}{10.9091}{mi}{85}%
}\hskip.1em\allowbreak{} ，则\allowbreak{}\hskip.1em \RB{30.228}{7.953}{0.934}{%
\RG{0.000}{0.002}{10.9091}{mi}{119}\RG{11.134}{0.002}{10.9091}{sy}{50}\RG{21.437}{0.002}{10.9091}{mi}{85}%
}\hskip.1em\allowbreak{} ；若\allowbreak{}\hskip.1em \RB{51.716}{7.953}{1.403}{%
\RG{0.000}{0.002}{10.9091}{mi}{117}\RG{8.053}{0.002}{10.9091}{rm}{43}\RG{18.347}{0.002}{10.9091}{mi}{119}\RG{29.480}{0.002}{10.9091}{sy}{50}\RG{39.783}{0.002}{10.9091}{mi}{87}%
}\hskip.1em\allowbreak{} ，则\allowbreak{}\hskip.1em \RB{31.511}{7.953}{0.934}{%
\RG{0.000}{0.002}{10.9091}{mi}{117}\RG{9.275}{0.002}{10.9091}{sy}{50}\RG{19.578}{0.002}{10.9091}{mi}{87}%
}\hskip.1em\allowbreak{} 。两种情况均矛盾。

记\allowbreak{}\hskip.1em \RB{55.221}{8.683}{3.226}{%
\RG{0.000}{0.001}{10.9091}{mi}{69}\RG{11.712}{0.001}{10.9091}{rm}{61}\RG{23.227}{0.001}{10.9091}{mi}{77}\RG{33.810}{-1.635}{7.9701}{mi}{110}\RG{39.446}{0.001}{10.9091}{rm}{40}\RG{43.689}{0.001}{10.9091}{bb}{82}\RG{51.568}{0.001}{10.9091}{rm}{41}%
}\hskip.1em\allowbreak{} 。固定\allowbreak{}\hskip.1em \RB{30.347}{8.312}{0.935}{%
\RG{0.000}{0.001}{10.9091}{mi}{65}\RG{11.212}{0.001}{10.9091}{sy}{50}\RG{21.515}{0.001}{10.9091}{mi}{69}%
}\hskip.1em\allowbreak{} ，令

\[
\RB{199.447}{8.683}{21.525}{%
\RG{2.855}{0.001}{10.9091}{mi}{85}\RG{20.175}{0.001}{10.9091}{rm}{61}\RG{31.690}{0.001}{10.9091}{sy}{102}\RG{37.144}{0.001}{10.9091}{mi}{66}\RG{48.997}{0.001}{10.9091}{sy}{50}\RG{59.300}{0.001}{10.9091}{mi}{69}\RG{71.013}{0.001}{10.9091}{rm}{58}\RG{77.073}{0.001}{10.9091}{mi}{39}\RG{84.210}{0.001}{10.9091}{rm}{40}\RG{88.452}{0.001}{10.9091}{mi}{65}\RG{96.634}{0.001}{10.9091}{mi}{66}\RG{105.456}{0.001}{10.9091}{rm}{41}\RG{112.123}{0.001}{10.9091}{sy}{0}\RG{123.032}{0.001}{10.9091}{mi}{39}\RG{130.168}{0.001}{10.9091}{rm}{40}\RG{134.410}{0.001}{10.9091}{mi}{65}\RG{142.592}{0.001}{10.9091}{rm}{41}\RG{146.835}{0.001}{10.9091}{mi}{39}\RG{153.971}{0.001}{10.9091}{rm}{40}\RG{158.213}{0.001}{10.9091}{mi}{66}\RG{167.035}{0.001}{10.9091}{rm}{41}\RG{174.290}{0.001}{10.9091}{rm}{61}\RG{185.799}{0.001}{10.9091}{rm}{48}\RG{191.279}{0.001}{10.9091}{sy}{103}\RG{196.733}{0.001}{10.9091}{mi}{59}\RG{10.303}{-1.687}{7.9701}{mi}{65}\RG{0.000}{-18.298}{10.9091}{mi}{87}\RG{20.175}{-18.298}{10.9091}{rm}{61}\RG{31.690}{-18.298}{10.9091}{sy}{102}\RG{37.144}{-18.298}{10.9091}{mi}{66}\RG{48.997}{-18.298}{10.9091}{sy}{50}\RG{59.300}{-18.298}{10.9091}{mi}{69}\RG{71.013}{-18.298}{10.9091}{rm}{58}\RG{77.073}{-18.298}{10.9091}{mi}{39}\RG{84.210}{-18.298}{10.9091}{rm}{40}\RG{88.452}{-18.298}{10.9091}{mi}{65}\RG{96.634}{-18.298}{10.9091}{mi}{66}\RG{105.456}{-18.298}{10.9091}{rm}{41}\RG{112.123}{-18.298}{10.9091}{sy}{0}\RG{123.032}{-18.298}{10.9091}{mi}{39}\RG{130.168}{-18.298}{10.9091}{rm}{40}\RG{134.410}{-18.298}{10.9091}{mi}{66}\RG{143.232}{-18.298}{10.9091}{rm}{41}\RG{147.475}{-18.298}{10.9091}{mi}{39}\RG{154.611}{-18.298}{10.9091}{rm}{40}\RG{158.854}{-18.298}{10.9091}{mi}{65}\RG{167.035}{-18.298}{10.9091}{rm}{41}\RG{174.290}{-18.298}{10.9091}{rm}{61}\RG{185.799}{-18.298}{10.9091}{rm}{48}\RG{191.279}{-18.298}{10.9091}{sy}{103}\RG{196.733}{-18.298}{10.9091}{mi}{58}\RG{10.303}{-19.987}{7.9701}{mi}{65}%
}
\]
其中关于\allowbreak{}\hskip.1em \RB{8.747}{7.949}{0.502}{%
\RG{0.000}{-0.002}{10.9091}{mi}{66}%
}\hskip.1em\allowbreak{} 的两个映射都是线性的，故\allowbreak{}\hskip.1em \RB{36.030}{7.949}{2.607}{%
\RG{0.000}{-0.002}{10.9091}{mi}{85}\RG{7.448}{-1.691}{7.9701}{mi}{65}\RG{14.289}{-0.002}{10.9091}{mi}{59}\RG{19.122}{-0.002}{10.9091}{mi}{87}\RG{29.441}{-1.691}{7.9701}{mi}{65}%
}\hskip.1em\allowbreak{} 都是子空间。由题设，\allowbreak{}\hskip.1em \RB{65.702}{7.949}{2.191}{%
\RG{0.000}{-0.002}{10.9091}{mi}{69}\RG{11.712}{-0.002}{10.9091}{rm}{61}\RG{23.227}{-0.002}{10.9091}{mi}{85}\RG{30.676}{-1.691}{7.9701}{mi}{65}\RG{39.526}{-0.002}{10.9091}{sy}{91}\RG{48.809}{-0.002}{10.9091}{mi}{87}\RG{59.113}{-1.691}{7.9701}{mi}{65}%
}\hskip.1em\allowbreak{} ，所以对每个固定的\allowbreak{}\hskip.1em \RB{8.366}{8.308}{0.503}{%
\RG{0.000}{-0.003}{10.9091}{mi}{65}%
}\hskip.1em\allowbreak{} ，必有

\allowbreak{}\hskip.1em \RB{37.667}{7.948}{2.192}{%
\RG{0.000}{-0.003}{10.9091}{mi}{85}\RG{7.449}{-1.692}{7.9701}{mi}{65}\RG{17.320}{-0.003}{10.9091}{rm}{61}\RG{28.835}{-0.003}{10.9091}{mi}{69}%
}\hskip.1em\allowbreak{} 或\allowbreak{}\hskip.1em \RB{42.958}{7.948}{2.192}{%
\RG{0.000}{-0.003}{10.9091}{mi}{87}\RG{10.303}{-1.692}{7.9701}{mi}{65}\RG{20.175}{-0.003}{10.9091}{rm}{61}\RG{31.690}{-0.003}{10.9091}{mi}{69}\RG{40.363}{-0.003}{10.9091}{mi}{58}%
}\hskip.1em\allowbreak{}

再令

\[
\RB{256.219}{8.682}{3.227}{%
\RG{0.000}{0.000}{10.9091}{sy}{85}\RG{10.940}{0.000}{10.9091}{rm}{61}\RG{22.456}{0.000}{10.9091}{sy}{102}\RG{27.910}{0.000}{10.9091}{mi}{65}\RG{39.122}{0.000}{10.9091}{sy}{50}\RG{49.425}{0.000}{10.9091}{mi}{69}\RG{61.137}{0.000}{10.9091}{rm}{58}\RG{67.199}{0.000}{10.9091}{mi}{85}\RG{84.519}{0.000}{10.9091}{rm}{61}\RG{96.034}{0.000}{10.9091}{mi}{69}\RG{104.716}{0.000}{10.9091}{sy}{103}\RG{110.170}{0.000}{10.9091}{mi}{59}\RG{136.837}{0.000}{10.9091}{sy}{87}\RG{151.540}{0.000}{10.9091}{rm}{61}\RG{163.055}{0.000}{10.9091}{sy}{102}\RG{168.510}{0.000}{10.9091}{mi}{65}\RG{179.722}{0.000}{10.9091}{sy}{50}\RG{190.026}{0.000}{10.9091}{mi}{69}\RG{201.738}{0.000}{10.9091}{rm}{58}\RG{207.798}{0.000}{10.9091}{mi}{87}\RG{227.973}{0.000}{10.9091}{rm}{61}\RG{239.488}{0.000}{10.9091}{mi}{69}\RG{248.170}{0.000}{10.9091}{sy}{103}\RG{253.624}{0.000}{10.9091}{mi}{58}\RG{74.647}{-1.689}{7.9701}{mi}{65}\RG{218.101}{-1.689}{7.9701}{mi}{65}%
}
\]
它们也都是子空间。例如，固定\allowbreak{}\hskip.1em \RB{8.747}{7.954}{0.500}{%
\RG{0.000}{0.003}{10.9091}{mi}{66}%
}\hskip.1em\allowbreak{} 后，映射\allowbreak{}\hskip.1em \RB{118.178}{8.685}{3.224}{%
\RG{0.000}{0.003}{10.9091}{mi}{65}\RG{11.213}{0.003}{10.9091}{sy}{55}\RG{11.213}{0.003}{10.9091}{sy}{33}\RG{25.152}{0.003}{10.9091}{mi}{39}\RG{32.288}{0.003}{10.9091}{rm}{40}\RG{36.531}{0.003}{10.9091}{mi}{65}\RG{44.713}{0.003}{10.9091}{mi}{66}\RG{53.535}{0.003}{10.9091}{rm}{41}\RG{59.907}{0.003}{10.9091}{sy}{0}\RG{70.522}{0.003}{10.9091}{mi}{39}\RG{77.658}{0.003}{10.9091}{rm}{40}\RG{81.900}{0.003}{10.9091}{mi}{65}\RG{90.082}{0.003}{10.9091}{rm}{41}\RG{94.325}{0.003}{10.9091}{mi}{39}\RG{101.461}{0.003}{10.9091}{rm}{40}\RG{105.704}{0.003}{10.9091}{mi}{66}\RG{114.525}{0.003}{10.9091}{rm}{41}%
}\hskip.1em\allowbreak{} 是线性的，而\allowbreak{}\hskip.1em \RB{8.006}{7.954}{0.802}{%
\RG{0.000}{0.003}{10.9091}{sy}{85}%
}\hskip.1em\allowbreak{} 正是这些映射的核的交；\allowbreak{}\hskip.1em \RB{11.802}{7.954}{0.988}{%
\RG{0.000}{0.003}{10.9091}{sy}{87}%
}\hskip.1em\allowbreak{} 同理。

由上一步，\allowbreak{}\hskip.1em \RB{55.061}{7.953}{0.989}{%
\RG{0.000}{0.002}{10.9091}{mi}{69}\RG{11.719}{0.002}{10.9091}{rm}{61}\RG{23.241}{0.002}{10.9091}{sy}{85}\RG{33.561}{0.002}{10.9091}{sy}{91}\RG{43.260}{0.002}{10.9091}{sy}{87}%
}\hskip.1em\allowbreak{} 。再次应用基本事实，得\allowbreak{}\hskip.1em \RB{31.302}{7.953}{0.803}{%
\RG{0.000}{0.002}{10.9091}{sy}{85}\RG{10.948}{0.002}{10.9091}{rm}{61}\RG{22.470}{0.002}{10.9091}{mi}{69}%
}\hskip.1em\allowbreak{} 或\allowbreak{}\hskip.1em \RB{35.064}{7.953}{0.989}{%
\RG{0.000}{0.002}{10.9091}{sy}{87}\RG{14.710}{0.002}{10.9091}{rm}{61}\RG{26.232}{0.002}{10.9091}{mi}{69}%
}\hskip.1em\allowbreak{} ，即题目所要求的两个等式中至少有一个恒成立。\allowbreak{}\hskip.1em \RB{8.366}{7.996}{0.500}{%
\RG{0.000}{0.001}{10.9091}{am}{3}%
}\hskip.1em\allowbreak{}

小结：关键是两次使用“两个子空间的并覆盖全空间”。第一次固定\allowbreak{}\hskip.1em \RB{7.683}{7.634}{0.500}{%
\RG{0.000}{0.001}{10.9091}{mi}{65}%
}\hskip.1em\allowbreak{} ，消除对\allowbreak{}\hskip.1em \RB{8.032}{7.306}{0.500}{%
\RG{0.000}{0.001}{10.9091}{mi}{66}%
}\hskip.1em\allowbreak{} 的选择依赖；第二次消除对\allowbreak{}\hskip.1em \RB{7.683}{7.635}{0.500}{%
\RG{0.000}{0.002}{10.9091}{mi}{65}%
}\hskip.1em\allowbreak{} 的选择依赖。不能直接交换“任意”与“或者”的次序。

% Original page 2.
\section*{2015S02三对角矩阵的对角化}
题目设复数\allowbreak{}\hskip.1em \RB{25.307}{8.072}{2.604}{%
\RG{0.000}{0.001}{10.9091}{mi}{97}\RG{5.760}{0.001}{10.9091}{mi}{59}\RG{10.593}{0.001}{10.9091}{mi}{98}\RG{15.273}{0.001}{10.9091}{mi}{59}\RG{20.117}{0.001}{10.9091}{mi}{99}%
}\hskip.1em\allowbreak{} 满足\allowbreak{}\hskip.1em \RB{29.450}{8.312}{2.855}{%
\RG{0.000}{0.001}{10.9091}{mi}{98}\RG{4.680}{0.001}{10.9091}{mi}{99}\RG{12.433}{0.001}{10.9091}{sy}{54}\RG{12.433}{0.001}{10.9091}{rm}{61}\RG{23.931}{0.001}{10.9091}{rm}{48}%
}\hskip.1em\allowbreak{} ，证明下列\allowbreak{}\hskip.1em \RB{6.729}{5.323}{0.619}{%
\RG{0.000}{0.001}{10.9091}{mi}{110}%
}\hskip.1em\allowbreak{} 阶方阵可对角化：

\[
\RB{135.819}{25.317}{54.208}{%
\RG{22.727}{24.391}{10.9092}{ex}{48}\RG{121.861}{24.391}{10.9092}{ex}{49}\RG{32.273}{15.307}{10.9091}{mi}{97}\RG{52.019}{15.307}{10.9091}{mi}{98}\RG{22.727}{5.191}{10.9092}{ex}{66}\RG{121.861}{5.191}{10.9092}{ex}{67}\RG{32.796}{-0.003}{10.9091}{mi}{99}\RG{51.473}{-0.003}{10.9091}{mi}{97}\RG{74.709}{-0.003}{10.9091}{mi}{98}\RG{22.727}{-1.354}{10.9092}{ex}{66}\RG{121.861}{-1.354}{10.9092}{ex}{67}\RG{22.727}{-7.900}{10.9092}{ex}{66}\RG{121.861}{-7.900}{10.9092}{ex}{67}\RT{48.608}{-12.727}{10.9091}{.}\RT{71.306}{-12.727}{10.9091}{.}\RT{94.003}{-12.727}{10.9091}{.}\RG{22.727}{-14.445}{10.9092}{ex}{66}\RG{121.861}{-14.445}{10.9092}{ex}{67}\RT{52.853}{-15.716}{10.9091}{.}\RT{75.551}{-15.716}{10.9091}{.}\RT{98.248}{-15.716}{10.9091}{.}\RG{0.000}{-17.173}{10.9091}{mi}{65}\RG{11.212}{-17.173}{10.9091}{rm}{61}\RG{133.224}{-17.173}{10.9091}{mi}{58}\RT{57.098}{-18.704}{10.9091}{.}\RT{79.796}{-18.704}{10.9091}{.}\RT{102.493}{-18.704}{10.9091}{.}\RG{22.727}{-20.991}{10.9092}{ex}{66}\RG{121.861}{-20.991}{10.9092}{ex}{67}\RG{22.727}{-27.536}{10.9092}{ex}{66}\RG{121.861}{-27.536}{10.9092}{ex}{67}\RG{22.727}{-34.518}{10.9092}{ex}{64}\RG{121.861}{-34.518}{10.9092}{ex}{65}\RG{74.707}{-35.011}{10.9091}{mi}{99}\RG{96.863}{-35.011}{10.9091}{mi}{97}\RG{116.620}{-35.011}{10.9091}{mi}{98}\RG{97.404}{-50.322}{10.9091}{mi}{99}\RG{116.080}{-50.322}{10.9091}{mi}{97}%
}
\]
解：取\allowbreak{}\hskip.1em \RB{26.420}{8.175}{0.941}{%
\RG{0.000}{-0.005}{10.9091}{mi}{115}\RG{8.144}{-0.005}{10.9091}{sy}{50}\RG{18.447}{-0.005}{10.9091}{bb}{67}%
}\hskip.1em\allowbreak{} 满足\allowbreak{}\hskip.1em \RB{34.262}{9.753}{0.625}{%
\RG{0.000}{-0.005}{10.9091}{mi}{115}\RG{5.113}{3.953}{7.9701}{rm}{50}\RG{12.876}{-0.005}{10.9091}{rm}{61}\RG{24.391}{-0.005}{10.9091}{mi}{98}\RG{29.071}{-0.005}{10.9091}{mi}{99}%
}\hskip.1em\allowbreak{} ，并令

\[
\RB{169.755}{12.697}{8.108}{%
\RG{20.965}{7.375}{10.9091}{mi}{115}\RG{146.461}{4.499}{7.9701}{mi}{110}\RG{151.599}{4.499}{7.9701}{sy}{0}\RG{158.186}{4.499}{7.9701}{rm}{49}\RG{0.000}{-0.005}{10.9091}{mi}{114}\RG{8.255}{-0.005}{10.9091}{rm}{61}\RG{27.275}{-0.005}{10.9091}{mi}{59}\RG{53.926}{-0.005}{10.9091}{mi}{68}\RG{66.307}{-0.005}{10.9091}{rm}{61}\RT{77.823}{-0.005}{10.9091}{d}\RT{83.888}{-0.005}{10.9091}{i}\RT{86.921}{-0.005}{10.9091}{a}\RT{92.375}{-0.005}{10.9091}{g}\RG{97.830}{-0.005}{10.9091}{rm}{40}\RG{102.063}{-0.005}{10.9091}{rm}{49}\RG{107.527}{-0.005}{10.9091}{mi}{59}\RG{112.360}{-0.005}{10.9091}{mi}{114}\RG{116.985}{-0.005}{10.9091}{mi}{59}\RG{121.818}{-0.005}{10.9091}{mi}{58}\RG{126.662}{-0.005}{10.9091}{mi}{58}\RG{131.494}{-0.005}{10.9091}{mi}{58}\RG{136.338}{-0.005}{10.9091}{mi}{59}\RG{141.182}{-0.005}{10.9091}{mi}{114}\RG{162.918}{-0.005}{10.9091}{rm}{41}\RG{167.160}{-0.005}{10.9091}{mi}{58}\RG{21.181}{-7.488}{10.9091}{mi}{98}\RR{20.966}{2.722}{5.113}{0.436}%
}
\]
因为\allowbreak{}\hskip.1em \RB{29.450}{8.311}{2.856}{%
\RG{0.000}{0.000}{10.9091}{mi}{98}\RG{4.680}{0.000}{10.9091}{mi}{99}\RG{12.433}{0.000}{10.9091}{sy}{54}\RG{12.433}{0.000}{10.9091}{rm}{61}\RG{23.931}{0.000}{10.9091}{rm}{48}%
}\hskip.1em\allowbreak{} ，所以\allowbreak{}\hskip.1em \RB{35.234}{8.311}{2.856}{%
\RG{0.000}{0.000}{10.9091}{mi}{115}\RG{5.105}{0.000}{10.9091}{mi}{59}\RG{9.938}{0.000}{10.9091}{mi}{114}\RG{18.217}{0.000}{10.9091}{sy}{54}\RG{18.217}{0.000}{10.9091}{rm}{61}\RG{29.715}{0.000}{10.9091}{rm}{48}%
}\hskip.1em\allowbreak{} ，\allowbreak{}\hskip.1em \RB{9.271}{7.951}{0.500}{%
\RG{0.000}{0.000}{10.9091}{mi}{68}%
}\hskip.1em\allowbreak{} 可逆，且

\[
\RB{86.095}{12.703}{8.102}{%
\RG{57.680}{7.381}{10.9091}{mi}{99}\RG{0.000}{0.001}{10.9091}{mi}{98}\RG{4.680}{0.001}{10.9091}{mi}{114}\RG{12.937}{0.001}{10.9091}{rm}{61}\RG{24.452}{0.001}{10.9091}{mi}{115}\RG{29.558}{0.001}{10.9091}{mi}{59}\RG{66.879}{0.001}{10.9091}{rm}{61}\RG{78.395}{0.001}{10.9091}{mi}{115}\RG{83.500}{0.001}{10.9091}{mi}{58}\RG{57.428}{-7.482}{10.9091}{mi}{114}\RR{57.428}{2.728}{5.225}{0.436}%
}
\]
因此\allowbreak{}\hskip.1em \RB{38.107}{9.756}{0.503}{%
\RG{0.000}{-0.003}{10.9091}{mi}{68}\RG{9.335}{3.956}{7.9701}{sy}{0}\RG{15.922}{3.956}{7.9701}{rm}{49}\RG{20.654}{-0.003}{10.9091}{mi}{65}\RG{28.836}{-0.003}{10.9091}{mi}{68}%
}\hskip.1em\allowbreak{} 的主对角元为\allowbreak{}\hskip.1em \RB{5.933}{5.319}{0.623}{%
\RG{0.000}{-0.003}{10.9091}{mi}{97}%
}\hskip.1em\allowbreak{} ，上下两条邻接对角线的元素都为\allowbreak{}\hskip.1em \RB{5.082}{5.319}{0.623}{%
\RG{0.000}{-0.003}{10.9091}{mi}{115}%
}\hskip.1em\allowbreak{} ，即

\[
\RB{256.216}{42.490}{37.035}{%
\RG{143.748}{41.564}{10.9092}{ex}{48}\RG{242.258}{41.564}{10.9092}{ex}{49}\RG{153.294}{32.479}{10.9091}{rm}{48}\RG{172.342}{32.479}{10.9091}{rm}{49}\RG{143.748}{22.363}{10.9092}{ex}{66}\RG{242.258}{22.363}{10.9092}{ex}{67}\RG{153.294}{17.169}{10.9091}{rm}{49}\RG{172.342}{17.169}{10.9091}{rm}{48}\RG{195.044}{17.169}{10.9091}{rm}{49}\RG{143.748}{15.818}{10.9092}{ex}{66}\RG{242.258}{15.818}{10.9092}{ex}{67}\RG{143.748}{9.272}{10.9092}{ex}{66}\RG{242.258}{9.272}{10.9092}{ex}{67}\RG{9.335}{4.504}{7.9701}{sy}{0}\RG{15.921}{4.504}{7.9701}{rm}{49}\RT{169.317}{4.445}{10.9091}{.}\RT{192.015}{4.445}{10.9091}{.}\RT{214.712}{4.445}{10.9091}{.}\RG{143.748}{2.727}{10.9092}{ex}{66}\RG{242.258}{2.727}{10.9092}{ex}{67}\RT{173.562}{1.456}{10.9091}{.}\RT{196.260}{1.456}{10.9091}{.}\RT{218.957}{1.456}{10.9091}{.}\RG{0.000}{0.000}{10.9091}{mi}{68}\RG{20.654}{0.000}{10.9091}{mi}{65}\RG{28.836}{0.000}{10.9091}{mi}{68}\RG{41.201}{0.000}{10.9091}{rm}{61}\RG{52.716}{0.000}{10.9091}{mi}{97}\RG{58.476}{0.000}{10.9091}{mi}{73}\RG{71.339}{0.000}{10.9091}{rm}{43}\RG{82.248}{0.000}{10.9091}{mi}{115}\RG{87.353}{0.000}{10.9091}{mi}{84}\RG{94.630}{0.000}{10.9091}{mi}{59}\RG{121.292}{0.000}{10.9091}{mi}{84}\RG{132.233}{0.000}{10.9091}{rm}{61}\RG{253.621}{-0.000}{10.9091}{mi}{58}\RT{177.807}{-1.532}{10.9091}{.}\RT{200.505}{-1.532}{10.9091}{.}\RT{223.202}{-1.532}{10.9091}{.}\RG{63.279}{-1.637}{7.9701}{mi}{110}\RG{143.748}{-3.819}{10.9092}{ex}{66}\RG{242.258}{-3.819}{10.9092}{ex}{67}\RG{143.748}{-10.364}{10.9092}{ex}{66}\RG{242.258}{-10.364}{10.9092}{ex}{67}\RG{143.748}{-17.346}{10.9092}{ex}{64}\RG{242.258}{-17.346}{10.9092}{ex}{65}\RG{195.049}{-17.839}{10.9091}{rm}{49}\RG{217.740}{-17.839}{10.9091}{rm}{48}\RG{236.798}{-17.839}{10.9091}{rm}{49}\RG{217.746}{-33.150}{10.9091}{rm}{49}\RG{236.793}{-33.150}{10.9091}{rm}{48}%
}
\]
\allowbreak{}\hskip.1em \RB{8.180}{7.886}{0.500}{%
\RG{0.000}{0.001}{10.9091}{mi}{84}%
}\hskip.1em\allowbreak{} 是实对称矩阵，故可正交对角化，从而\allowbreak{}\hskip.1em \RB{42.817}{7.952}{2.215}{%
\RG{0.000}{0.001}{10.9091}{mi}{97}\RG{5.760}{0.001}{10.9091}{mi}{73}\RG{10.561}{-1.635}{7.9701}{mi}{110}\RG{18.622}{0.001}{10.9091}{rm}{43}\RG{29.531}{0.001}{10.9091}{mi}{115}\RG{34.637}{0.001}{10.9091}{mi}{84}%
}\hskip.1em\allowbreak{} 以及\allowbreak{}\hskip.1em \RB{8.366}{8.312}{0.500}{%
\RG{0.000}{0.001}{10.9091}{mi}{65}%
}\hskip.1em\allowbreak{} 都可对角化。

下面进一步写出特征值与对角化矩阵。对\allowbreak{}\hskip.1em \RB{56.949}{8.071}{2.606}{%
\RG{0.000}{-0.000}{10.9091}{mi}{107}\RG{9.053}{-0.000}{10.9091}{rm}{61}\RG{20.551}{-0.000}{10.9091}{rm}{49}\RG{26.023}{-0.000}{10.9091}{mi}{59}\RG{30.856}{-0.000}{10.9091}{mi}{58}\RG{35.700}{-0.000}{10.9091}{mi}{58}\RG{40.532}{-0.000}{10.9091}{mi}{58}\RG{45.376}{-0.000}{10.9091}{mi}{59}\RG{50.220}{-0.000}{10.9091}{mi}{110}%
}\hskip.1em\allowbreak{} ，令

\[
\RB{239.020}{15.449}{8.891}{%
\RG{32.334}{7.378}{10.9091}{mi}{107}\RG{38.345}{7.378}{10.9091}{mi}{25}\RT{230.140}{4.502}{7.9701}{T}\RG{0.000}{-0.002}{10.9091}{mi}{18}\RG{13.272}{-0.002}{10.9091}{rm}{61}\RG{52.515}{-0.002}{10.9091}{mi}{59}\RG{79.166}{-0.002}{10.9091}{mi}{121}\RG{92.680}{-0.002}{10.9091}{rm}{61}\RG{104.178}{-0.002}{10.9091}{rm}{40}\RT{108.437}{-0.002}{10.9091}{s}\RT{112.735}{-0.002}{10.9091}{i}\RT{115.768}{-0.002}{10.9091}{n}\RG{123.652}{-0.002}{10.9091}{mi}{18}\RG{133.893}{-0.002}{10.9091}{mi}{59}\RT{138.741}{-0.002}{10.9091}{s}\RT{143.039}{-0.002}{10.9091}{i}\RT{146.072}{-0.002}{10.9091}{n}\RG{153.956}{-0.002}{10.9091}{rm}{50}\RG{159.410}{-0.002}{10.9091}{mi}{18}\RG{169.651}{-0.002}{10.9091}{mi}{59}\RG{174.484}{-0.002}{10.9091}{mi}{58}\RG{179.328}{-0.002}{10.9091}{mi}{58}\RG{184.172}{-0.002}{10.9091}{mi}{58}\RG{189.004}{-0.002}{10.9091}{mi}{59}\RT{193.895}{-0.002}{10.9091}{s}\RT{198.193}{-0.002}{10.9091}{i}\RT{201.226}{-0.002}{10.9091}{n}\RG{209.109}{-0.002}{10.9091}{mi}{110}\RG{215.654}{-0.002}{10.9091}{mi}{18}\RG{225.898}{-0.002}{10.9091}{rm}{41}\RG{236.425}{-0.002}{10.9091}{mi}{58}\RG{5.122}{-1.779}{7.9701}{mi}{107}\RG{84.530}{-1.779}{7.9701}{mi}{107}\RG{128.773}{-1.779}{7.9701}{mi}{107}\RG{164.532}{-1.779}{7.9701}{mi}{107}\RG{220.778}{-1.779}{7.9701}{mi}{107}\RG{25.982}{-7.485}{10.9091}{mi}{110}\RG{34.954}{-7.485}{10.9091}{rm}{43}\RG{45.853}{-7.485}{10.9091}{rm}{49}\RR{25.982}{2.725}{25.336}{0.436}%
}
\]
\allowbreak{}\hskip.1em \RB{30.515}{8.311}{2.856}{%
\RG{0.000}{0.000}{10.9091}{mi}{121}\RG{5.348}{-1.778}{7.9701}{mi}{107}\RG{13.498}{0.000}{10.9091}{sy}{54}\RG{13.498}{0.000}{10.9091}{rm}{61}\RG{24.996}{0.000}{10.9091}{rm}{48}%
}\hskip.1em\allowbreak{} ，因为\allowbreak{}\hskip.1em \RB{51.455}{8.191}{2.358}{%
\RG{0.000}{0.000}{10.9091}{rm}{48}\RG{8.485}{0.000}{10.9091}{mi}{60}\RG{19.983}{0.000}{10.9091}{mi}{18}\RG{25.121}{-1.778}{7.9701}{mi}{107}\RG{33.271}{0.000}{10.9091}{mi}{60}\RG{44.770}{0.000}{10.9091}{mi}{25}%
}\hskip.1em\allowbreak{} 。利用sin\allowbreak{}\hskip.1em \RB{16.842}{7.766}{0.740}{%
\RG{0.000}{0.000}{10.9091}{rm}{48}\RG{8.476}{0.000}{10.9091}{rm}{61}%
}\hskip.1em\allowbreak{} sin\allowbreak{}\hskip.1em \RB{72.586}{8.682}{3.227}{%
\RG{0.000}{0.000}{10.9091}{rm}{40}\RG{4.233}{0.000}{10.9091}{rm}{40}\RG{8.485}{0.000}{10.9091}{mi}{110}\RG{17.457}{0.000}{10.9091}{rm}{43}\RG{28.355}{0.000}{10.9091}{rm}{49}\RG{33.810}{0.000}{10.9091}{rm}{41}\RG{38.063}{0.000}{10.9091}{mi}{18}\RG{43.184}{-1.778}{7.9701}{mi}{107}\RG{48.304}{0.000}{10.9091}{rm}{41}\RG{55.559}{0.000}{10.9091}{rm}{61}\RG{67.068}{0.000}{10.9091}{rm}{48}%
}\hskip.1em\allowbreak{} 和恒等式

\[
\RB{232.604}{9.293}{3.227}{%
\RT{0.000}{0.000}{10.9091}{s}\RT{4.298}{0.000}{10.9091}{i}\RT{7.331}{0.000}{10.9091}{n}\RG{13.396}{0.000}{10.9091}{rm}{40}\RG{17.629}{0.000}{10.9091}{rm}{40}\RG{21.881}{0.000}{10.9091}{mi}{106}\RG{29.422}{0.000}{10.9091}{sy}{0}\RG{40.331}{0.000}{10.9091}{rm}{49}\RG{45.786}{0.000}{10.9091}{rm}{41}\RG{50.028}{0.000}{10.9091}{mi}{18}\RG{60.269}{0.000}{10.9091}{rm}{41}\RG{66.924}{0.000}{10.9091}{rm}{43}\RT{77.846}{0.000}{10.9091}{s}\RT{82.144}{0.000}{10.9091}{i}\RT{85.177}{0.000}{10.9091}{n}\RG{91.242}{0.000}{10.9091}{rm}{40}\RG{95.474}{0.000}{10.9091}{rm}{40}\RG{99.727}{0.000}{10.9091}{mi}{106}\RG{107.268}{0.000}{10.9091}{rm}{43}\RG{118.166}{0.000}{10.9091}{rm}{49}\RG{123.621}{0.000}{10.9091}{rm}{41}\RG{127.874}{0.000}{10.9091}{mi}{18}\RG{138.115}{0.000}{10.9091}{rm}{41}\RG{145.370}{0.000}{10.9091}{rm}{61}\RG{156.879}{0.000}{10.9091}{rm}{50}\RT{164.176}{0.000}{10.9091}{c}\RT{169.020}{0.000}{10.9091}{o}\RT{174.475}{0.000}{10.9091}{s}\RG{180.590}{0.000}{10.9091}{mi}{18}\RT{192.651}{0.000}{10.9091}{s}\RT{196.949}{0.000}{10.9091}{i}\RT{199.982}{0.000}{10.9091}{n}\RG{206.047}{0.000}{10.9091}{rm}{40}\RG{210.289}{0.000}{10.9091}{mi}{106}\RG{215.394}{0.000}{10.9091}{mi}{18}\RG{225.647}{0.000}{10.9091}{rm}{41}\RG{229.890}{0.000}{10.9091}{mi}{59}\RG{55.150}{-1.777}{7.9701}{mi}{107}\RG{132.996}{-1.777}{7.9701}{mi}{107}\RG{185.713}{-1.777}{7.9701}{mi}{107}\RG{220.527}{-1.777}{7.9701}{mi}{107}%
}
\]
逐分量可得\allowbreak{}\hskip.1em \RB{38.285}{7.886}{2.735}{%
\RG{0.000}{0.001}{10.9091}{mi}{84}\RG{7.876}{0.001}{10.9091}{mi}{121}\RG{13.238}{-1.777}{7.9701}{mi}{107}\RG{21.388}{0.001}{10.9091}{rm}{61}\RG{32.886}{0.001}{10.9091}{rm}{50}%
}\hskip.1em\allowbreak{} cos\allowbreak{}\hskip.1em \RB{22.155}{8.192}{2.735}{%
\RG{0.000}{0.001}{10.9091}{mi}{18}\RG{5.121}{-1.777}{7.9701}{mi}{107}\RG{12.059}{0.001}{10.9091}{mi}{121}\RG{17.407}{-1.777}{7.9701}{mi}{107}%
}\hskip.1em\allowbreak{} ，故

\[
\RB{219.995}{15.452}{8.888}{%
\RG{193.834}{7.381}{10.9091}{mi}{107}\RG{199.845}{7.381}{10.9091}{mi}{25}\RG{0.000}{0.001}{10.9091}{mi}{65}\RG{8.182}{0.001}{10.9091}{rm}{40}\RG{12.424}{0.001}{10.9091}{mi}{68}\RG{21.741}{0.001}{10.9091}{mi}{121}\RG{32.228}{0.001}{10.9091}{rm}{41}\RG{39.483}{0.001}{10.9091}{rm}{61}\RG{51.015}{0.001}{10.9091}{mi}{21}\RG{62.500}{0.001}{10.9091}{mi}{68}\RG{71.816}{0.001}{10.9091}{mi}{121}\RG{82.303}{0.001}{10.9091}{mi}{59}\RG{112.347}{0.001}{10.9091}{mi}{21}\RG{126.871}{0.001}{10.9091}{rm}{61}\RG{138.386}{0.001}{10.9091}{mi}{97}\RG{146.576}{0.001}{10.9091}{rm}{43}\RG{157.474}{0.001}{10.9091}{rm}{50}\RG{162.940}{0.001}{10.9091}{mi}{115}\RT{169.872}{0.001}{10.9091}{c}\RT{174.716}{0.001}{10.9091}{o}\RT{180.171}{0.001}{10.9091}{s}\RG{217.401}{0.001}{10.9091}{mi}{58}\RG{27.108}{-1.776}{7.9701}{mi}{107}\RG{57.379}{-1.776}{7.9701}{mi}{107}\RG{77.183}{-1.776}{7.9701}{mi}{107}\RG{118.721}{-1.776}{7.9701}{mi}{107}\RG{187.483}{-7.482}{10.9091}{mi}{110}\RG{196.455}{-7.482}{10.9091}{rm}{43}\RG{207.353}{-7.482}{10.9091}{rm}{49}\RR{187.483}{2.728}{25.335}{0.436}%
}
\]
这些特征值两两不同：cos\allowbreak{}\hskip.1em \RB{6.249}{5.326}{0.616}{%
\RG{0.000}{0.004}{10.9091}{mi}{120}%
}\hskip.1em\allowbreak{} 在\allowbreak{}\hskip.1em \RB{24.809}{8.686}{3.223}{%
\RG{0.000}{0.004}{10.9091}{rm}{40}\RG{4.233}{0.004}{10.9091}{rm}{48}\RG{9.697}{0.004}{10.9091}{mi}{59}\RG{14.530}{0.004}{10.9091}{mi}{25}\RG{21.156}{0.004}{10.9091}{rm}{41}%
}\hskip.1em\allowbreak{} 上严格递减，且\allowbreak{}\hskip.1em \RB{25.161}{8.315}{2.852}{%
\RG{0.000}{0.004}{10.9091}{mi}{115}\RG{8.144}{0.004}{10.9091}{sy}{54}\RG{8.144}{0.004}{10.9091}{rm}{61}\RG{19.642}{0.004}{10.9091}{rm}{48}%
}\hskip.1em\allowbreak{} 。于是\allowbreak{}\hskip.1em \RB{63.719}{7.955}{2.732}{%
\RG{0.000}{0.004}{10.9091}{mi}{68}\RG{9.316}{0.004}{10.9091}{mi}{121}\RG{14.683}{-1.633}{7.9701}{rm}{49}\RG{19.415}{0.004}{10.9091}{mi}{59}\RG{24.248}{0.004}{10.9091}{mi}{58}\RG{29.092}{0.004}{10.9091}{mi}{58}\RG{33.935}{0.004}{10.9091}{mi}{58}\RG{38.768}{0.004}{10.9091}{mi}{59}\RG{43.612}{0.004}{10.9091}{mi}{68}\RG{52.939}{0.004}{10.9091}{mi}{121}\RG{58.341}{-1.633}{7.9701}{mi}{110}%
}\hskip.1em\allowbreak{} 线性无关。特别地，令

\[
\RB{142.277}{16.322}{12.985}{%
\RG{23.064}{15.385}{10.9092}{ex}{18}\RG{96.283}{15.385}{10.9092}{ex}{19}\RG{73.545}{7.383}{10.9091}{mi}{106}\RG{78.650}{7.383}{10.9091}{mi}{107}\RG{84.661}{7.383}{10.9091}{mi}{25}\RG{36.319}{4.507}{7.9701}{mi}{106}\RG{40.203}{4.507}{7.9701}{sy}{0}\RG{46.790}{4.507}{7.9701}{rm}{49}\RG{0.000}{0.003}{10.9091}{mi}{80}\RG{11.549}{0.003}{10.9091}{rm}{61}\RG{31.095}{0.003}{10.9091}{mi}{114}\RT{53.340}{0.003}{10.9091}{s}\RT{57.639}{0.003}{10.9091}{i}\RT{60.671}{0.003}{10.9091}{n}\RG{139.562}{0.003}{10.9091}{mi}{59}\RG{69.751}{-7.480}{10.9091}{mi}{110}\RG{78.724}{-7.480}{10.9091}{rm}{43}\RG{89.622}{-7.480}{10.9091}{rm}{49}\RG{104.313}{-10.859}{7.9701}{rm}{49}\RG{108.547}{-10.859}{7.9701}{sy}{20}\RG{115.134}{-10.859}{7.9701}{mi}{106}\RG{118.545}{-10.859}{7.9701}{mi}{59}\RG{120.896}{-10.859}{7.9701}{mi}{107}\RG{125.521}{-10.859}{7.9701}{sy}{20}\RG{132.108}{-10.859}{7.9701}{mi}{110}\RR{69.751}{2.730}{25.336}{0.436}%
}
\]
则\allowbreak{}\hskip.1em \RB{8.726}{7.955}{0.500}{%
\RG{0.000}{0.004}{10.9091}{mi}{80}%
}\hskip.1em\allowbreak{} 可逆，且\allowbreak{}\hskip.1em \RB{47.935}{9.763}{0.500}{%
\RG{0.000}{0.004}{10.9091}{mi}{80}\RG{8.519}{3.963}{7.9701}{sy}{0}\RG{15.106}{3.963}{7.9701}{rm}{49}\RG{19.838}{0.004}{10.9091}{mi}{65}\RG{28.020}{0.004}{10.9091}{mi}{80}\RG{39.569}{0.004}{10.9091}{rm}{61}%
}\hskip.1em\allowbreak{} diag\allowbreak{}\hskip.1em \RB{55.234}{8.686}{3.223}{%
\RG{0.000}{0.004}{10.9091}{rm}{40}\RG{4.243}{0.004}{10.9091}{mi}{21}\RG{10.606}{-1.632}{7.9701}{rm}{49}\RG{15.339}{0.004}{10.9091}{mi}{59}\RG{20.172}{0.004}{10.9091}{mi}{58}\RG{25.016}{0.004}{10.9091}{mi}{58}\RG{29.848}{0.004}{10.9091}{mi}{58}\RG{34.692}{0.004}{10.9091}{mi}{59}\RG{39.536}{0.004}{10.9091}{mi}{21}\RG{45.945}{-1.632}{7.9701}{mi}{110}\RG{51.581}{0.004}{10.9091}{rm}{41}%
}\hskip.1em\allowbreak{} 。\allowbreak{}\hskip.1em \RB{8.366}{7.999}{0.500}{%
\RG{0.000}{0.004}{10.9091}{am}{3}%
}\hskip.1em\allowbreak{}

小结：先用对角相似变换平衡上下对角线，再归结为实对称矩阵；正弦向量则给出显式对角化。注意这里不能把“复对称矩阵”直接当作可对角化矩阵。

% Original page 3.
\section*{2015S03不可约多项式与循环向量}
题目设

\[
\RB{224.233}{10.305}{3.227}{%
\RG{41.095}{4.505}{7.9701}{mi}{110}\RG{76.799}{4.505}{7.9701}{mi}{110}\RG{81.937}{4.505}{7.9701}{sy}{0}\RG{88.524}{4.505}{7.9701}{rm}{49}\RG{0.000}{0.000}{10.9091}{mi}{103}\RG{5.595}{0.000}{10.9091}{rm}{40}\RG{9.837}{0.000}{10.9091}{mi}{120}\RG{16.072}{0.000}{10.9091}{rm}{41}\RG{23.327}{0.000}{10.9091}{rm}{61}\RG{34.860}{0.000}{10.9091}{mi}{120}\RG{49.155}{0.000}{10.9091}{rm}{43}\RG{60.064}{0.000}{10.9091}{mi}{97}\RG{70.564}{0.000}{10.9091}{mi}{120}\RG{95.680}{0.000}{10.9091}{rm}{43}\RG{106.589}{0.000}{10.9091}{sy}{1}\RG{111.422}{0.000}{10.9091}{sy}{1}\RG{116.265}{0.000}{10.9091}{sy}{1}\RG{121.741}{0.000}{10.9091}{rm}{43}\RG{132.650}{0.000}{10.9091}{mi}{97}\RG{154.873}{0.000}{10.9091}{mi}{120}\RG{163.532}{0.000}{10.9091}{rm}{43}\RG{174.441}{0.000}{10.9091}{mi}{97}\RG{188.874}{0.000}{10.9091}{sy}{50}\RG{199.178}{0.000}{10.9091}{mi}{75}\RG{209.224}{0.000}{10.9091}{rm}{91}\RG{212.254}{0.000}{10.9091}{mi}{120}\RG{218.489}{0.000}{10.9091}{rm}{93}\RG{221.519}{0.000}{10.9091}{mi}{59}\RG{65.832}{-1.636}{7.9701}{rm}{49}\RG{138.416}{-1.636}{7.9701}{mi}{110}\RG{143.554}{-1.636}{7.9701}{sy}{0}\RG{150.141}{-1.636}{7.9701}{rm}{49}\RG{180.208}{-1.636}{7.9701}{mi}{110}%
}
\]
\allowbreak{}\hskip.1em \RB{8.889}{7.955}{0.736}{%
\RG{0.000}{0.004}{10.9091}{mi}{86}%
}\hskip.1em\allowbreak{} 是数域\allowbreak{}\hskip.1em \RB{10.198}{7.955}{0.500}{%
\RG{0.000}{0.004}{10.9091}{mi}{75}%
}\hskip.1em\allowbreak{} 上的\allowbreak{}\hskip.1em \RB{6.729}{5.326}{0.616}{%
\RG{0.000}{0.004}{10.9091}{mi}{110}%
}\hskip.1em\allowbreak{} 维线性空间，\allowbreak{}\hskip.1em \RB{7.253}{5.326}{2.874}{%
\RG{0.000}{0.004}{10.9091}{mi}{39}%
}\hskip.1em\allowbreak{} 是\allowbreak{}\hskip.1em \RB{8.889}{7.955}{0.736}{%
\RG{0.000}{0.004}{10.9091}{mi}{86}%
}\hskip.1em\allowbreak{} 上的线性变换。向量\allowbreak{}\hskip.1em \RB{107.350}{8.315}{2.852}{%
\RG{0.000}{0.004}{10.9091}{mi}{11}\RG{6.978}{-1.632}{7.9701}{rm}{49}\RG{14.741}{0.004}{10.9091}{sy}{54}\RG{14.741}{0.004}{10.9091}{rm}{61}\RG{26.239}{0.004}{10.9091}{rm}{48}\RG{31.710}{0.004}{10.9091}{mi}{59}\RG{36.543}{0.004}{10.9091}{mi}{11}\RG{43.537}{-1.632}{7.9701}{rm}{50}\RG{48.270}{0.004}{10.9091}{mi}{59}\RG{53.103}{0.004}{10.9091}{mi}{58}\RG{57.947}{0.004}{10.9091}{mi}{58}\RG{62.780}{0.004}{10.9091}{mi}{58}\RG{67.623}{0.004}{10.9091}{mi}{59}\RG{72.467}{0.004}{10.9091}{mi}{11}\RG{79.492}{-1.632}{7.9701}{mi}{110}\RG{88.158}{0.004}{10.9091}{sy}{50}\RG{98.461}{0.004}{10.9091}{mi}{86}%
}\hskip.1em\allowbreak{} 满足

\[
\RB{186.751}{26.982}{3.226}{%
\RG{2.255}{18.300}{10.9091}{mi}{39}\RG{9.392}{18.300}{10.9091}{rm}{40}\RG{13.634}{18.300}{10.9091}{mi}{11}\RG{23.994}{18.300}{10.9091}{rm}{41}\RG{31.249}{18.300}{10.9091}{rm}{61}\RG{42.782}{18.300}{10.9091}{mi}{11}\RG{74.872}{18.300}{10.9091}{rm}{40}\RG{79.104}{18.300}{10.9091}{rm}{49}\RG{87.600}{18.300}{10.9091}{sy}{20}\RG{99.115}{18.300}{10.9091}{mi}{105}\RG{105.889}{18.300}{10.9091}{mi}{60}\RG{117.399}{18.300}{10.9091}{mi}{110}\RG{123.966}{18.300}{10.9091}{rm}{41}\RG{128.209}{18.300}{10.9091}{mi}{59}\RG{20.613}{16.664}{7.9701}{mi}{105}\RG{49.760}{16.664}{7.9701}{mi}{105}\RG{52.643}{16.664}{7.9701}{rm}{43}\RG{59.227}{16.664}{7.9701}{rm}{49}\RG{0.000}{0.001}{10.9091}{mi}{39}\RG{7.137}{0.001}{10.9091}{rm}{40}\RG{11.379}{0.001}{10.9091}{mi}{11}\RG{23.994}{0.001}{10.9091}{rm}{41}\RG{31.249}{0.001}{10.9091}{rm}{61}\RG{42.782}{0.001}{10.9091}{sy}{0}\RG{51.267}{0.001}{10.9091}{mi}{97}\RG{62.670}{0.001}{10.9091}{mi}{11}\RG{76.805}{0.001}{10.9091}{sy}{0}\RG{87.714}{0.001}{10.9091}{mi}{97}\RG{109.938}{0.001}{10.9091}{mi}{11}\RG{124.073}{0.001}{10.9091}{sy}{0}\RG{134.971}{0.001}{10.9091}{sy}{1}\RG{139.804}{0.001}{10.9091}{sy}{1}\RG{144.648}{0.001}{10.9091}{sy}{1}\RG{150.091}{0.001}{10.9091}{sy}{0}\RG{161.043}{0.001}{10.9091}{mi}{97}\RG{171.541}{0.001}{10.9091}{mi}{11}\RG{184.157}{0.001}{10.9091}{mi}{58}\RG{18.358}{-1.636}{7.9701}{mi}{110}\RG{57.033}{-1.636}{7.9701}{mi}{110}\RG{69.649}{-1.636}{7.9701}{rm}{49}\RG{93.481}{-1.636}{7.9701}{mi}{110}\RG{98.619}{-1.636}{7.9701}{sy}{0}\RG{105.206}{-1.636}{7.9701}{rm}{49}\RG{116.917}{-1.636}{7.9701}{rm}{50}\RG{166.809}{-1.636}{7.9701}{rm}{49}\RG{178.520}{-1.636}{7.9701}{mi}{110}%
}
\]
证明：若\allowbreak{}\hskip.1em \RB{19.725}{8.686}{3.223}{%
\RG{0.000}{0.004}{10.9091}{mi}{103}\RG{5.595}{0.004}{10.9091}{rm}{40}\RG{9.837}{0.004}{10.9091}{mi}{120}\RG{16.072}{0.004}{10.9091}{rm}{41}%
}\hskip.1em\allowbreak{} 在\allowbreak{}\hskip.1em \RB{10.198}{7.955}{0.500}{%
\RG{0.000}{0.004}{10.9091}{mi}{75}%
}\hskip.1em\allowbreak{} 上不可约，则\allowbreak{}\hskip.1em \RB{59.182}{8.686}{3.223}{%
\RG{0.000}{0.004}{10.9091}{sy}{102}\RG{5.455}{0.004}{10.9091}{mi}{11}\RG{12.433}{-1.632}{7.9701}{rm}{49}\RG{17.166}{0.004}{10.9091}{mi}{59}\RG{21.999}{0.004}{10.9091}{mi}{58}\RG{26.843}{0.004}{10.9091}{mi}{58}\RG{31.675}{0.004}{10.9091}{mi}{58}\RG{36.519}{0.004}{10.9091}{mi}{59}\RG{41.363}{0.004}{10.9091}{mi}{11}\RG{48.387}{-1.632}{7.9701}{mi}{110}\RG{54.023}{0.004}{10.9091}{sy}{103}%
}\hskip.1em\allowbreak{} 是\allowbreak{}\hskip.1em \RB{8.889}{7.955}{0.736}{%
\RG{0.000}{0.004}{10.9091}{mi}{86}%
}\hskip.1em\allowbreak{} 的一组基。举例说明：若\allowbreak{}\hskip.1em \RB{19.725}{8.686}{3.223}{%
\RG{0.000}{0.004}{10.9091}{mi}{103}\RG{5.595}{0.004}{10.9091}{rm}{40}\RG{9.837}{0.004}{10.9091}{mi}{120}\RG{16.072}{0.004}{10.9091}{rm}{41}%
}\hskip.1em\allowbreak{} 可约，则上述结论一般不成立。

解：由递推关系，

\[
\RB{224.720}{10.307}{3.760}{%
\RG{33.043}{4.507}{7.9701}{mi}{106}\RG{36.927}{4.506}{7.9701}{sy}{0}\RG{43.513}{4.507}{7.9701}{rm}{49}\RG{0.000}{0.002}{10.9091}{mi}{11}\RG{14.391}{0.002}{10.9091}{rm}{61}\RG{25.906}{0.002}{10.9091}{mi}{39}\RG{48.245}{0.002}{10.9091}{rm}{40}\RG{52.488}{0.002}{10.9091}{mi}{11}\RG{64.200}{0.002}{10.9091}{rm}{41}\RG{79.331}{0.002}{10.9091}{rm}{40}\RG{83.563}{0.002}{10.9091}{rm}{49}\RG{92.078}{0.002}{10.9091}{sy}{20}\RG{103.594}{0.002}{10.9091}{mi}{106}\RG{111.741}{0.002}{10.9091}{sy}{20}\RG{123.256}{0.002}{10.9091}{mi}{110}\RG{129.804}{0.002}{10.9091}{rm}{41}\RG{134.046}{0.002}{10.9091}{mi}{59}\RG{160.697}{0.002}{10.9091}{mi}{103}\RG{166.308}{0.002}{10.9091}{rm}{40}\RG{170.550}{0.002}{10.9091}{mi}{39}\RG{177.686}{0.002}{10.9091}{rm}{41}\RG{181.919}{0.002}{10.9091}{rm}{40}\RG{186.172}{0.002}{10.9091}{mi}{11}\RG{197.883}{0.002}{10.9091}{rm}{41}\RG{205.137}{0.002}{10.9091}{rm}{61}\RG{216.646}{0.002}{10.9091}{rm}{48}\RG{222.126}{0.002}{10.9091}{mi}{58}\RG{6.979}{-1.634}{7.9701}{mi}{106}\RG{59.466}{-1.634}{7.9701}{rm}{49}\RG{193.151}{-1.634}{7.9701}{rm}{49}%
}
\]
若\allowbreak{}\hskip.1em \RB{48.310}{5.318}{2.609}{%
\RG{0.000}{-0.004}{10.9091}{mi}{11}\RG{6.979}{-1.641}{7.9701}{rm}{49}\RG{11.711}{-0.004}{10.9091}{mi}{59}\RG{16.544}{-0.004}{10.9091}{mi}{58}\RG{21.388}{-0.004}{10.9091}{mi}{58}\RG{26.231}{-0.004}{10.9091}{mi}{58}\RG{31.064}{-0.004}{10.9091}{mi}{59}\RG{35.908}{-0.004}{10.9091}{mi}{11}\RG{42.932}{-1.641}{7.9701}{mi}{110}%
}\hskip.1em\allowbreak{} 线性相关，则存在不全为零的\allowbreak{}\hskip.1em \RB{78.406}{7.947}{2.609}{%
\RG{0.000}{-0.004}{10.9091}{mi}{99}\RG{4.721}{-1.641}{7.9701}{rm}{48}\RG{9.453}{-0.004}{10.9091}{mi}{59}\RG{14.286}{-0.004}{10.9091}{mi}{58}\RG{19.130}{-0.004}{10.9091}{mi}{58}\RG{23.973}{-0.004}{10.9091}{mi}{58}\RG{28.806}{-0.004}{10.9091}{mi}{59}\RG{33.650}{-0.004}{10.9091}{mi}{99}\RG{38.417}{-1.641}{7.9701}{mi}{110}\RG{43.555}{-1.640}{7.9701}{sy}{0}\RG{50.141}{-1.641}{7.9701}{rm}{49}\RG{57.904}{-0.004}{10.9091}{sy}{50}\RG{68.208}{-0.004}{10.9091}{mi}{75}%
}\hskip.1em\allowbreak{} ，使得

\[
\RB{81.189}{19.437}{15.213}{%
\RG{0.000}{13.637}{7.9701}{mi}{110}\RG{5.138}{13.637}{7.9701}{sy}{0}\RG{11.724}{13.637}{7.9701}{rm}{49}\RG{0.100}{10.364}{10.9092}{ex}{88}\RG{34.016}{4.505}{7.9701}{mi}{106}\RG{17.777}{0.000}{10.9091}{mi}{99}\RG{26.880}{0.000}{10.9091}{mi}{39}\RG{38.398}{0.000}{10.9091}{rm}{40}\RG{42.641}{0.000}{10.9091}{mi}{11}\RG{54.352}{0.000}{10.9091}{rm}{41}\RG{61.607}{0.000}{10.9091}{rm}{61}\RG{73.116}{0.000}{10.9091}{rm}{48}\RG{78.595}{0.000}{10.9091}{mi}{58}\RG{22.498}{-1.636}{7.9701}{mi}{106}\RG{49.619}{-1.636}{7.9701}{rm}{49}\RG{0.627}{-13.087}{7.9701}{mi}{106}\RG{4.511}{-13.087}{7.9701}{rm}{61}\RG{11.094}{-13.087}{7.9701}{rm}{48}%
}
\]
令\allowbreak{}\hskip.1em \RB{84.471}{10.831}{5.347}{%
\RG{0.000}{0.004}{10.9091}{mi}{104}\RG{6.286}{0.004}{10.9091}{rm}{40}\RG{10.528}{0.004}{10.9091}{mi}{120}\RG{16.763}{0.004}{10.9091}{rm}{41}\RG{24.018}{0.004}{10.9091}{rm}{61}\RG{35.551}{8.186}{10.9092}{ex}{80}\RG{47.066}{5.031}{7.9701}{mi}{110}\RG{47.066}{-3.221}{7.9701}{mi}{106}\RG{50.950}{-3.221}{7.9701}{rm}{61}\RG{52.204}{5.031}{7.9701}{sy}{0}\RG{57.534}{-3.221}{7.9701}{rm}{48}\RG{58.791}{5.031}{7.9701}{rm}{49}\RG{65.342}{0.004}{10.9091}{mi}{99}\RG{70.063}{-1.632}{7.9701}{mi}{106}\RG{74.445}{0.004}{10.9091}{mi}{120}\RG{80.680}{3.963}{7.9701}{mi}{106}%
}\hskip.1em\allowbreak{} ，则\allowbreak{}\hskip.1em \RB{26.333}{8.315}{2.852}{%
\RG{0.000}{0.004}{10.9091}{mi}{104}\RG{9.316}{0.004}{10.9091}{sy}{54}\RG{9.316}{0.004}{10.9091}{rm}{61}\RG{20.814}{0.004}{10.9091}{rm}{48}%
}\hskip.1em\allowbreak{} 、deg\allowbreak{}\hskip.1em \RB{27.544}{8.075}{0.932}{%
\RG{0.000}{0.004}{10.9091}{mi}{104}\RG{9.306}{0.004}{10.9091}{mi}{60}\RG{20.815}{0.004}{10.9091}{mi}{110}%
}\hskip.1em\allowbreak{} ，且\allowbreak{}\hskip.1em \RB{62.142}{8.686}{3.223}{%
\RG{0.000}{0.004}{10.9091}{mi}{104}\RG{6.285}{0.004}{10.9091}{rm}{40}\RG{10.528}{0.004}{10.9091}{mi}{39}\RG{17.664}{0.004}{10.9091}{rm}{41}\RG{21.897}{0.004}{10.9091}{rm}{40}\RG{26.149}{0.004}{10.9091}{mi}{11}\RG{33.128}{-1.632}{7.9701}{rm}{49}\RG{37.860}{0.004}{10.9091}{rm}{41}\RG{45.115}{0.004}{10.9091}{rm}{61}\RG{56.624}{0.004}{10.9091}{rm}{48}%
}\hskip.1em\allowbreak{} 。

因为\allowbreak{}\hskip.1em \RB{5.671}{5.325}{2.733}{%
\RG{0.000}{0.003}{10.9091}{mi}{103}%
}\hskip.1em\allowbreak{} 是\allowbreak{}\hskip.1em \RB{6.729}{5.325}{0.617}{%
\RG{0.000}{0.003}{10.9091}{mi}{110}%
}\hskip.1em\allowbreak{} 次不可约多项式，故gcd\allowbreak{}\hskip.1em \RB{44.806}{8.685}{3.224}{%
\RG{0.000}{0.003}{10.9091}{rm}{40}\RG{4.242}{0.003}{10.9091}{mi}{103}\RG{9.817}{0.003}{10.9091}{mi}{59}\RG{14.660}{0.003}{10.9091}{mi}{104}\RG{20.971}{0.003}{10.9091}{rm}{41}\RG{28.226}{0.003}{10.9091}{rm}{61}\RG{39.735}{0.003}{10.9091}{rm}{49}%
}\hskip.1em\allowbreak{} 。由Bézout等式，存在\allowbreak{}\hskip.1em \RB{51.652}{8.685}{3.224}{%
\RG{0.000}{0.003}{10.9091}{mi}{117}\RG{6.240}{0.003}{10.9091}{mi}{59}\RG{11.073}{0.003}{10.9091}{mi}{118}\RG{19.803}{0.003}{10.9091}{sy}{50}\RG{30.106}{0.003}{10.9091}{mi}{75}\RG{40.152}{0.003}{10.9091}{rm}{91}\RG{43.182}{0.003}{10.9091}{mi}{120}\RG{49.417}{0.003}{10.9091}{rm}{93}%
}\hskip.1em\allowbreak{} ，使

\[
\RB{118.612}{8.686}{3.223}{%
\RG{0.000}{0.004}{10.9091}{mi}{117}\RG{6.245}{0.004}{10.9091}{rm}{40}\RG{10.487}{0.004}{10.9091}{mi}{120}\RG{16.722}{0.004}{10.9091}{rm}{41}\RG{20.964}{0.004}{10.9091}{mi}{103}\RG{26.559}{0.004}{10.9091}{rm}{40}\RG{30.801}{0.004}{10.9091}{mi}{120}\RG{37.036}{0.004}{10.9091}{rm}{41}\RG{43.691}{0.004}{10.9091}{rm}{43}\RG{54.612}{0.004}{10.9091}{mi}{118}\RG{60.292}{0.004}{10.9091}{rm}{40}\RG{64.535}{0.004}{10.9091}{mi}{120}\RG{70.770}{0.004}{10.9091}{rm}{41}\RG{75.012}{0.004}{10.9091}{mi}{104}\RG{81.297}{0.004}{10.9091}{rm}{40}\RG{85.540}{0.004}{10.9091}{mi}{120}\RG{91.775}{0.004}{10.9091}{rm}{41}\RG{99.029}{0.004}{10.9091}{rm}{61}\RG{110.539}{0.004}{10.9091}{rm}{49}\RG{116.017}{0.004}{10.9091}{mi}{58}%
}
\]
代入\allowbreak{}\hskip.1em \RB{28.033}{5.326}{2.874}{%
\RG{0.000}{0.004}{10.9091}{mi}{120}\RG{9.265}{0.004}{10.9091}{rm}{61}\RG{20.780}{0.004}{10.9091}{mi}{39}%
}\hskip.1em\allowbreak{} ，再作用于\allowbreak{}\hskip.1em \RB{11.042}{5.326}{2.132}{%
\RG{0.000}{0.004}{10.9091}{mi}{11}\RG{6.979}{-1.632}{7.9701}{rm}{49}%
}\hskip.1em\allowbreak{} ，得到

\[
\RB{188.986}{8.677}{3.232}{%
\RG{0.000}{-0.005}{10.9091}{mi}{11}\RG{14.741}{-0.005}{10.9091}{rm}{61}\RG{26.256}{-0.005}{10.9091}{mi}{117}\RG{32.501}{-0.005}{10.9091}{rm}{40}\RG{36.743}{-0.005}{10.9091}{mi}{39}\RG{43.880}{-0.005}{10.9091}{rm}{41}\RG{48.122}{-0.005}{10.9091}{mi}{103}\RG{53.717}{-0.005}{10.9091}{rm}{40}\RG{57.959}{-0.005}{10.9091}{mi}{39}\RG{65.097}{-0.005}{10.9091}{rm}{41}\RG{69.329}{-0.005}{10.9091}{rm}{40}\RG{73.582}{-0.005}{10.9091}{mi}{11}\RG{85.292}{-0.005}{10.9091}{rm}{41}\RG{91.946}{-0.005}{10.9091}{rm}{43}\RG{102.868}{-0.005}{10.9091}{mi}{118}\RG{108.547}{-0.005}{10.9091}{rm}{40}\RG{112.790}{-0.005}{10.9091}{mi}{39}\RG{119.926}{-0.005}{10.9091}{rm}{41}\RG{124.169}{-0.005}{10.9091}{mi}{104}\RG{130.454}{-0.005}{10.9091}{rm}{40}\RG{134.697}{-0.005}{10.9091}{mi}{39}\RG{141.833}{-0.005}{10.9091}{rm}{41}\RG{146.066}{-0.005}{10.9091}{rm}{40}\RG{150.318}{-0.005}{10.9091}{mi}{11}\RG{162.029}{-0.005}{10.9091}{rm}{41}\RG{169.284}{-0.005}{10.9091}{rm}{61}\RG{180.793}{-0.005}{10.9091}{rm}{48}\RG{186.272}{-0.005}{10.9091}{mi}{59}\RG{6.978}{-1.642}{7.9701}{rm}{49}\RG{80.560}{-1.642}{7.9701}{rm}{49}\RG{157.296}{-1.642}{7.9701}{rm}{49}%
}
\]
与\allowbreak{}\hskip.1em \RB{31.759}{8.306}{2.861}{%
\RG{0.000}{-0.005}{10.9091}{mi}{11}\RG{6.979}{-1.641}{7.9701}{rm}{49}\RG{14.742}{-0.005}{10.9091}{sy}{54}\RG{14.742}{-0.005}{10.9091}{rm}{61}\RG{26.240}{-0.005}{10.9091}{rm}{48}%
}\hskip.1em\allowbreak{} 矛盾。因此上述\allowbreak{}\hskip.1em \RB{6.729}{5.317}{0.625}{%
\RG{0.000}{-0.005}{10.9091}{mi}{110}%
}\hskip.1em\allowbreak{} 个向量线性无关，构成\allowbreak{}\hskip.1em \RB{8.889}{7.946}{0.745}{%
\RG{0.000}{-0.005}{10.9091}{mi}{86}%
}\hskip.1em\allowbreak{} 的一组基。

\noindent \RB{384.748}{32.035}{14.122}{%
\RT{0.000}{21.935}{10.9091}{反}\RT{10.909}{21.935}{10.9091}{例}\RT{36.360}{21.935}{10.9091}{取}\RG{357.607}{18.658}{10.9092}{ex}{32}\RG{371.698}{18.658}{10.9092}{ex}{33}\RG{366.244}{7.324}{10.9091}{rm}{49}\RG{154.684}{4.508}{7.9701}{rm}{50}\RG{216.270}{4.508}{7.9701}{rm}{50}\RG{74.033}{0.004}{10.9091}{mi}{75}\RG{87.109}{0.004}{10.9091}{rm}{61}\RG{98.624}{0.004}{10.9091}{bb}{81}\RG{107.109}{0.004}{10.9091}{mi}{59}\RG{122.851}{0.004}{10.9091}{mi}{86}\RG{134.684}{0.004}{10.9091}{rm}{61}\RG{146.199}{0.004}{10.9091}{bb}{81}\RG{159.417}{0.004}{10.9091}{mi}{59}\RG{175.159}{0.004}{10.9091}{mi}{103}\RG{180.769}{0.004}{10.9091}{rm}{40}\RG{185.012}{0.004}{10.9091}{mi}{120}\RG{191.247}{0.004}{10.9091}{rm}{41}\RG{198.501}{0.004}{10.9091}{rm}{61}\RG{210.035}{0.004}{10.9091}{mi}{120}\RG{223.426}{0.004}{10.9091}{sy}{0}\RG{234.335}{0.004}{10.9091}{rm}{49}\RG{239.790}{0.004}{10.9091}{mi}{59}\RG{255.532}{0.004}{10.9091}{mi}{39}\RG{265.714}{0.004}{10.9091}{rm}{61}\RG{277.229}{0.004}{10.9091}{mi}{73}\RG{289.336}{0.004}{10.9091}{mi}{59}\RG{305.078}{0.004}{10.9091}{mi}{11}\RG{319.836}{0.004}{10.9091}{rm}{61}\RG{331.351}{0.004}{10.9091}{mi}{11}\RG{346.092}{0.004}{10.9091}{rm}{61}\RG{382.153}{0.004}{10.9091}{mi}{58}\RG{282.025}{-1.685}{7.9701}{mi}{86}\RG{312.073}{-1.633}{7.9701}{rm}{49}\RG{338.330}{-1.633}{7.9701}{rm}{50}\RG{366.244}{-7.987}{10.9091}{rm}{48}%
}

此时\allowbreak{}\hskip.1em \RB{73.948}{7.762}{2.609}{%
\RG{0.000}{-0.004}{10.9091}{mi}{97}\RG{5.766}{-1.640}{7.9701}{rm}{49}\RG{13.529}{-0.004}{10.9091}{rm}{61}\RG{25.027}{-0.004}{10.9091}{rm}{48}\RG{30.498}{-0.004}{10.9091}{mi}{59}\RG{35.331}{-0.004}{10.9091}{mi}{97}\RG{41.113}{-1.640}{7.9701}{rm}{50}\RG{48.876}{-0.004}{10.9091}{rm}{61}\RG{60.391}{-0.004}{10.9091}{sy}{0}\RG{68.877}{-0.004}{10.9091}{rm}{49}%
}\hskip.1em\allowbreak{} ，因而

\[
\RB{217.222}{8.678}{3.231}{%
\RG{0.000}{-0.004}{10.9091}{mi}{39}\RG{7.136}{-0.004}{10.9091}{rm}{40}\RG{11.379}{-0.004}{10.9091}{mi}{11}\RG{23.089}{-0.004}{10.9091}{rm}{41}\RG{30.344}{-0.004}{10.9091}{rm}{61}\RG{41.877}{-0.004}{10.9091}{mi}{11}\RG{53.588}{-0.004}{10.9091}{mi}{59}\RG{80.240}{-0.004}{10.9091}{mi}{39}\RG{87.392}{-0.004}{10.9091}{rm}{40}\RG{91.635}{-0.004}{10.9091}{mi}{11}\RG{103.345}{-0.004}{10.9091}{rm}{41}\RG{110.599}{-0.004}{10.9091}{rm}{61}\RG{122.133}{-0.004}{10.9091}{mi}{11}\RG{136.874}{-0.004}{10.9091}{rm}{61}\RG{148.390}{-0.004}{10.9091}{sy}{0}\RG{156.874}{-0.004}{10.9091}{mi}{97}\RG{167.373}{-0.004}{10.9091}{mi}{11}\RG{181.508}{-0.004}{10.9091}{sy}{0}\RG{192.418}{-0.004}{10.9091}{mi}{97}\RG{202.917}{-0.004}{10.9091}{mi}{11}\RG{214.628}{-0.004}{10.9091}{mi}{58}\RG{18.357}{-1.640}{7.9701}{rm}{49}\RG{48.856}{-1.640}{7.9701}{rm}{50}\RG{98.613}{-1.640}{7.9701}{rm}{50}\RG{129.112}{-1.640}{7.9701}{rm}{49}\RG{162.641}{-1.640}{7.9701}{rm}{50}\RG{174.352}{-1.640}{7.9701}{rm}{49}\RG{198.185}{-1.640}{7.9701}{rm}{49}\RG{209.896}{-1.640}{7.9701}{rm}{50}%
}
\]
题设关系全部成立，但\allowbreak{}\hskip.1em \RB{27.840}{5.319}{2.608}{%
\RG{0.000}{-0.003}{10.9091}{mi}{11}\RG{6.979}{-1.639}{7.9701}{rm}{49}\RG{11.711}{-0.003}{10.9091}{mi}{59}\RG{16.544}{-0.003}{10.9091}{mi}{11}\RG{23.538}{-1.639}{7.9701}{rm}{50}%
}\hskip.1em\allowbreak{} 线性相关。\allowbreak{}\hskip.1em \RB{8.366}{7.992}{0.503}{%
\RG{0.000}{-0.003}{10.9091}{am}{3}%
}\hskip.1em\allowbreak{}

小结：将向量的线性关系写成\allowbreak{}\hskip.1em \RB{49.200}{7.971}{2.992}{%
\RG{0.000}{-0.001}{10.9091}{mi}{104}\RG{5.740}{-0.001}{10.9091}{rm}{40}\RG{9.614}{-0.001}{10.9091}{mi}{39}\RG{16.131}{-0.001}{10.9091}{rm}{41}\RG{20.006}{-0.001}{10.9091}{mi}{11}\RG{26.379}{-1.495}{7.9701}{rm}{49}\RG{33.616}{-0.001}{10.9091}{rm}{61}\RG{44.117}{-0.001}{10.9091}{rm}{48}%
}\hskip.1em\allowbreak{} ，再用不可约性与Bézout等式排除低次关系。题设起初只给出\allowbreak{}\hskip.1em \RB{48.569}{7.972}{2.991}{%
\RG{0.000}{0.000}{10.9091}{mi}{103}\RG{5.110}{0.000}{10.9091}{rm}{40}\RG{8.984}{0.000}{10.9091}{mi}{39}\RG{15.501}{0.000}{10.9091}{rm}{41}\RG{19.376}{0.000}{10.9091}{mi}{11}\RG{25.749}{-1.495}{7.9701}{rm}{49}\RG{32.985}{0.000}{10.9091}{rm}{61}\RG{43.486}{0.000}{10.9091}{rm}{48}%
}\hskip.1em\allowbreak{} ，不能未经证明就写成\allowbreak{}\hskip.1em \RB{37.720}{7.972}{2.991}{%
\RG{0.000}{0.000}{10.9091}{mi}{103}\RG{5.109}{0.000}{10.9091}{rm}{40}\RG{8.983}{0.000}{10.9091}{mi}{39}\RG{15.501}{0.000}{10.9091}{rm}{41}\RG{22.126}{0.000}{10.9091}{rm}{61}\RG{32.637}{0.000}{10.9091}{rm}{48}%
}\hskip.1em\allowbreak{} 。

% Original page 4.
\section*{2015S04列满秩条件与不变子空间}
题目设\allowbreak{}\hskip.1em \RB{8.366}{8.312}{0.500}{%
\RG{0.000}{0.001}{10.9091}{mi}{65}%
}\hskip.1em\allowbreak{} 为\allowbreak{}\hskip.1em \RB{6.729}{5.323}{0.619}{%
\RG{0.000}{0.001}{10.9091}{mi}{110}%
}\hskip.1em\allowbreak{} 阶方阵，\allowbreak{}\hskip.1em \RB{8.791}{8.192}{0.739}{%
\RG{0.000}{0.001}{10.9091}{mi}{67}%
}\hskip.1em\allowbreak{} 为\allowbreak{}\hskip.1em \RB{26.085}{8.072}{0.619}{%
\RG{0.000}{0.001}{10.9091}{mi}{107}\RG{8.447}{0.001}{10.9091}{sy}{2}\RG{19.356}{0.001}{10.9091}{mi}{110}%
}\hskip.1em\allowbreak{} 矩阵。对任意\allowbreak{}\hskip.1em \RB{27.879}{8.181}{2.855}{%
\RG{0.000}{0.001}{10.9091}{mi}{22}\RG{9.603}{0.001}{10.9091}{sy}{50}\RG{19.906}{0.001}{10.9091}{bb}{67}%
}\hskip.1em\allowbreak{} ，矩阵

\[
\RB{53.712}{12.274}{21.442}{%
\RG{0.000}{11.338}{10.9092}{ex}{32}\RG{47.157}{11.338}{10.9092}{ex}{33}\RG{8.637}{0.004}{10.9091}{mi}{65}\RG{19.243}{0.004}{10.9091}{sy}{0}\RG{30.152}{0.004}{10.9091}{mi}{22}\RG{36.720}{0.004}{10.9091}{mi}{73}\RG{41.521}{-1.632}{7.9701}{mi}{110}\RG{23.608}{-15.307}{10.9091}{mi}{67}%
}
\]
均为列满秩阵。证明：对任意\allowbreak{}\hskip.1em \RB{27.670}{8.176}{0.940}{%
\RG{0.000}{-0.004}{10.9091}{mi}{21}\RG{9.394}{-0.004}{10.9091}{sy}{50}\RG{19.697}{-0.004}{10.9091}{bb}{67}%
}\hskip.1em\allowbreak{} ，矩阵

\[
\RB{130.731}{25.318}{54.208}{%
\RG{42.495}{24.392}{10.9092}{ex}{48}\RG{123.871}{24.392}{10.9092}{ex}{49}\RG{83.668}{15.308}{10.9091}{mi}{67}\RG{42.495}{5.192}{10.9092}{ex}{66}\RG{123.871}{5.192}{10.9092}{ex}{67}\RG{60.269}{-0.003}{10.9091}{mi}{67}\RG{68.847}{-0.003}{10.9091}{rm}{40}\RG{73.090}{-0.003}{10.9091}{mi}{65}\RG{83.696}{-0.003}{10.9091}{sy}{0}\RG{94.605}{-0.003}{10.9091}{mi}{21}\RG{100.965}{-0.003}{10.9091}{mi}{73}\RG{111.400}{-0.003}{10.9091}{rm}{41}\RG{42.495}{-1.354}{10.9092}{ex}{66}\RG{123.871}{-1.354}{10.9092}{ex}{67}\RG{105.764}{-1.639}{7.9701}{mi}{110}\RG{42.495}{-7.899}{10.9092}{ex}{66}\RG{123.871}{-7.899}{10.9092}{ex}{67}\RG{113.276}{-11.355}{7.9701}{rm}{50}\RG{42.495}{-14.445}{10.9092}{ex}{66}\RG{123.871}{-14.445}{10.9092}{ex}{67}\RG{57.903}{-15.314}{10.9091}{mi}{67}\RG{66.481}{-15.314}{10.9091}{rm}{40}\RG{70.723}{-15.314}{10.9091}{mi}{65}\RG{81.329}{-15.314}{10.9091}{sy}{0}\RG{92.239}{-15.314}{10.9091}{mi}{21}\RG{98.599}{-15.314}{10.9091}{mi}{73}\RG{109.034}{-15.314}{10.9091}{rm}{41}\RG{103.398}{-16.950}{7.9701}{mi}{110}\RG{0.000}{-17.172}{10.9091}{sy}{77}\RG{13.101}{-17.172}{10.9091}{rm}{40}\RG{17.343}{-17.172}{10.9091}{mi}{21}\RG{23.707}{-17.172}{10.9091}{rm}{41}\RG{30.962}{-17.172}{10.9091}{rm}{61}\RG{42.495}{-20.990}{10.9092}{ex}{66}\RG{123.871}{-20.990}{10.9092}{ex}{67}\RT{86.440}{-27.041}{10.9091}{.}\RG{42.495}{-27.536}{10.9092}{ex}{66}\RG{123.871}{-27.536}{10.9092}{ex}{67}\RT{86.440}{-31.026}{10.9091}{.}\RG{42.495}{-34.518}{10.9092}{ex}{64}\RG{123.871}{-34.518}{10.9092}{ex}{65}\RT{86.440}{-35.011}{10.9091}{.}\RG{107.414}{-46.363}{7.9701}{mi}{110}\RG{112.552}{-46.363}{7.9701}{sy}{0}\RG{119.139}{-46.363}{7.9701}{rm}{49}\RG{52.040}{-50.322}{10.9091}{mi}{67}\RG{60.618}{-50.322}{10.9091}{rm}{40}\RG{64.861}{-50.322}{10.9091}{mi}{65}\RG{75.467}{-50.322}{10.9091}{sy}{0}\RG{86.376}{-50.322}{10.9091}{mi}{21}\RG{92.736}{-50.322}{10.9091}{mi}{73}\RG{103.172}{-50.322}{10.9091}{rm}{41}\RG{97.535}{-51.958}{7.9701}{mi}{110}%
}
\]
均为列满秩阵。解：任意固定\allowbreak{}\hskip.1em \RB{27.670}{8.183}{0.933}{%
\RG{0.000}{0.003}{10.9091}{mi}{21}\RG{9.394}{0.003}{10.9091}{sy}{50}\RG{19.697}{0.003}{10.9091}{bb}{67}%
}\hskip.1em\allowbreak{} ，记\allowbreak{}\hskip.1em \RB{61.420}{8.314}{2.213}{%
\RG{0.000}{0.003}{10.9091}{mi}{66}\RG{11.852}{0.003}{10.9091}{rm}{61}\RG{23.368}{0.003}{10.9091}{mi}{65}\RG{33.974}{0.003}{10.9091}{sy}{0}\RG{44.883}{0.003}{10.9091}{mi}{21}\RG{51.243}{0.003}{10.9091}{mi}{73}\RG{56.042}{-1.633}{7.9701}{mi}{110}%
}\hskip.1em\allowbreak{} 。假设\allowbreak{}\hskip.1em \RB{27.360}{8.685}{3.224}{%
\RG{0.000}{0.003}{10.9091}{sy}{77}\RG{13.101}{0.003}{10.9091}{rm}{40}\RG{17.343}{0.003}{10.9091}{mi}{21}\RG{23.707}{0.003}{10.9091}{rm}{41}%
}\hskip.1em\allowbreak{} 不列满秩，则存在\allowbreak{}\hskip.1em \RB{52.824}{8.314}{2.853}{%
\RG{0.000}{0.003}{10.9091}{rm}{48}\RG{8.484}{0.003}{10.9091}{sy}{54}\RG{8.484}{0.003}{10.9091}{rm}{61}\RG{20.000}{0.003}{10.9091}{mi}{120}\RG{29.265}{0.003}{10.9091}{sy}{50}\RG{39.568}{0.003}{10.9091}{bb}{67}\RG{47.446}{3.962}{7.9701}{mi}{110}%
}\hskip.1em\allowbreak{} ，使

\[
\RB{162.277}{10.284}{2.733}{%
\RG{17.400}{4.508}{7.9701}{mi}{106}\RG{0.000}{0.003}{10.9091}{mi}{67}\RG{8.564}{0.003}{10.9091}{mi}{66}\RG{21.782}{0.003}{10.9091}{mi}{120}\RG{31.047}{0.003}{10.9091}{rm}{61}\RG{42.546}{0.003}{10.9091}{rm}{48}\RG{48.017}{0.003}{10.9091}{mi}{59}\RG{74.669}{0.003}{10.9091}{mi}{106}\RG{82.831}{0.003}{10.9091}{rm}{61}\RG{94.329}{0.003}{10.9091}{rm}{48}\RG{99.801}{0.003}{10.9091}{mi}{59}\RG{104.650}{0.003}{10.9091}{rm}{49}\RG{110.104}{0.003}{10.9091}{mi}{59}\RG{114.937}{0.003}{10.9091}{mi}{58}\RG{119.780}{0.003}{10.9091}{mi}{58}\RG{124.613}{0.003}{10.9091}{mi}{58}\RG{129.457}{0.003}{10.9091}{mi}{59}\RG{134.300}{0.003}{10.9091}{mi}{110}\RG{143.319}{0.003}{10.9091}{sy}{0}\RG{154.228}{0.003}{10.9091}{rm}{49}\RG{159.682}{0.003}{10.9091}{mi}{58}%
}
\]
由于空间维数为\allowbreak{}\hskip.1em \RB{6.729}{5.326}{0.616}{%
\RG{0.000}{0.004}{10.9091}{mi}{110}%
}\hskip.1em\allowbreak{} ，存在最小的\allowbreak{}\hskip.1em \RB{69.770}{8.686}{3.223}{%
\RG{0.000}{0.004}{10.9091}{mi}{109}\RG{12.609}{0.004}{10.9091}{sy}{50}\RG{22.896}{0.004}{10.9091}{sy}{102}\RG{28.366}{0.004}{10.9091}{rm}{49}\RG{33.821}{0.004}{10.9091}{mi}{59}\RG{38.654}{0.004}{10.9091}{mi}{58}\RG{43.498}{0.004}{10.9091}{mi}{58}\RG{48.330}{0.004}{10.9091}{mi}{58}\RG{53.174}{0.004}{10.9091}{mi}{59}\RG{58.018}{0.004}{10.9091}{mi}{110}\RG{64.612}{0.004}{10.9091}{sy}{103}%
}\hskip.1em\allowbreak{} ，满足

\[
\RB{182.230}{10.309}{3.223}{%
\RG{8.822}{4.509}{7.9701}{mi}{109}\RG{149.136}{4.509}{7.9701}{mi}{109}\RG{156.626}{4.508}{7.9701}{sy}{0}\RG{163.213}{4.509}{7.9701}{rm}{49}\RG{0.000}{0.004}{10.9091}{mi}{66}\RG{16.811}{0.004}{10.9091}{mi}{120}\RG{26.076}{0.004}{10.9091}{sy}{50}\RG{36.379}{0.004}{10.9091}{mi}{85}\RG{48.047}{0.004}{10.9091}{rm}{58}\RG{51.069}{0.004}{10.9091}{rm}{61}\RT{62.592}{0.004}{10.9091}{s}\RT{66.891}{0.004}{10.9091}{p}\RT{72.956}{0.004}{10.9091}{a}\RT{78.411}{0.004}{10.9091}{n}\RG{84.477}{0.004}{10.9091}{sy}{102}\RG{89.931}{0.004}{10.9091}{mi}{120}\RG{96.160}{0.004}{10.9091}{mi}{59}\RG{100.993}{0.004}{10.9091}{mi}{66}\RG{109.807}{0.004}{10.9091}{mi}{120}\RG{116.036}{0.004}{10.9091}{mi}{59}\RG{120.880}{0.004}{10.9091}{mi}{58}\RG{125.724}{0.004}{10.9091}{mi}{58}\RG{130.557}{0.004}{10.9091}{mi}{58}\RG{135.400}{0.004}{10.9091}{mi}{59}\RG{140.244}{0.004}{10.9091}{mi}{66}\RG{167.945}{0.004}{10.9091}{mi}{120}\RG{174.180}{0.004}{10.9091}{sy}{103}\RG{179.635}{0.004}{10.9091}{mi}{58}%
}
\]
由\allowbreak{}\hskip.1em \RB{9.751}{5.327}{0.615}{%
\RG{0.000}{0.005}{10.9091}{mi}{109}%
}\hskip.1em\allowbreak{} 的最小性，这\allowbreak{}\hskip.1em \RB{9.751}{5.327}{0.615}{%
\RG{0.000}{0.005}{10.9091}{mi}{109}%
}\hskip.1em\allowbreak{} 个向量线性无关。\allowbreak{}\hskip.1em \RB{8.747}{7.956}{0.500}{%
\RG{0.000}{0.005}{10.9091}{mi}{66}%
}\hskip.1em\allowbreak{} 将前\allowbreak{}\hskip.1em \RB{27.268}{7.770}{0.615}{%
\RG{0.000}{0.005}{10.9091}{mi}{109}\RG{11.645}{0.005}{10.9091}{sy}{0}\RG{22.197}{0.005}{10.9091}{rm}{49}%
}\hskip.1em\allowbreak{} 个生成元送到下一个生成元，又将最后一个生成元送入\allowbreak{}\hskip.1em \RB{8.791}{7.955}{0.736}{%
\RG{0.000}{0.004}{10.9091}{mi}{85}%
}\hskip.1em\allowbreak{} ，故

\[
\RB{178.096}{8.685}{3.224}{%
\RG{0.000}{0.003}{10.9091}{mi}{66}\RG{8.815}{0.003}{10.9091}{mi}{85}\RG{20.490}{0.003}{10.9091}{sy}{18}\RG{32.005}{0.003}{10.9091}{mi}{85}\RG{39.424}{0.003}{10.9091}{mi}{59}\RG{66.075}{0.003}{10.9091}{mi}{65}\RG{74.257}{0.003}{10.9091}{mi}{85}\RG{85.949}{0.003}{10.9091}{rm}{61}\RG{97.447}{0.003}{10.9091}{rm}{40}\RG{101.706}{0.003}{10.9091}{mi}{66}\RG{112.952}{0.003}{10.9091}{rm}{43}\RG{123.862}{0.003}{10.9091}{mi}{21}\RG{130.222}{0.003}{10.9091}{mi}{73}\RG{140.657}{0.003}{10.9091}{rm}{41}\RG{144.899}{0.003}{10.9091}{mi}{85}\RG{156.568}{0.003}{10.9091}{sy}{18}\RG{168.083}{0.003}{10.9091}{mi}{85}\RG{175.502}{0.003}{10.9091}{mi}{58}\RG{135.021}{-1.633}{7.9701}{mi}{110}%
}
\]
并且由\allowbreak{}\hskip.1em \RB{30.852}{7.435}{1.998}{%
\RG{0.000}{-0.003}{10.9091}{mi}{109}\RG{12.608}{-0.003}{10.9091}{sy}{20}\RG{24.123}{-0.003}{10.9091}{mi}{110}%
}\hskip.1em\allowbreak{} 和前面的等式，每个生成元都被\allowbreak{}\hskip.1em \RB{8.791}{8.188}{0.743}{%
\RG{0.000}{-0.003}{10.9091}{mi}{67}%
}\hskip.1em\allowbreak{} 零化，所以

\[
\RB{70.429}{9.290}{2.859}{%
\RG{0.000}{-0.003}{10.9091}{rm}{48}\RG{8.484}{-0.003}{10.9091}{sy}{54}\RG{8.484}{-0.003}{10.9091}{rm}{61}\RG{19.999}{-0.003}{10.9091}{mi}{85}\RG{31.668}{-0.003}{10.9091}{sy}{18}\RT{43.183}{-0.003}{10.9091}{k}\RT{48.943}{-0.003}{10.9091}{e}\RT{53.787}{-0.003}{10.9091}{r}\RG{59.881}{-0.003}{10.9091}{mi}{67}\RG{67.834}{-0.003}{10.9091}{mi}{58}%
}
\]
在非零有限维复空间\allowbreak{}\hskip.1em \RB{8.791}{7.949}{0.742}{%
\RG{0.000}{-0.002}{10.9091}{mi}{85}%
}\hskip.1em\allowbreak{} 上，限制变换\allowbreak{}\hskip.1em \RB{18.016}{8.680}{3.229}{%
\RG{0.000}{-0.002}{10.9091}{mi}{65}\RG{8.182}{-0.002}{10.9091}{sy}{106}\RG{11.212}{-1.691}{7.9701}{mi}{85}%
}\hskip.1em\allowbreak{} 存在特征值\allowbreak{}\hskip.1em \RB{27.880}{8.178}{2.858}{%
\RG{0.000}{-0.002}{10.9091}{mi}{22}\RG{9.604}{-0.002}{10.9091}{sy}{50}\RG{19.907}{-0.002}{10.9091}{bb}{67}%
}\hskip.1em\allowbreak{} 及特征向量\allowbreak{}\hskip.1em \RB{47.865}{8.309}{2.858}{%
\RG{0.000}{-0.002}{10.9091}{rm}{48}\RG{8.485}{-0.002}{10.9091}{sy}{54}\RG{8.485}{-0.002}{10.9091}{rm}{61}\RG{20.000}{-0.002}{10.9091}{mi}{121}\RG{28.771}{-0.002}{10.9091}{sy}{50}\RG{39.074}{-0.002}{10.9091}{mi}{85}%
}\hskip.1em\allowbreak{} 。于是

\[
\RB{136.444}{8.680}{3.229}{%
\RG{0.000}{-0.002}{10.9091}{rm}{40}\RG{4.242}{-0.002}{10.9091}{mi}{65}\RG{14.848}{-0.002}{10.9091}{sy}{0}\RG{25.757}{-0.002}{10.9091}{mi}{22}\RG{32.325}{-0.002}{10.9091}{mi}{73}\RG{42.762}{-0.002}{10.9091}{rm}{41}\RG{47.005}{-0.002}{10.9091}{mi}{121}\RG{55.775}{-0.002}{10.9091}{rm}{61}\RG{67.274}{-0.002}{10.9091}{rm}{48}\RG{72.745}{-0.002}{10.9091}{mi}{59}\RG{99.396}{-0.002}{10.9091}{mi}{67}\RG{107.970}{-0.002}{10.9091}{mi}{121}\RG{116.760}{-0.002}{10.9091}{rm}{61}\RG{128.258}{-0.002}{10.9091}{rm}{48}\RG{133.729}{-0.002}{10.9091}{mi}{59}\RG{37.126}{-1.638}{7.9701}{mi}{110}%
}
\]
即\allowbreak{}\hskip.1em \RB{66.218}{10.483}{4.953}{%
\RG{0.000}{8.835}{10.9092}{ex}{0}\RG{6.818}{4.292}{7.9701}{mi}{65}\RG{13.161}{4.292}{7.9701}{sy}{0}\RG{16.916}{-4.286}{7.9701}{mi}{67}\RG{19.748}{4.292}{7.9701}{mi}{22}\RG{24.841}{4.292}{7.9701}{mi}{73}\RG{28.514}{3.296}{5.9776}{mi}{110}\RG{35.431}{8.835}{10.9092}{ex}{1}\RG{40.431}{-0.001}{10.9091}{mi}{121}\RG{49.201}{-0.001}{10.9091}{rm}{61}\RG{60.700}{-0.001}{10.9091}{rm}{48}%
}\hskip.1em\allowbreak{} ，与题设的列满秩条件矛盾。因此ker\allowbreak{}\hskip.1em \RB{58.563}{8.681}{3.228}{%
\RG{0.000}{-0.001}{10.9091}{sy}{77}\RG{13.101}{-0.001}{10.9091}{rm}{40}\RG{17.343}{-0.001}{10.9091}{mi}{21}\RG{23.707}{-0.001}{10.9091}{rm}{41}\RG{30.962}{-0.001}{10.9091}{rm}{61}\RG{42.494}{-0.001}{10.9091}{sy}{102}\RG{47.949}{-0.001}{10.9091}{rm}{48}\RG{53.404}{-0.001}{10.9091}{sy}{103}%
}\hskip.1em\allowbreak{} ，也就是rank\allowbreak{}\hskip.1em \RB{49.225}{8.681}{3.228}{%
\RG{0.000}{-0.001}{10.9091}{sy}{77}\RG{13.101}{-0.001}{10.9091}{rm}{40}\RG{17.343}{-0.001}{10.9091}{mi}{21}\RG{23.707}{-0.001}{10.9091}{rm}{41}\RG{30.962}{-0.001}{10.9091}{rm}{61}\RG{42.495}{-0.001}{10.9091}{mi}{110}%
}\hskip.1em\allowbreak{} 。由\allowbreak{}\hskip.1em \RB{6.478}{8.069}{0.644}{%
\RG{0.000}{-0.002}{10.9091}{mi}{21}%
}\hskip.1em\allowbreak{} 的任意性，结论成立。\allowbreak{}\hskip.1em \RB{8.366}{7.993}{0.502}{%
\RG{0.000}{-0.002}{10.9091}{am}{3}%
}\hskip.1em\allowbreak{}

小结：“不列满秩”先产生非零核向量，再由它的迭代向量生成一个包含于ker\allowbreak{}\hskip.1em \RB{8.072}{7.519}{0.724}{%
\RG{0.000}{-0.005}{10.9091}{mi}{67}%
}\hskip.1em\allowbreak{} 的不变子空间，最后在该子空间内取特征向量。这里不必预先引用Cayley–Hamilton定理。

% Original page 5.
\section*{2015S05可对角化矩阵与循环矩阵}
题目设\allowbreak{}\hskip.1em \RB{8.366}{8.312}{0.500}{%
\RG{0.000}{0.001}{10.9091}{mi}{65}%
}\hskip.1em\allowbreak{} 是\allowbreak{}\hskip.1em \RB{6.729}{5.323}{0.619}{%
\RG{0.000}{0.001}{10.9091}{mi}{110}%
}\hskip.1em\allowbreak{} 阶复方阵，证明：\allowbreak{}\hskip.1em \RB{8.366}{8.312}{0.500}{%
\RG{0.000}{0.001}{10.9091}{mi}{65}%
}\hskip.1em\allowbreak{} 可对角化的充分必要条件是\allowbreak{}\hskip.1em \RB{8.366}{8.312}{0.500}{%
\RG{0.000}{0.001}{10.9091}{mi}{65}%
}\hskip.1em\allowbreak{} 相似于某个如下的循环矩阵：

\[
\RB{167.372}{42.265}{37.260}{%
\RG{23.122}{41.339}{10.9092}{ex}{48}\RG{153.413}{41.339}{10.9092}{ex}{49}\RG{38.530}{32.255}{10.9091}{mi}{97}\RG{65.307}{32.255}{10.9091}{mi}{97}\RG{86.220}{32.255}{10.9091}{mi}{97}\RG{106.682}{32.255}{10.9091}{sy}{1}\RG{111.515}{32.255}{10.9091}{sy}{1}\RG{116.358}{32.255}{10.9091}{sy}{1}\RG{136.600}{32.255}{10.9091}{mi}{97}\RG{44.297}{30.619}{7.9701}{rm}{49}\RG{71.073}{30.619}{7.9701}{rm}{50}\RG{91.987}{30.619}{7.9701}{rm}{51}\RG{142.366}{30.619}{7.9701}{mi}{110}\RG{23.122}{22.139}{10.9092}{ex}{66}\RG{153.413}{22.139}{10.9092}{ex}{67}\RG{38.078}{16.944}{10.9091}{mi}{97}\RG{65.307}{16.944}{10.9091}{mi}{97}\RG{86.220}{16.944}{10.9091}{mi}{97}\RG{106.682}{16.944}{10.9091}{sy}{1}\RG{111.515}{16.944}{10.9091}{sy}{1}\RG{116.358}{16.944}{10.9091}{sy}{1}\RG{131.190}{16.944}{10.9091}{mi}{97}\RG{23.122}{15.593}{10.9092}{ex}{66}\RG{153.413}{15.593}{10.9092}{ex}{67}\RG{43.845}{15.308}{7.9701}{mi}{110}\RG{71.073}{15.308}{7.9701}{rm}{49}\RG{91.987}{15.308}{7.9701}{rm}{50}\RG{136.956}{15.308}{7.9701}{mi}{110}\RG{142.094}{15.308}{7.9701}{sy}{0}\RG{148.681}{15.308}{7.9701}{rm}{49}\RG{23.122}{9.048}{10.9092}{ex}{66}\RG{153.413}{9.048}{10.9092}{ex}{67}\RG{23.122}{2.502}{10.9092}{ex}{66}\RG{153.413}{2.502}{10.9092}{ex}{67}\RG{32.668}{1.634}{10.9091}{mi}{97}\RG{64.855}{1.634}{10.9091}{mi}{97}\RG{86.220}{1.634}{10.9091}{mi}{97}\RG{106.682}{1.634}{10.9091}{sy}{1}\RG{111.515}{1.634}{10.9091}{sy}{1}\RG{116.358}{1.634}{10.9091}{sy}{1}\RG{131.190}{1.634}{10.9091}{mi}{97}\RG{38.434}{-0.003}{7.9701}{mi}{110}\RG{43.572}{-0.003}{7.9701}{sy}{0}\RG{50.159}{-0.003}{7.9701}{rm}{49}\RG{70.621}{-0.003}{7.9701}{mi}{110}\RG{91.987}{-0.003}{7.9701}{rm}{49}\RG{136.956}{-0.003}{7.9701}{mi}{110}\RG{142.094}{-0.003}{7.9701}{sy}{0}\RG{148.681}{-0.003}{7.9701}{rm}{50}\RG{0.000}{-0.225}{10.9091}{mi}{67}\RG{11.607}{-0.225}{10.9091}{rm}{61}\RG{164.777}{-0.225}{10.9091}{mi}{58}\RG{23.122}{-4.043}{10.9092}{ex}{66}\RG{153.413}{-4.043}{10.9092}{ex}{67}\RT{42.263}{-10.094}{10.9091}{.}\RT{69.040}{-10.094}{10.9091}{.}\RT{89.953}{-10.094}{10.9091}{.}\RT{140.785}{-10.094}{10.9091}{.}\RG{23.122}{-10.589}{10.9092}{ex}{66}\RG{153.413}{-10.589}{10.9092}{ex}{67}\RT{42.263}{-14.079}{10.9091}{.}\RT{69.040}{-14.079}{10.9091}{.}\RT{89.953}{-14.079}{10.9091}{.}\RT{140.785}{-14.079}{10.9091}{.}\RG{23.122}{-17.570}{10.9092}{ex}{64}\RG{153.413}{-17.571}{10.9092}{ex}{65}\RT{42.263}{-18.064}{10.9091}{.}\RT{69.040}{-18.064}{10.9091}{.}\RT{89.953}{-18.064}{10.9091}{.}\RT{140.785}{-18.064}{10.9091}{.}\RG{38.530}{-33.374}{10.9091}{mi}{97}\RG{65.307}{-33.374}{10.9091}{mi}{97}\RG{86.220}{-33.374}{10.9091}{mi}{97}\RG{106.682}{-33.374}{10.9091}{sy}{1}\RG{111.515}{-33.374}{10.9091}{sy}{1}\RG{116.358}{-33.374}{10.9091}{sy}{1}\RG{137.052}{-33.374}{10.9091}{mi}{97}\RG{44.297}{-35.011}{7.9701}{rm}{50}\RG{71.073}{-35.011}{7.9701}{rm}{51}\RG{91.987}{-35.011}{7.9701}{rm}{52}\RG{142.818}{-35.011}{7.9701}{rm}{49}%
}
\]
解：令\allowbreak{}\hskip.1em \RB{48.685}{10.441}{0.616}{%
\RG{0.000}{0.004}{10.9091}{mi}{33}\RG{10.212}{0.004}{10.9091}{rm}{61}\RG{21.728}{0.004}{10.9091}{mi}{101}\RG{26.807}{3.963}{7.9701}{rm}{50}\RG{31.041}{3.963}{7.9701}{mi}{25}\RG{36.182}{3.963}{7.9701}{mi}{105}\RG{39.073}{3.963}{7.9701}{rm}{47}\RG{43.307}{3.963}{7.9701}{mi}{110}%
}\hskip.1em\allowbreak{} 。记\allowbreak{}\hskip.1em \RB{7.536}{8.195}{0.736}{%
\RG{0.000}{0.004}{10.9091}{mi}{83}%
}\hskip.1em\allowbreak{} 为循环移位矩阵，即

\[
\RB{195.606}{11.433}{3.223}{%
\RT{78.700}{4.509}{7.9701}{T}\RT{186.728}{4.509}{7.9701}{T}\RG{0.000}{0.004}{10.9091}{mi}{83}\RG{7.318}{0.004}{10.9091}{rm}{40}\RG{11.560}{0.004}{10.9091}{mi}{120}\RG{22.527}{0.004}{10.9091}{mi}{59}\RG{27.360}{0.004}{10.9091}{mi}{120}\RG{38.343}{0.004}{10.9091}{mi}{59}\RG{43.176}{0.004}{10.9091}{mi}{58}\RG{48.020}{0.004}{10.9091}{mi}{58}\RG{52.853}{0.004}{10.9091}{mi}{58}\RG{57.696}{0.004}{10.9091}{mi}{59}\RG{62.540}{0.004}{10.9091}{mi}{120}\RG{74.457}{0.004}{10.9091}{rm}{41}\RG{88.015}{0.004}{10.9091}{rm}{61}\RG{99.513}{0.004}{10.9091}{rm}{40}\RG{103.772}{0.004}{10.9091}{mi}{120}\RG{114.739}{0.004}{10.9091}{mi}{59}\RG{119.572}{0.004}{10.9091}{mi}{120}\RG{130.555}{0.004}{10.9091}{mi}{59}\RG{135.388}{0.004}{10.9091}{mi}{58}\RG{140.232}{0.004}{10.9091}{mi}{58}\RG{145.064}{0.004}{10.9091}{mi}{58}\RG{149.908}{0.004}{10.9091}{mi}{59}\RG{154.752}{0.004}{10.9091}{mi}{120}\RG{166.668}{0.004}{10.9091}{mi}{59}\RG{171.501}{0.004}{10.9091}{mi}{120}\RG{182.484}{0.004}{10.9091}{rm}{41}\RG{193.012}{0.004}{10.9091}{mi}{58}\RG{17.795}{-1.632}{7.9701}{rm}{49}\RG{33.611}{-1.632}{7.9701}{rm}{50}\RG{68.821}{-1.632}{7.9701}{mi}{110}\RG{110.007}{-1.632}{7.9701}{rm}{50}\RG{125.823}{-1.632}{7.9701}{rm}{51}\RG{161.032}{-1.632}{7.9701}{mi}{110}\RG{177.752}{-1.632}{7.9701}{rm}{49}%
}
\]
题中的循环矩阵恰可写为

\[
\RB{153.548}{19.436}{13.622}{%
\RG{103.540}{13.636}{7.9701}{mi}{110}\RG{108.678}{13.636}{7.9701}{sy}{0}\RG{115.265}{13.636}{7.9701}{rm}{49}\RG{103.641}{10.363}{10.9092}{ex}{88}\RG{146.398}{4.504}{7.9701}{mi}{114}\RG{0.000}{0.000}{10.9091}{mi}{67}\RG{11.607}{0.000}{10.9091}{rm}{61}\RG{23.122}{0.000}{10.9091}{mi}{112}\RG{28.611}{0.000}{10.9091}{rm}{40}\RG{32.853}{0.000}{10.9091}{mi}{83}\RG{40.171}{0.000}{10.9091}{rm}{41}\RG{44.414}{0.000}{10.9091}{mi}{59}\RG{71.065}{0.000}{10.9091}{mi}{112}\RG{76.570}{0.000}{10.9091}{rm}{40}\RG{80.813}{0.000}{10.9091}{mi}{116}\RG{84.752}{0.000}{10.9091}{rm}{41}\RG{92.006}{0.000}{10.9091}{rm}{61}\RG{121.317}{0.000}{10.9091}{mi}{97}\RG{142.459}{0.000}{10.9091}{mi}{116}\RG{150.953}{0.000}{10.9091}{mi}{58}\RG{127.083}{-1.637}{7.9701}{mi}{114}\RG{131.140}{-1.637}{7.9701}{rm}{43}\RG{137.723}{-1.637}{7.9701}{rm}{49}\RG{104.081}{-12.955}{7.9701}{mi}{114}\RG{108.137}{-12.955}{7.9701}{rm}{61}\RG{114.720}{-12.955}{7.9701}{rm}{48}%
}
\]
第一步：每个循环矩阵均可对角化。取

\[
\RB{194.648}{11.430}{3.761}{%
\RG{45.942}{4.506}{7.9701}{mi}{106}\RG{81.750}{4.506}{7.9701}{rm}{40}\RG{85.043}{4.506}{7.9701}{mi}{110}\RG{90.181}{4.506}{7.9701}{sy}{0}\RG{96.767}{4.506}{7.9701}{rm}{49}\RG{100.999}{4.506}{7.9701}{rm}{41}\RG{104.295}{4.506}{7.9701}{mi}{106}\RT{112.920}{4.506}{7.9701}{T}\RG{0.000}{0.002}{10.9091}{mi}{118}\RG{12.700}{0.002}{10.9091}{rm}{61}\RG{24.198}{0.002}{10.9091}{rm}{40}\RG{28.431}{0.002}{10.9091}{rm}{49}\RG{33.912}{0.002}{10.9091}{mi}{59}\RG{38.745}{0.002}{10.9091}{mi}{33}\RG{50.324}{0.002}{10.9091}{mi}{59}\RG{55.157}{0.002}{10.9091}{mi}{58}\RG{60.001}{0.002}{10.9091}{mi}{58}\RG{64.845}{0.002}{10.9091}{mi}{58}\RG{69.677}{0.002}{10.9091}{mi}{59}\RG{74.521}{0.002}{10.9091}{mi}{33}\RG{108.677}{0.002}{10.9091}{rm}{41}\RG{119.204}{0.002}{10.9091}{mi}{59}\RG{145.871}{0.002}{10.9091}{rm}{48}\RG{154.355}{0.002}{10.9091}{sy}{20}\RG{165.871}{0.002}{10.9091}{mi}{106}\RG{173.999}{0.002}{10.9091}{mi}{60}\RG{185.508}{0.002}{10.9091}{mi}{110}\RG{192.053}{0.002}{10.9091}{mi}{58}\RG{5.287}{-1.635}{7.9701}{mi}{106}%
}
\]
由\allowbreak{}\hskip.1em \RB{36.302}{9.737}{0.618}{%
\RG{0.000}{0.002}{10.9091}{mi}{33}\RG{7.182}{3.961}{7.9701}{mi}{110}\RG{12.315}{3.961}{7.9701}{mi}{106}\RG{19.733}{0.002}{10.9091}{rm}{61}\RG{31.231}{0.002}{10.9091}{rm}{49}%
}\hskip.1em\allowbreak{} ，有\allowbreak{}\hskip.1em \RB{52.178}{9.737}{3.760}{%
\RG{0.000}{0.002}{10.9091}{mi}{83}\RG{7.309}{0.002}{10.9091}{mi}{118}\RG{12.606}{-1.634}{7.9701}{mi}{106}\RG{20.019}{0.002}{10.9091}{rm}{61}\RG{31.534}{0.002}{10.9091}{mi}{33}\RG{38.716}{3.961}{7.9701}{mi}{106}\RG{43.098}{0.002}{10.9091}{mi}{118}\RG{48.386}{-1.634}{7.9701}{mi}{106}%
}\hskip.1em\allowbreak{} ，故\allowbreak{}\hskip.1em \RB{67.410}{9.737}{3.760}{%
\RG{0.000}{0.002}{10.9091}{mi}{67}\RG{8.564}{0.002}{10.9091}{mi}{118}\RG{13.865}{-1.634}{7.9701}{mi}{106}\RG{21.277}{0.002}{10.9091}{rm}{61}\RG{32.792}{0.002}{10.9091}{mi}{112}\RG{38.281}{0.002}{10.9091}{rm}{40}\RG{42.524}{0.002}{10.9091}{mi}{33}\RG{49.705}{3.961}{7.9701}{mi}{106}\RG{54.088}{0.002}{10.9091}{rm}{41}\RG{58.330}{0.002}{10.9091}{mi}{118}\RG{63.619}{-1.634}{7.9701}{mi}{106}%
}\hskip.1em\allowbreak{} 。矩阵\allowbreak{}\hskip.1em \RB{86.979}{8.684}{3.225}{%
\RG{0.000}{0.002}{10.9091}{mi}{70}\RG{11.561}{0.002}{10.9091}{rm}{61}\RG{23.059}{0.002}{10.9091}{rm}{40}\RG{27.318}{0.002}{10.9091}{mi}{118}\RG{32.606}{-1.634}{7.9701}{rm}{48}\RG{37.338}{0.002}{10.9091}{mi}{59}\RG{42.171}{0.002}{10.9091}{mi}{58}\RG{47.015}{0.002}{10.9091}{mi}{58}\RG{51.858}{0.002}{10.9091}{mi}{58}\RG{56.691}{0.002}{10.9091}{mi}{59}\RG{61.535}{0.002}{10.9091}{mi}{118}\RG{66.870}{-1.634}{7.9701}{mi}{110}\RG{72.008}{-1.634}{7.9701}{sy}{0}\RG{78.594}{-1.634}{7.9701}{rm}{49}\RG{83.327}{0.002}{10.9091}{rm}{41}%
}\hskip.1em\allowbreak{} 的行列式为

\[
\RB{142.607}{10.866}{15.271}{%
\RG{51.027}{10.366}{10.9092}{ex}{89}\RG{87.371}{4.507}{7.9701}{mi}{106}\RG{112.268}{4.507}{7.9701}{mi}{105}\RT{0.000}{0.003}{10.9091}{d}\RT{6.066}{0.003}{10.9091}{e}\RT{10.909}{0.003}{10.9091}{t}\RG{16.971}{0.003}{10.9091}{mi}{70}\RG{28.532}{0.003}{10.9091}{rm}{61}\RG{75.947}{0.003}{10.9091}{rm}{40}\RG{80.189}{0.003}{10.9091}{mi}{33}\RG{94.178}{0.003}{10.9091}{sy}{0}\RG{105.087}{0.003}{10.9091}{mi}{33}\RG{115.650}{0.003}{10.9091}{rm}{41}\RG{122.922}{0.003}{10.9091}{sy}{54}\RG{122.922}{0.003}{10.9091}{rm}{61}\RG{134.420}{0.003}{10.9091}{rm}{48}\RG{139.892}{0.003}{10.9091}{mi}{59}\RG{40.047}{-13.145}{7.9701}{rm}{48}\RG{44.281}{-13.145}{7.9701}{sy}{20}\RG{50.868}{-13.145}{7.9701}{mi}{105}\RG{53.746}{-13.145}{7.9701}{mi}{60}\RG{60.329}{-13.145}{7.9701}{mi}{106}\RG{64.203}{-13.145}{7.9701}{mi}{60}\RG{70.786}{-13.145}{7.9701}{mi}{110}%
}
\]
所以

\[
\RB{191.392}{10.301}{4.310}{%
\RG{71.510}{8.834}{10.9092}{ex}{0}\RG{183.798}{8.834}{10.9092}{ex}{1}\RG{8.530}{4.501}{7.9701}{sy}{0}\RG{15.117}{4.501}{7.9701}{rm}{49}\RG{163.097}{4.501}{7.9701}{mi}{110}\RG{168.235}{4.501}{7.9701}{sy}{0}\RG{174.822}{4.501}{7.9701}{rm}{49}\RG{0.000}{-0.003}{10.9091}{mi}{70}\RG{19.849}{-0.003}{10.9091}{mi}{67}\RG{28.413}{-0.003}{10.9091}{mi}{70}\RG{39.987}{-0.003}{10.9091}{rm}{61}\RT{51.502}{-0.003}{10.9091}{d}\RT{57.568}{-0.003}{10.9091}{i}\RT{60.601}{-0.003}{10.9091}{a}\RT{66.055}{-0.003}{10.9091}{g}\RG{76.510}{-0.003}{10.9091}{mi}{112}\RG{81.999}{-0.003}{10.9091}{rm}{40}\RG{86.231}{-0.003}{10.9091}{rm}{49}\RG{91.686}{-0.003}{10.9091}{rm}{41}\RG{95.938}{-0.003}{10.9091}{mi}{59}\RG{100.771}{-0.003}{10.9091}{mi}{112}\RG{106.275}{-0.003}{10.9091}{rm}{40}\RG{110.518}{-0.003}{10.9091}{mi}{33}\RG{117.700}{-0.003}{10.9091}{rm}{41}\RG{121.942}{-0.003}{10.9091}{mi}{59}\RG{126.775}{-0.003}{10.9091}{mi}{58}\RG{131.618}{-0.003}{10.9091}{mi}{58}\RG{136.451}{-0.003}{10.9091}{mi}{58}\RG{141.295}{-0.003}{10.9091}{mi}{59}\RG{146.139}{-0.003}{10.9091}{mi}{112}\RG{151.673}{-0.003}{10.9091}{rm}{40}\RG{155.915}{-0.003}{10.9091}{mi}{33}\RG{179.554}{-0.003}{10.9091}{rm}{41}\RG{188.798}{-0.003}{10.9091}{mi}{58}%
}
\]
因此，若\allowbreak{}\hskip.1em \RB{8.366}{8.312}{0.500}{%
\RG{0.000}{0.001}{10.9091}{mi}{65}%
}\hskip.1em\allowbreak{} 相似于循环矩阵，则\allowbreak{}\hskip.1em \RB{8.366}{8.312}{0.500}{%
\RG{0.000}{0.001}{10.9091}{mi}{65}%
}\hskip.1em\allowbreak{} 可对角化。第二步：由任意对角矩阵构造循环矩阵。若\allowbreak{}\hskip.1em \RB{8.366}{8.311}{0.500}{%
\RG{0.000}{0.000}{10.9091}{mi}{65}%
}\hskip.1em\allowbreak{} 可对角化，取可逆阵\allowbreak{}\hskip.1em \RB{8.726}{7.951}{0.500}{%
\RG{0.000}{0.000}{10.9091}{mi}{80}%
}\hskip.1em\allowbreak{} ，使

\[
\RB{162.452}{10.305}{3.227}{%
\RG{8.519}{4.505}{7.9701}{sy}{0}\RG{15.105}{4.505}{7.9701}{rm}{49}\RG{0.000}{0.000}{10.9091}{mi}{80}\RG{19.837}{0.000}{10.9091}{mi}{65}\RG{28.019}{0.000}{10.9091}{mi}{80}\RG{39.568}{0.000}{10.9091}{rm}{61}\RG{51.067}{0.000}{10.9091}{rm}{3}\RG{61.670}{0.000}{10.9091}{rm}{61}\RT{73.206}{0.000}{10.9091}{d}\RT{79.271}{0.000}{10.9091}{i}\RT{82.304}{0.000}{10.9091}{a}\RT{87.758}{0.000}{10.9091}{g}\RG{93.213}{0.000}{10.9091}{rm}{40}\RG{97.455}{0.000}{10.9091}{mi}{21}\RG{108.551}{0.000}{10.9091}{mi}{59}\RG{113.384}{0.000}{10.9091}{mi}{58}\RG{118.227}{0.000}{10.9091}{mi}{58}\RG{123.071}{0.000}{10.9091}{mi}{58}\RG{127.904}{0.000}{10.9091}{mi}{59}\RG{132.747}{0.000}{10.9091}{mi}{21}\RG{155.614}{0.000}{10.9091}{rm}{41}\RG{159.857}{0.000}{10.9091}{mi}{58}\RG{103.819}{-1.636}{7.9701}{rm}{48}\RG{139.157}{-1.636}{7.9701}{mi}{110}\RG{144.295}{-1.636}{7.9701}{sy}{0}\RG{150.882}{-1.636}{7.9701}{rm}{49}%
}
\]
定义

\[
\RB{179.899}{19.436}{15.214}{%
\RG{46.444}{13.636}{7.9701}{mi}{110}\RG{51.583}{13.636}{7.9701}{sy}{0}\RG{58.169}{13.636}{7.9701}{rm}{49}\RG{46.545}{10.363}{10.9092}{ex}{88}\RG{37.429}{7.379}{10.9091}{rm}{49}\RG{82.150}{4.504}{7.9701}{sy}{0}\RG{88.737}{4.504}{7.9701}{mi}{106}\RG{92.610}{4.504}{7.9701}{mi}{114}\RG{0.000}{-0.001}{10.9091}{mi}{97}\RG{24.172}{-0.001}{10.9091}{rm}{61}\RG{64.222}{-0.001}{10.9091}{mi}{21}\RG{74.968}{-0.001}{10.9091}{mi}{33}\RG{122.381}{-0.001}{10.9091}{rm}{40}\RG{126.614}{-0.001}{10.9091}{rm}{48}\RG{135.108}{-0.001}{10.9091}{sy}{20}\RG{146.623}{-0.001}{10.9091}{mi}{114}\RG{154.871}{-0.001}{10.9091}{mi}{60}\RG{166.380}{-0.001}{10.9091}{mi}{110}\RG{172.941}{-0.001}{10.9091}{rm}{41}\RG{177.184}{-0.001}{10.9091}{mi}{59}\RG{5.766}{-1.637}{7.9701}{mi}{114}\RG{9.823}{-1.637}{7.9701}{rm}{43}\RG{16.406}{-1.637}{7.9701}{rm}{49}\RG{70.586}{-1.637}{7.9701}{mi}{106}\RG{36.883}{-7.484}{10.9091}{mi}{110}\RG{47.071}{-13.088}{7.9701}{mi}{106}\RG{50.955}{-13.088}{7.9701}{rm}{61}\RG{57.539}{-13.088}{7.9701}{rm}{48}\RR{36.883}{2.727}{6.548}{0.436}%
}
\]
并按题目所示形状构造\allowbreak{}\hskip.1em \RB{43.825}{8.684}{3.225}{%
\RG{0.000}{0.002}{10.9091}{mi}{67}\RG{11.608}{0.002}{10.9091}{rm}{61}\RG{23.123}{0.002}{10.9091}{mi}{112}\RG{28.611}{0.002}{10.9091}{rm}{40}\RG{32.854}{0.002}{10.9091}{mi}{83}\RG{40.172}{0.002}{10.9091}{rm}{41}%
}\hskip.1em\allowbreak{} 。对\allowbreak{}\hskip.1em \RB{47.282}{8.073}{1.993}{%
\RG{0.000}{0.002}{10.9091}{rm}{48}\RG{8.485}{0.002}{10.9091}{sy}{20}\RG{20.000}{0.002}{10.9091}{mi}{107}\RG{29.044}{0.002}{10.9091}{mi}{60}\RG{40.553}{0.002}{10.9091}{mi}{110}%
}\hskip.1em\allowbreak{} ，有

\[
\RB{161.735}{19.432}{15.217}{%
\RG{51.577}{13.632}{7.9701}{mi}{110}\RG{56.715}{13.632}{7.9701}{sy}{0}\RG{63.303}{13.632}{7.9701}{rm}{49}\RG{81.919}{13.632}{7.9701}{mi}{110}\RG{87.057}{13.632}{7.9701}{sy}{0}\RG{93.644}{13.632}{7.9701}{rm}{49}\RG{51.678}{10.360}{10.9092}{ex}{88}\RG{82.020}{10.360}{10.9092}{ex}{88}\RG{42.562}{7.376}{10.9091}{rm}{49}\RG{16.912}{4.500}{7.9701}{mi}{107}\RG{106.878}{4.500}{7.9701}{rm}{40}\RG{110.171}{4.500}{7.9701}{mi}{107}\RG{114.793}{4.500}{7.9701}{sy}{0}\RG{121.379}{4.500}{7.9701}{mi}{106}\RG{125.263}{4.500}{7.9701}{rm}{41}\RG{128.557}{4.500}{7.9701}{mi}{114}\RG{0.000}{-0.004}{10.9091}{mi}{112}\RG{5.488}{-0.004}{10.9091}{rm}{40}\RG{9.731}{-0.004}{10.9091}{mi}{33}\RG{22.032}{-0.004}{10.9091}{rm}{41}\RG{29.287}{-0.004}{10.9091}{rm}{61}\RG{69.355}{-0.004}{10.9091}{mi}{21}\RG{99.696}{-0.004}{10.9091}{mi}{33}\RG{136.142}{-0.004}{10.9091}{rm}{61}\RG{147.657}{-0.004}{10.9091}{mi}{21}\RG{159.140}{-0.004}{10.9091}{mi}{58}\RG{75.719}{-1.640}{7.9701}{mi}{106}\RG{154.020}{-1.781}{7.9701}{mi}{107}\RG{42.015}{-7.487}{10.9091}{mi}{110}\RG{82.460}{-12.958}{7.9701}{mi}{114}\RG{86.517}{-12.958}{7.9701}{rm}{61}\RG{93.100}{-12.958}{7.9701}{rm}{48}\RG{52.204}{-13.091}{7.9701}{mi}{106}\RG{56.088}{-13.091}{7.9701}{rm}{61}\RG{62.672}{-13.091}{7.9701}{rm}{48}\RR{42.015}{2.723}{6.548}{0.436}%
}
\]
最后一步是因为内层等比数列的和在\allowbreak{}\hskip.1em \RB{25.704}{8.069}{2.738}{%
\RG{0.000}{-0.002}{10.9091}{mi}{106}\RG{8.147}{-0.002}{10.9091}{rm}{61}\RG{19.662}{-0.002}{10.9091}{mi}{107}%
}\hskip.1em\allowbreak{} 时为\allowbreak{}\hskip.1em \RB{6.729}{5.320}{0.622}{%
\RG{0.000}{-0.002}{10.9091}{mi}{110}%
}\hskip.1em\allowbreak{} ，否则为\allowbreak{}\hskip.1em \RB{5.518}{7.764}{0.742}{%
\RG{0.000}{-0.002}{10.9091}{rm}{48}%
}\hskip.1em\allowbreak{} 。于是\allowbreak{}\hskip.1em \RB{110.369}{9.757}{0.742}{%
\RG{0.000}{-0.002}{10.9091}{mi}{70}\RG{8.530}{3.957}{7.9701}{sy}{0}\RG{15.116}{3.957}{7.9701}{rm}{49}\RG{19.849}{-0.002}{10.9091}{mi}{67}\RG{28.413}{-0.002}{10.9091}{mi}{70}\RG{39.987}{-0.002}{10.9091}{rm}{61}\RG{51.486}{-0.002}{10.9091}{rm}{3}\RG{62.089}{-0.002}{10.9091}{rm}{61}\RG{73.624}{-0.002}{10.9091}{mi}{80}\RG{82.143}{3.957}{7.9701}{sy}{0}\RG{88.729}{3.957}{7.9701}{rm}{49}\RG{93.462}{-0.002}{10.9091}{mi}{65}\RG{101.644}{-0.002}{10.9091}{mi}{80}%
}\hskip.1em\allowbreak{} ，即

\[
\RB{118.300}{10.303}{3.229}{%
\RG{21.292}{4.503}{7.9701}{sy}{0}\RG{27.878}{4.503}{7.9701}{rm}{49}\RG{36.853}{4.503}{7.9701}{sy}{0}\RG{43.440}{4.503}{7.9701}{rm}{49}\RG{77.646}{4.503}{7.9701}{sy}{0}\RG{84.233}{4.503}{7.9701}{rm}{49}\RG{0.000}{-0.002}{10.9091}{rm}{40}\RG{4.243}{-0.002}{10.9091}{mi}{80}\RG{12.752}{-0.002}{10.9091}{mi}{70}\RG{32.611}{-0.002}{10.9091}{rm}{41}\RG{48.172}{-0.002}{10.9091}{mi}{65}\RG{56.354}{-0.002}{10.9091}{rm}{40}\RG{60.597}{-0.002}{10.9091}{mi}{80}\RG{69.106}{-0.002}{10.9091}{mi}{70}\RG{88.965}{-0.002}{10.9091}{rm}{41}\RG{96.219}{-0.002}{10.9091}{rm}{61}\RG{107.753}{-0.002}{10.9091}{mi}{67}\RG{115.706}{-0.002}{10.9091}{mi}{58}%
}
\]
故\allowbreak{}\hskip.1em \RB{8.366}{8.310}{0.501}{%
\RG{0.000}{-0.001}{10.9091}{mi}{65}%
}\hskip.1em\allowbreak{} 相似于循环矩阵，充分必要性得证。\allowbreak{}\hskip.1em \RB{8.366}{7.994}{0.501}{%
\RG{0.000}{-0.001}{10.9091}{am}{3}%
}\hskip.1em\allowbreak{}

小结：单位根向量同时对角化全部循环矩阵；逆向则用上面的有限求和公式，把任意给定的特征值排列还原成循环矩阵的第一行。特征值不必互异。

% Original page 6.
\section*{2015S06循环子空间与Cayley–Hamilton定理}
题目设\allowbreak{}\hskip.1em \RB{8.889}{7.952}{0.739}{%
\RG{0.000}{0.001}{10.9091}{mi}{86}%
}\hskip.1em\allowbreak{} 是数域\allowbreak{}\hskip.1em \RB{10.198}{7.952}{0.500}{%
\RG{0.000}{0.001}{10.9091}{mi}{75}%
}\hskip.1em\allowbreak{} 上的\allowbreak{}\hskip.1em \RB{6.729}{5.323}{0.619}{%
\RG{0.000}{0.001}{10.9091}{mi}{110}%
}\hskip.1em\allowbreak{} 维线性空间，\allowbreak{}\hskip.1em \RB{7.253}{5.323}{2.877}{%
\RG{0.000}{0.001}{10.9091}{mi}{39}%
}\hskip.1em\allowbreak{} 是\allowbreak{}\hskip.1em \RB{8.889}{7.952}{0.739}{%
\RG{0.000}{0.001}{10.9091}{mi}{86}%
}\hskip.1em\allowbreak{} 上的线性变换。

(1)对任意\allowbreak{}\hskip.1em \RB{49.242}{8.308}{2.859}{%
\RG{0.000}{-0.003}{10.9091}{rm}{48}\RG{8.485}{-0.003}{10.9091}{sy}{54}\RG{8.485}{-0.003}{10.9091}{rm}{61}\RG{20.000}{-0.003}{10.9091}{mi}{11}\RG{30.049}{-0.003}{10.9091}{sy}{50}\RG{40.352}{-0.003}{10.9091}{mi}{86}%
}\hskip.1em\allowbreak{} ，证明\allowbreak{}\hskip.1em \RB{122.808}{9.756}{3.230}{%
\RG{0.000}{-0.003}{10.9091}{mi}{85}\RG{11.668}{-0.003}{10.9091}{rm}{61}\RG{23.183}{-0.003}{10.9091}{mi}{76}\RG{30.607}{-0.003}{10.9091}{rm}{40}\RG{34.850}{-0.003}{10.9091}{mi}{11}\RG{41.854}{-0.003}{10.9091}{mi}{59}\RG{46.698}{-0.003}{10.9091}{mi}{39}\RG{53.853}{-0.003}{10.9091}{rm}{40}\RG{58.096}{-0.003}{10.9091}{mi}{11}\RG{65.116}{-0.003}{10.9091}{rm}{41}\RG{69.358}{-0.003}{10.9091}{mi}{59}\RG{74.190}{-0.003}{10.9091}{mi}{39}\RG{81.343}{3.956}{7.9701}{rm}{50}\RG{86.075}{-0.003}{10.9091}{rm}{40}\RG{90.318}{-0.003}{10.9091}{mi}{11}\RG{97.337}{-0.003}{10.9091}{rm}{41}\RG{101.579}{-0.003}{10.9091}{mi}{59}\RG{106.412}{-0.003}{10.9091}{mi}{58}\RG{111.255}{-0.003}{10.9091}{mi}{58}\RG{116.088}{-0.003}{10.9091}{mi}{58}\RG{119.155}{-0.003}{10.9091}{rm}{41}%
}\hskip.1em\allowbreak{} 是包含\allowbreak{}\hskip.1em \RB{7.067}{5.319}{0.623}{%
\RG{0.000}{-0.003}{10.9091}{mi}{11}%
}\hskip.1em\allowbreak{} 的最小的\allowbreak{}\hskip.1em \RB{7.253}{5.319}{2.881}{%
\RG{0.000}{-0.003}{10.9091}{mi}{39}%
}\hskip.1em\allowbreak{} –不变子空间。称\allowbreak{}\hskip.1em \RB{8.791}{7.948}{0.743}{%
\RG{0.000}{-0.003}{10.9091}{mi}{85}%
}\hskip.1em\allowbreak{}

为\allowbreak{}\hskip.1em \RB{7.067}{5.318}{0.624}{%
\RG{0.000}{-0.004}{10.9091}{mi}{11}%
}\hskip.1em\allowbreak{} 关于\allowbreak{}\hskip.1em \RB{7.253}{5.318}{2.882}{%
\RG{0.000}{-0.004}{10.9091}{mi}{39}%
}\hskip.1em\allowbreak{} 的循环子空间，\allowbreak{}\hskip.1em \RB{7.067}{5.318}{0.624}{%
\RG{0.000}{-0.004}{10.9091}{mi}{11}%
}\hskip.1em\allowbreak{} 为其循环向量。若dim\allowbreak{}\hskip.1em \RB{32.934}{7.947}{0.744}{%
\RG{0.000}{-0.004}{10.9091}{mi}{85}\RG{11.668}{-0.004}{10.9091}{rm}{61}\RG{23.183}{-0.004}{10.9091}{mi}{109}%
}\hskip.1em\allowbreak{} ，证明\allowbreak{}\hskip.1em \RB{110.813}{9.755}{3.231}{%
\RG{0.000}{-0.004}{10.9091}{sy}{102}\RG{5.454}{-0.004}{10.9091}{mi}{11}\RG{12.458}{-0.004}{10.9091}{mi}{59}\RG{17.301}{-0.004}{10.9091}{mi}{39}\RG{24.458}{-0.004}{10.9091}{rm}{40}\RG{28.700}{-0.004}{10.9091}{mi}{11}\RG{35.719}{-0.004}{10.9091}{rm}{41}\RG{39.962}{-0.004}{10.9091}{mi}{59}\RG{44.795}{-0.004}{10.9091}{mi}{58}\RG{49.639}{-0.004}{10.9091}{mi}{58}\RG{54.472}{-0.004}{10.9091}{mi}{58}\RG{59.315}{-0.004}{10.9091}{mi}{59}\RG{64.159}{-0.004}{10.9091}{mi}{39}\RG{71.341}{3.955}{7.9701}{mi}{109}\RG{78.832}{3.955}{7.9701}{sy}{0}\RG{85.418}{3.955}{7.9701}{rm}{49}\RG{90.151}{-0.004}{10.9091}{rm}{40}\RG{94.393}{-0.004}{10.9091}{mi}{11}\RG{101.412}{-0.004}{10.9091}{rm}{41}\RG{105.655}{-0.004}{10.9091}{sy}{103}%
}\hskip.1em\allowbreak{}是\allowbreak{}\hskip.1em \RB{8.791}{7.947}{0.744}{%
\RG{0.000}{-0.004}{10.9091}{mi}{85}%
}\hskip.1em\allowbreak{} 的一组基。(2)利用(1)证明Cayley–Hamilton定理。(3)证明：\allowbreak{}\hskip.1em \RB{8.889}{7.951}{0.740}{%
\RG{0.000}{0.000}{10.9091}{mi}{86}%
}\hskip.1em\allowbreak{} 只有平凡的\allowbreak{}\hskip.1em \RB{7.253}{5.322}{2.878}{%
\RG{0.000}{0.000}{10.9091}{mi}{39}%
}\hskip.1em\allowbreak{} –不变子空间，当且仅当每个非零向量的循环子空间都等于\allowbreak{}\hskip.1em \RB{8.889}{7.951}{0.740}{%
\RG{0.000}{0.000}{10.9091}{mi}{86}%
}\hskip.1em\allowbreak{} 。(4)证明：\allowbreak{}\hskip.1em \RB{8.889}{7.953}{0.738}{%
\RG{0.000}{0.002}{10.9091}{mi}{86}%
}\hskip.1em\allowbreak{} 只有平凡的\allowbreak{}\hskip.1em \RB{7.253}{5.324}{2.876}{%
\RG{0.000}{0.002}{10.9091}{mi}{39}%
}\hskip.1em\allowbreak{} –不变子空间，当且仅当\allowbreak{}\hskip.1em \RB{7.253}{5.324}{2.876}{%
\RG{0.000}{0.002}{10.9091}{mi}{39}%
}\hskip.1em\allowbreak{} 的特征多项式\allowbreak{}\hskip.1em \RB{32.741}{8.684}{3.225}{%
\RG{0.000}{0.002}{10.9091}{mi}{102}\RG{6.515}{0.002}{10.9091}{rm}{40}\RG{10.758}{0.002}{10.9091}{mi}{21}\RG{17.121}{0.002}{10.9091}{rm}{41}\RG{24.376}{0.002}{10.9091}{rm}{61}%
}\hskip.1em\allowbreak{} det\allowbreak{}\hskip.1em \RB{46.837}{8.684}{3.225}{%
\RG{0.000}{0.002}{10.9091}{rm}{40}\RG{4.242}{0.002}{10.9091}{mi}{21}\RG{10.602}{0.002}{10.9091}{mi}{73}\RG{15.402}{-1.687}{7.9701}{mi}{86}\RG{25.138}{0.002}{10.9091}{sy}{0}\RG{36.047}{0.002}{10.9091}{mi}{39}\RG{43.184}{0.002}{10.9091}{rm}{41}%
}\hskip.1em\allowbreak{} 在\allowbreak{}\hskip.1em \RB{10.198}{7.953}{0.500}{%
\RG{0.000}{0.002}{10.9091}{mi}{75}%
}\hskip.1em\allowbreak{}

上不可约。

这里“平凡子空间”指\allowbreak{}\hskip.1em \RB{16.067}{8.679}{3.230}{%
\RG{0.000}{-0.003}{10.9091}{sy}{102}\RG{5.455}{-0.003}{10.9091}{rm}{48}\RG{10.909}{-0.003}{10.9091}{sy}{103}%
}\hskip.1em\allowbreak{} 和\allowbreak{}\hskip.1em \RB{8.889}{7.948}{0.743}{%
\RG{0.000}{-0.003}{10.9091}{mi}{86}%
}\hskip.1em\allowbreak{} ；\allowbreak{}\hskip.1em \RB{18.350}{8.679}{3.230}{%
\RG{0.000}{-0.003}{10.9091}{mi}{76}\RG{7.425}{-0.003}{10.9091}{rm}{40}\RG{11.667}{-0.003}{10.9091}{sy}{1}\RG{14.697}{-0.003}{10.9091}{rm}{41}%
}\hskip.1em\allowbreak{} 表示所列向量的线性包。

(1)最小不变性与循环基

\allowbreak{}\hskip.1em \RB{7.253}{5.325}{2.875}{%
\RG{0.000}{0.003}{10.9091}{mi}{39}%
}\hskip.1em\allowbreak{} 将每个生成元\allowbreak{}\hskip.1em \RB{26.432}{9.738}{3.224}{%
\RG{0.000}{0.003}{10.9091}{mi}{39}\RG{7.136}{3.962}{7.9701}{mi}{106}\RG{11.518}{0.003}{10.9091}{rm}{40}\RG{15.760}{0.003}{10.9091}{mi}{11}\RG{22.779}{0.003}{10.9091}{rm}{41}%
}\hskip.1em\allowbreak{} 送到下一个生成元，故\allowbreak{}\hskip.1em \RB{47.595}{8.685}{3.224}{%
\RG{0.000}{0.003}{10.9091}{mi}{39}\RG{7.136}{0.003}{10.9091}{rm}{40}\RG{11.378}{0.003}{10.9091}{mi}{85}\RG{20.016}{0.003}{10.9091}{rm}{41}\RG{27.289}{0.003}{10.9091}{sy}{18}\RG{38.804}{0.003}{10.9091}{mi}{85}%
}\hskip.1em\allowbreak{} 。若不变子空间\allowbreak{}\hskip.1em \RB{11.933}{7.954}{0.737}{%
\RG{0.000}{0.003}{10.9091}{mi}{87}%
}\hskip.1em\allowbreak{} 包含\allowbreak{}\hskip.1em \RB{7.067}{5.325}{0.617}{%
\RG{0.000}{0.003}{10.9091}{mi}{11}%
}\hskip.1em\allowbreak{} ，则由归纳知所有\allowbreak{}\hskip.1em \RB{52.288}{9.737}{3.225}{%
\RG{0.000}{0.002}{10.9091}{mi}{39}\RG{7.136}{3.961}{7.9701}{mi}{106}\RG{11.518}{0.002}{10.9091}{rm}{40}\RG{15.761}{0.002}{10.9091}{mi}{11}\RG{22.779}{0.002}{10.9091}{rm}{41}\RG{30.052}{0.002}{10.9091}{sy}{50}\RG{40.355}{0.002}{10.9091}{mi}{87}%
}\hskip.1em\allowbreak{} ，从而\allowbreak{}\hskip.1em \RB{35.117}{7.953}{1.993}{%
\RG{0.000}{0.002}{10.9091}{mi}{85}\RG{11.668}{0.002}{10.9091}{sy}{18}\RG{23.184}{0.002}{10.9091}{mi}{87}%
}\hskip.1em\allowbreak{} 。所以\allowbreak{}\hskip.1em \RB{8.791}{7.953}{0.738}{%
\RG{0.000}{0.002}{10.9091}{mi}{85}%
}\hskip.1em\allowbreak{} 的确是最小者。

取最小的正整数\allowbreak{}\hskip.1em \RB{5.256}{5.323}{0.619}{%
\RG{0.000}{0.001}{10.9091}{mi}{114}%
}\hskip.1em\allowbreak{} ，使

\[
\RB{189.747}{10.306}{4.306}{%
\RG{80.387}{8.838}{10.9092}{ex}{0}\RG{182.153}{8.838}{10.9092}{ex}{1}\RG{7.136}{4.506}{7.9701}{mi}{114}\RG{151.273}{4.506}{7.9701}{mi}{114}\RG{155.330}{4.506}{7.9701}{sy}{0}\RG{161.916}{4.506}{7.9701}{rm}{49}\RG{0.000}{0.002}{10.9091}{mi}{39}\RG{11.691}{0.002}{10.9091}{rm}{40}\RG{15.933}{0.002}{10.9091}{mi}{11}\RG{22.952}{0.002}{10.9091}{rm}{41}\RG{30.225}{0.002}{10.9091}{sy}{50}\RG{40.528}{0.002}{10.9091}{mi}{87}\RG{58.416}{0.002}{10.9091}{rm}{58}\RG{61.438}{0.002}{10.9091}{rm}{61}\RG{72.963}{0.002}{10.9091}{mi}{76}\RG{85.387}{0.002}{10.9091}{mi}{11}\RG{92.390}{0.002}{10.9091}{mi}{59}\RG{97.234}{0.002}{10.9091}{mi}{39}\RG{104.391}{0.002}{10.9091}{rm}{40}\RG{108.633}{0.002}{10.9091}{mi}{11}\RG{115.652}{0.002}{10.9091}{rm}{41}\RG{119.894}{0.002}{10.9091}{mi}{59}\RG{124.727}{0.002}{10.9091}{mi}{58}\RG{129.570}{0.002}{10.9091}{mi}{58}\RG{134.414}{0.002}{10.9091}{mi}{58}\RG{139.247}{0.002}{10.9091}{mi}{59}\RG{144.091}{0.002}{10.9091}{mi}{39}\RG{166.648}{0.002}{10.9091}{rm}{40}\RG{170.891}{0.002}{10.9091}{mi}{11}\RG{177.910}{0.002}{10.9091}{rm}{41}\RG{187.153}{0.002}{10.9091}{mi}{58}\RG{50.831}{-1.634}{7.9701}{mi}{114}%
}
\]
这样的\allowbreak{}\hskip.1em \RB{5.256}{5.324}{0.618}{%
\RG{0.000}{0.002}{10.9091}{mi}{114}%
}\hskip.1em\allowbreak{} 存在且\allowbreak{}\hskip.1em \RB{26.499}{7.440}{1.993}{%
\RG{0.000}{0.002}{10.9091}{mi}{114}\RG{8.255}{0.002}{10.9091}{sy}{20}\RG{19.770}{0.002}{10.9091}{mi}{110}%
}\hskip.1em\allowbreak{} ，因为\allowbreak{}\hskip.1em \RB{22.155}{7.768}{1.403}{%
\RG{0.000}{0.002}{10.9091}{mi}{110}\RG{7.582}{0.002}{10.9091}{rm}{43}\RG{17.084}{0.002}{10.9091}{rm}{49}%
}\hskip.1em\allowbreak{} 个迭代向量必线性相关。由最小性，\allowbreak{}\hskip.1em \RB{14.485}{7.953}{2.214}{%
\RG{0.000}{0.002}{10.9091}{mi}{87}\RG{10.303}{-1.634}{7.9701}{mi}{114}%
}\hskip.1em\allowbreak{} 的这\allowbreak{}\hskip.1em \RB{5.256}{5.324}{0.618}{%
\RG{0.000}{0.002}{10.9091}{mi}{114}%
}\hskip.1em\allowbreak{} 个生成元线性无关，且上式说明\allowbreak{}\hskip.1em \RB{59.510}{8.684}{3.225}{%
\RG{0.000}{0.002}{10.9091}{mi}{39}\RG{7.136}{0.002}{10.9091}{rm}{40}\RG{11.378}{0.002}{10.9091}{mi}{87}\RG{21.682}{-1.635}{7.9701}{mi}{114}\RG{26.236}{0.002}{10.9091}{rm}{41}\RG{33.509}{0.002}{10.9091}{sy}{18}\RG{45.024}{0.002}{10.9091}{mi}{87}\RG{55.327}{-1.635}{7.9701}{mi}{114}%
}\hskip.1em\allowbreak{} 。于是所有后续迭代向量仍在\allowbreak{}\hskip.1em \RB{14.485}{7.953}{2.215}{%
\RG{0.000}{0.002}{10.9091}{mi}{87}\RG{10.303}{-1.635}{7.9701}{mi}{114}%
}\hskip.1em\allowbreak{} 中，故\allowbreak{}\hskip.1em \RB{37.669}{7.953}{2.215}{%
\RG{0.000}{0.002}{10.9091}{mi}{85}\RG{11.668}{0.002}{10.9091}{rm}{61}\RG{23.183}{0.002}{10.9091}{mi}{87}\RG{33.486}{-1.635}{7.9701}{mi}{114}%
}\hskip.1em\allowbreak{} ，\allowbreak{}\hskip.1em \RB{20.974}{5.324}{0.618}{%
\RG{0.000}{0.002}{10.9091}{mi}{109}\RG{12.608}{0.002}{10.9091}{rm}{61}%
}\hskip.1em\allowbreak{} dim\allowbreak{}\hskip.1em \RB{28.440}{7.953}{0.738}{%
\RG{0.000}{0.002}{10.9091}{mi}{85}\RG{11.668}{0.002}{10.9091}{rm}{61}\RG{23.183}{0.002}{10.9091}{mi}{114}%
}\hskip.1em\allowbreak{} ，所求基即得。

(2)用循环子空间证明Cayley–Hamilton定理

设\allowbreak{}\hskip.1em \RB{32.613}{8.679}{3.230}{%
\RG{0.000}{-0.003}{10.9091}{mi}{102}\RG{6.516}{-0.003}{10.9091}{rm}{40}\RG{10.758}{-0.003}{10.9091}{mi}{120}\RG{16.993}{-0.003}{10.9091}{rm}{41}\RG{24.248}{-0.003}{10.9091}{rm}{61}%
}\hskip.1em\allowbreak{} det\allowbreak{}\hskip.1em \RB{46.707}{8.679}{3.230}{%
\RG{0.000}{-0.003}{10.9091}{rm}{40}\RG{4.243}{-0.003}{10.9091}{mi}{120}\RG{10.472}{-0.003}{10.9091}{mi}{73}\RG{15.273}{-1.692}{7.9701}{mi}{86}\RG{25.009}{-0.003}{10.9091}{sy}{0}\RG{35.918}{-0.003}{10.9091}{mi}{39}\RG{43.054}{-0.003}{10.9091}{rm}{41}%
}\hskip.1em\allowbreak{} 。要证\allowbreak{}\hskip.1em \RB{42.176}{8.679}{3.230}{%
\RG{0.000}{-0.003}{10.9091}{mi}{102}\RG{6.515}{-0.003}{10.9091}{rm}{40}\RG{10.757}{-0.003}{10.9091}{mi}{39}\RG{17.894}{-0.003}{10.9091}{rm}{41}\RG{25.149}{-0.003}{10.9091}{rm}{61}\RG{36.658}{-0.003}{10.9091}{rm}{48}%
}\hskip.1em\allowbreak{} ，只须对任意\allowbreak{}\hskip.1em \RB{49.242}{8.308}{2.859}{%
\RG{0.000}{-0.003}{10.9091}{rm}{48}\RG{8.484}{-0.003}{10.9091}{sy}{54}\RG{8.484}{-0.003}{10.9091}{rm}{61}\RG{19.999}{-0.003}{10.9091}{mi}{11}\RG{30.049}{-0.003}{10.9091}{sy}{50}\RG{40.352}{-0.003}{10.9091}{mi}{86}%
}\hskip.1em\allowbreak{} 证明\allowbreak{}\hskip.1em \RB{57.680}{8.679}{3.230}{%
\RG{0.000}{-0.003}{10.9091}{mi}{102}\RG{6.515}{-0.003}{10.9091}{rm}{40}\RG{10.758}{-0.003}{10.9091}{mi}{39}\RG{17.894}{-0.003}{10.9091}{rm}{41}\RG{22.127}{-0.003}{10.9091}{rm}{40}\RG{26.379}{-0.003}{10.9091}{mi}{11}\RG{33.398}{-0.003}{10.9091}{rm}{41}\RG{40.653}{-0.003}{10.9091}{rm}{61}\RG{52.162}{-0.003}{10.9091}{rm}{48}%
}\hskip.1em\allowbreak{} 。取其循环子空间\allowbreak{}\hskip.1em \RB{8.791}{7.947}{0.744}{%
\RG{0.000}{-0.004}{10.9091}{mi}{85}%
}\hskip.1em\allowbreak{} ，设dim\allowbreak{}\hskip.1em \RB{32.935}{7.947}{0.744}{%
\RG{0.000}{-0.004}{10.9091}{mi}{85}\RG{11.668}{-0.004}{10.9091}{rm}{61}\RG{23.184}{-0.004}{10.9091}{mi}{109}%
}\hskip.1em\allowbreak{} 。由(1)，存在唯一的\allowbreak{}\hskip.1em \RB{80.758}{7.947}{2.609}{%
\RG{0.000}{-0.004}{10.9091}{mi}{99}\RG{4.721}{-1.640}{7.9701}{rm}{48}\RG{9.453}{-0.004}{10.9091}{mi}{59}\RG{14.286}{-0.004}{10.9091}{mi}{58}\RG{19.130}{-0.004}{10.9091}{mi}{58}\RG{23.973}{-0.004}{10.9091}{mi}{58}\RG{28.806}{-0.004}{10.9091}{mi}{59}\RG{33.650}{-0.004}{10.9091}{mi}{99}\RG{38.416}{-1.640}{7.9701}{mi}{109}\RG{45.907}{-1.640}{7.9701}{sy}{0}\RG{52.493}{-1.640}{7.9701}{rm}{49}\RG{60.257}{-0.004}{10.9091}{sy}{50}\RG{70.560}{-0.004}{10.9091}{mi}{75}%
}\hskip.1em\allowbreak{} ，使

\[
\RB{187.404}{10.301}{3.230}{%
\RG{7.136}{4.501}{7.9701}{mi}{109}\RG{74.630}{4.501}{7.9701}{mi}{109}\RG{82.120}{4.501}{7.9701}{sy}{0}\RG{88.707}{4.501}{7.9701}{rm}{49}\RG{0.000}{-0.003}{10.9091}{mi}{39}\RG{15.125}{-0.003}{10.9091}{rm}{40}\RG{19.367}{-0.003}{10.9091}{mi}{11}\RG{26.386}{-0.003}{10.9091}{rm}{41}\RG{33.041}{-0.003}{10.9091}{rm}{43}\RG{43.962}{-0.003}{10.9091}{mi}{99}\RG{67.493}{-0.003}{10.9091}{mi}{39}\RG{93.439}{-0.003}{10.9091}{rm}{40}\RG{97.681}{-0.003}{10.9091}{mi}{11}\RG{104.700}{-0.003}{10.9091}{rm}{41}\RG{111.354}{-0.003}{10.9091}{rm}{43}\RG{122.276}{-0.003}{10.9091}{sy}{1}\RG{127.109}{-0.003}{10.9091}{sy}{1}\RG{131.952}{-0.003}{10.9091}{sy}{1}\RG{137.427}{-0.003}{10.9091}{rm}{43}\RG{148.337}{-0.003}{10.9091}{mi}{99}\RG{157.790}{-0.003}{10.9091}{mi}{11}\RG{167.839}{-0.003}{10.9091}{rm}{61}\RG{179.338}{-0.003}{10.9091}{rm}{48}\RG{184.810}{-0.003}{10.9091}{mi}{58}\RG{48.683}{-1.640}{7.9701}{mi}{109}\RG{56.173}{-1.640}{7.9701}{sy}{0}\RG{62.761}{-1.640}{7.9701}{rm}{49}\RG{153.057}{-1.640}{7.9701}{rm}{48}%
}
\]
令\allowbreak{}\hskip.1em \RB{160.196}{9.756}{3.230}{%
\RG{0.000}{-0.003}{10.9091}{mi}{104}\RG{6.286}{-0.003}{10.9091}{rm}{40}\RG{10.528}{-0.003}{10.9091}{mi}{120}\RG{16.763}{-0.003}{10.9091}{rm}{41}\RG{24.018}{-0.003}{10.9091}{rm}{61}\RG{35.551}{-0.003}{10.9091}{mi}{120}\RG{41.786}{3.956}{7.9701}{mi}{109}\RG{52.199}{-0.003}{10.9091}{rm}{43}\RG{63.109}{-0.003}{10.9091}{mi}{99}\RG{67.830}{-1.639}{7.9701}{mi}{109}\RG{75.320}{-1.639}{7.9701}{sy}{0}\RG{81.907}{-1.639}{7.9701}{rm}{49}\RG{86.639}{-0.003}{10.9091}{mi}{120}\RG{92.874}{3.956}{7.9701}{mi}{109}\RG{100.364}{3.956}{7.9701}{sy}{0}\RG{106.951}{3.956}{7.9701}{rm}{49}\RG{114.107}{-0.003}{10.9091}{rm}{43}\RG{125.016}{-0.003}{10.9091}{sy}{1}\RG{129.849}{-0.003}{10.9091}{sy}{1}\RG{134.692}{-0.003}{10.9091}{sy}{1}\RG{140.168}{-0.003}{10.9091}{rm}{43}\RG{151.077}{-0.003}{10.9091}{mi}{99}\RG{155.798}{-1.639}{7.9701}{rm}{48}%
}\hskip.1em\allowbreak{} ，则\allowbreak{}\hskip.1em \RB{57.450}{8.679}{3.230}{%
\RG{0.000}{-0.003}{10.9091}{mi}{104}\RG{6.286}{-0.003}{10.9091}{rm}{40}\RG{10.528}{-0.003}{10.9091}{mi}{39}\RG{17.664}{-0.003}{10.9091}{rm}{41}\RG{21.897}{-0.003}{10.9091}{rm}{40}\RG{26.149}{-0.003}{10.9091}{mi}{11}\RG{33.168}{-0.003}{10.9091}{rm}{41}\RG{40.423}{-0.003}{10.9091}{rm}{61}\RG{51.932}{-0.003}{10.9091}{rm}{48}%
}\hskip.1em\allowbreak{} 。

在循环基\allowbreak{}\hskip.1em \RB{99.611}{9.755}{3.230}{%
\RG{0.000}{-0.003}{10.9091}{mi}{11}\RG{7.004}{-0.003}{10.9091}{mi}{59}\RG{11.847}{-0.003}{10.9091}{mi}{39}\RG{19.004}{-0.003}{10.9091}{rm}{40}\RG{23.246}{-0.003}{10.9091}{mi}{11}\RG{30.265}{-0.003}{10.9091}{rm}{41}\RG{34.507}{-0.003}{10.9091}{mi}{59}\RG{39.340}{-0.003}{10.9091}{mi}{58}\RG{44.184}{-0.003}{10.9091}{mi}{58}\RG{49.027}{-0.003}{10.9091}{mi}{58}\RG{53.860}{-0.003}{10.9091}{mi}{59}\RG{58.704}{-0.003}{10.9091}{mi}{39}\RG{65.887}{3.955}{7.9701}{mi}{109}\RG{73.378}{3.955}{7.9701}{sy}{0}\RG{79.964}{3.955}{7.9701}{rm}{49}\RG{84.696}{-0.003}{10.9091}{rm}{40}\RG{88.939}{-0.003}{10.9091}{mi}{11}\RG{95.958}{-0.003}{10.9091}{rm}{41}%
}\hskip.1em\allowbreak{} 下，\allowbreak{}\hskip.1em \RB{16.971}{8.679}{3.230}{%
\RG{0.000}{-0.003}{10.9091}{mi}{39}\RG{7.137}{-0.003}{10.9091}{sy}{106}\RG{10.167}{-1.692}{7.9701}{mi}{85}%
}\hskip.1em\allowbreak{} 的表示矩阵为友阵

\[
\RB{159.991}{11.085}{74.985}{%
\RG{24.500}{10.160}{10.9092}{ex}{48}\RG{145.912}{10.160}{10.9092}{ex}{49}\RG{34.046}{-0.003}{10.9091}{rm}{48}\RG{53.999}{-0.003}{10.9091}{rm}{48}\RG{73.972}{-0.003}{10.9091}{sy}{1}\RG{78.804}{-0.003}{10.9091}{sy}{1}\RG{83.648}{-0.003}{10.9091}{sy}{1}\RG{98.480}{-0.003}{10.9091}{rm}{48}\RG{120.936}{-0.003}{10.9091}{sy}{0}\RG{129.421}{-0.003}{10.9091}{mi}{99}\RG{134.142}{-1.640}{7.9701}{rm}{48}\RG{24.500}{-9.040}{10.9092}{ex}{66}\RG{145.912}{-9.040}{10.9092}{ex}{67}\RG{34.046}{-15.314}{10.9091}{rm}{49}\RG{53.999}{-15.314}{10.9091}{rm}{48}\RG{73.972}{-15.314}{10.9091}{sy}{1}\RG{78.804}{-15.314}{10.9091}{sy}{1}\RG{83.648}{-15.314}{10.9091}{sy}{1}\RG{98.480}{-15.314}{10.9091}{rm}{48}\RG{120.936}{-15.314}{10.9091}{sy}{0}\RG{129.421}{-15.314}{10.9091}{mi}{99}\RG{24.500}{-15.585}{10.9092}{ex}{66}\RG{145.912}{-15.585}{10.9092}{ex}{67}\RG{134.142}{-16.950}{7.9701}{rm}{49}\RG{24.500}{-22.131}{10.9092}{ex}{66}\RG{145.912}{-22.131}{10.9092}{ex}{67}\RT{99.691}{-27.041}{10.9091}{.}\RT{128.388}{-27.041}{10.9091}{.}\RT{75.484}{-28.037}{10.9091}{.}\RG{24.500}{-28.676}{10.9092}{ex}{66}\RG{145.912}{-28.676}{10.9092}{ex}{67}\RT{79.728}{-31.026}{10.9091}{.}\RT{99.691}{-31.026}{10.9091}{.}\RT{128.388}{-31.026}{10.9091}{.}\RT{83.973}{-34.015}{10.9091}{.}\RG{0.000}{-34.676}{10.9091}{mi}{72}\RG{12.985}{-34.676}{10.9091}{rm}{61}\RG{157.276}{-34.676}{10.9091}{mi}{59}\RG{34.046}{-35.011}{10.9091}{rm}{48}\RG{53.999}{-35.011}{10.9091}{rm}{49}\RT{99.691}{-35.011}{10.9091}{.}\RT{128.388}{-35.011}{10.9091}{.}\RG{24.500}{-35.222}{10.9092}{ex}{66}\RG{145.912}{-35.222}{10.9092}{ex}{67}\RG{24.500}{-41.767}{10.9092}{ex}{66}\RG{145.912}{-41.767}{10.9092}{ex}{67}\RT{35.257}{-46.738}{10.9091}{.}\RT{50.975}{-47.735}{10.9091}{.}\RT{75.484}{-47.735}{10.9091}{.}\RG{24.500}{-48.313}{10.9092}{ex}{66}\RG{145.912}{-48.314}{10.9092}{ex}{67}\RT{35.257}{-50.724}{10.9091}{.}\RT{55.219}{-50.724}{10.9091}{.}\RT{79.728}{-50.724}{10.9091}{.}\RT{59.464}{-53.712}{10.9091}{.}\RT{83.973}{-53.712}{10.9091}{.}\RT{35.257}{-54.710}{10.9091}{.}\RG{98.480}{-54.710}{10.9091}{rm}{48}\RG{113.897}{-54.710}{10.9091}{sy}{0}\RG{122.382}{-54.710}{10.9091}{mi}{99}\RG{24.500}{-55.296}{10.9092}{ex}{64}\RG{145.912}{-55.296}{10.9092}{ex}{65}\RG{127.103}{-56.346}{7.9701}{mi}{109}\RG{134.594}{-56.346}{7.9701}{sy}{0}\RG{141.180}{-56.346}{7.9701}{rm}{50}\RG{34.046}{-70.020}{10.9091}{rm}{48}\RG{49.463}{-70.020}{10.9091}{sy}{1}\RG{54.296}{-70.020}{10.9091}{sy}{1}\RG{59.140}{-70.020}{10.9091}{sy}{1}\RG{78.517}{-70.020}{10.9091}{rm}{48}\RG{98.470}{-70.020}{10.9091}{rm}{49}\RG{113.897}{-70.020}{10.9091}{sy}{0}\RG{122.382}{-70.020}{10.9091}{mi}{99}\RG{127.103}{-71.657}{7.9701}{mi}{109}\RG{134.594}{-71.657}{7.9701}{sy}{0}\RG{141.180}{-71.657}{7.9701}{rm}{49}%
}
\]
其中\allowbreak{}\hskip.1em \RB{29.177}{7.767}{0.618}{%
\RG{0.000}{0.002}{10.9091}{mi}{109}\RG{12.608}{0.002}{10.9091}{rm}{61}\RG{24.106}{0.002}{10.9091}{rm}{49}%
}\hskip.1em\allowbreak{} 时理解为\allowbreak{}\hskip.1em \RB{50.334}{8.684}{3.225}{%
\RG{0.000}{0.002}{10.9091}{mi}{72}\RG{12.985}{0.002}{10.9091}{rm}{61}\RG{24.483}{0.002}{10.9091}{rm}{40}\RG{28.742}{0.002}{10.9091}{sy}{0}\RG{37.227}{0.002}{10.9091}{mi}{99}\RG{41.948}{-1.634}{7.9701}{rm}{48}\RG{46.681}{0.002}{10.9091}{rm}{41}%
}\hskip.1em\allowbreak{} 。直接将det\allowbreak{}\hskip.1em \RB{50.191}{8.684}{3.225}{%
\RG{0.000}{0.002}{10.9091}{rm}{40}\RG{4.243}{0.002}{10.9091}{mi}{120}\RG{10.472}{0.002}{10.9091}{mi}{73}\RG{15.273}{-1.634}{7.9701}{mi}{109}\RG{25.680}{0.002}{10.9091}{sy}{0}\RG{36.584}{0.002}{10.9091}{mi}{72}\RG{46.538}{0.002}{10.9091}{rm}{41}%
}\hskip.1em\allowbreak{} 按末列展开：该列第\allowbreak{}\hskip.1em \RB{3.696}{7.713}{0.618}{%
\RG{0.000}{0.002}{10.9091}{mi}{105}%
}\hskip.1em\allowbreak{} 行的代数余子式为\allowbreak{}\hskip.1em \RB{19.768}{9.761}{0.618}{%
\RG{0.000}{0.002}{10.9091}{mi}{120}\RG{6.235}{3.961}{7.9701}{mi}{105}\RG{9.118}{3.961}{7.9701}{sy}{0}\RG{15.705}{3.961}{7.9701}{rm}{49}%
}\hskip.1em\allowbreak{} ，所以

\[
\RB{246.446}{19.433}{13.759}{%
\RG{80.491}{13.633}{7.9701}{mi}{109}\RG{87.982}{13.633}{7.9701}{sy}{0}\RG{94.568}{13.633}{7.9701}{rm}{50}\RG{81.768}{10.360}{10.9092}{ex}{88}\RG{114.958}{4.501}{7.9701}{mi}{105}\RG{189.491}{4.501}{7.9701}{mi}{109}\RG{196.982}{4.501}{7.9701}{sy}{0}\RG{203.568}{4.501}{7.9701}{rm}{49}\RT{0.000}{-0.004}{10.9091}{d}\RT{6.066}{-0.004}{10.9091}{e}\RT{10.909}{-0.004}{10.9091}{t}\RG{15.153}{-0.004}{10.9091}{rm}{40}\RG{19.396}{-0.004}{10.9091}{mi}{120}\RG{25.625}{-0.004}{10.9091}{mi}{73}\RG{40.839}{-0.004}{10.9091}{sy}{0}\RG{51.748}{-0.004}{10.9091}{mi}{72}\RG{61.702}{-0.004}{10.9091}{rm}{41}\RG{68.957}{-0.004}{10.9091}{rm}{61}\RG{100.621}{-0.004}{10.9091}{mi}{99}\RG{108.723}{-0.004}{10.9091}{mi}{120}\RG{120.763}{-0.004}{10.9091}{rm}{43}\RG{131.661}{-0.004}{10.9091}{rm}{40}\RG{135.915}{-0.004}{10.9091}{mi}{120}\RG{144.574}{-0.004}{10.9091}{rm}{43}\RG{155.483}{-0.004}{10.9091}{mi}{99}\RG{179.013}{-0.004}{10.9091}{rm}{41}\RG{183.256}{-0.004}{10.9091}{mi}{120}\RG{211.331}{-0.004}{10.9091}{rm}{61}\RG{222.846}{-0.004}{10.9091}{mi}{104}\RG{229.131}{-0.004}{10.9091}{rm}{40}\RG{233.374}{-0.004}{10.9091}{mi}{120}\RG{239.609}{-0.004}{10.9091}{rm}{41}\RG{243.851}{-0.004}{10.9091}{mi}{58}\RG{30.426}{-1.640}{7.9701}{mi}{109}\RG{105.342}{-1.640}{7.9701}{mi}{105}\RG{160.204}{-1.640}{7.9701}{mi}{109}\RG{167.694}{-1.640}{7.9701}{sy}{0}\RG{174.281}{-1.640}{7.9701}{rm}{49}\RG{82.795}{-13.091}{7.9701}{mi}{105}\RG{85.678}{-13.091}{7.9701}{rm}{61}\RG{92.261}{-13.091}{7.9701}{rm}{48}%
}
\]
这一步只用行列式计算，并未使用待证定理。

% Original page 7.
将循环基扩张为\allowbreak{}\hskip.1em \RB{8.889}{7.955}{0.736}{%
\RG{0.000}{0.004}{10.9091}{mi}{86}%
}\hskip.1em\allowbreak{} 的一组基。因为\allowbreak{}\hskip.1em \RB{8.791}{7.955}{0.736}{%
\RG{0.000}{0.004}{10.9091}{mi}{85}%
}\hskip.1em\allowbreak{} 不变，\allowbreak{}\hskip.1em \RB{7.253}{5.326}{2.874}{%
\RG{0.000}{0.004}{10.9091}{mi}{39}%
}\hskip.1em\allowbreak{} 的表示矩阵为

\[
\RB{50.425}{0.937}{32.779}{%
\RG{0.000}{0.001}{10.9092}{ex}{32}\RG{37.375}{0.001}{10.9092}{ex}{33}\RG{8.636}{-11.333}{10.9091}{mi}{72}\RG{30.237}{-11.333}{10.9091}{sy}{3}\RG{47.830}{-18.654}{10.9091}{mi}{58}\RG{10.886}{-26.644}{10.9091}{rm}{48}\RG{28.553}{-26.644}{10.9091}{mi}{66}%
}
\]
故

\[
\RB{238.327}{9.289}{3.231}{%
\RG{0.000}{-0.004}{10.9091}{mi}{102}\RG{6.515}{-0.004}{10.9091}{rm}{40}\RG{10.758}{-0.004}{10.9091}{mi}{120}\RG{16.993}{-0.004}{10.9091}{rm}{41}\RG{24.248}{-0.004}{10.9091}{rm}{61}\RT{35.781}{-0.004}{10.9091}{d}\RT{41.847}{-0.004}{10.9091}{e}\RT{46.690}{-0.004}{10.9091}{t}\RG{50.933}{-0.004}{10.9091}{rm}{40}\RG{55.176}{-0.004}{10.9091}{mi}{120}\RG{61.406}{-0.004}{10.9091}{mi}{73}\RG{76.620}{-0.004}{10.9091}{sy}{0}\RG{87.529}{-0.004}{10.9091}{mi}{72}\RG{97.484}{-0.004}{10.9091}{rm}{41}\RT{103.544}{-0.004}{10.9091}{d}\RT{109.609}{-0.004}{10.9091}{e}\RT{114.453}{-0.004}{10.9091}{t}\RG{118.697}{-0.004}{10.9091}{rm}{40}\RG{122.939}{-0.004}{10.9091}{mi}{120}\RG{129.168}{-0.004}{10.9091}{mi}{73}\RG{156.107}{-0.004}{10.9091}{sy}{0}\RG{167.016}{-0.004}{10.9091}{mi}{66}\RG{175.838}{-0.004}{10.9091}{rm}{41}\RG{183.093}{-0.004}{10.9091}{rm}{61}\RG{194.627}{-0.004}{10.9091}{mi}{104}\RG{200.912}{-0.004}{10.9091}{rm}{40}\RG{205.155}{-0.004}{10.9091}{mi}{120}\RG{211.390}{-0.004}{10.9091}{rm}{41}\RG{215.632}{-0.004}{10.9091}{mi}{113}\RG{220.893}{-0.004}{10.9091}{rm}{40}\RG{225.136}{-0.004}{10.9091}{mi}{120}\RG{231.371}{-0.004}{10.9091}{rm}{41}\RG{235.613}{-0.004}{10.9091}{mi}{59}\RG{66.207}{-1.641}{7.9701}{mi}{109}\RG{133.970}{-1.641}{7.9701}{mi}{110}\RG{139.108}{-1.641}{7.9701}{sy}{0}\RG{145.694}{-1.641}{7.9701}{mi}{109}%
}
\]
其中\allowbreak{}\hskip.1em \RB{30.853}{5.321}{0.621}{%
\RG{0.000}{-0.001}{10.9091}{mi}{109}\RG{12.609}{-0.001}{10.9091}{rm}{61}\RG{24.124}{-0.001}{10.9091}{mi}{110}%
}\hskip.1em\allowbreak{} 时约定\allowbreak{}\hskip.1em \RB{39.574}{8.681}{3.228}{%
\RG{0.000}{-0.001}{10.9091}{mi}{113}\RG{5.261}{-0.001}{10.9091}{rm}{40}\RG{9.504}{-0.001}{10.9091}{mi}{120}\RG{15.739}{-0.001}{10.9091}{rm}{41}\RG{22.994}{-0.001}{10.9091}{rm}{61}\RG{34.503}{-0.001}{10.9091}{rm}{49}%
}\hskip.1em\allowbreak{} 。于是

\[
\RB{133.075}{8.682}{3.227}{%
\RG{0.000}{-0.000}{10.9091}{mi}{102}\RG{6.516}{-0.000}{10.9091}{rm}{40}\RG{10.758}{-0.000}{10.9091}{mi}{39}\RG{17.894}{-0.000}{10.9091}{rm}{41}\RG{22.127}{-0.000}{10.9091}{rm}{40}\RG{26.379}{-0.000}{10.9091}{mi}{11}\RG{33.398}{-0.000}{10.9091}{rm}{41}\RG{40.653}{-0.000}{10.9091}{rm}{61}\RG{52.186}{-0.000}{10.9091}{mi}{113}\RG{57.447}{-0.000}{10.9091}{rm}{40}\RG{61.691}{-0.000}{10.9091}{mi}{39}\RG{68.827}{-0.000}{10.9091}{rm}{41}\RG{73.070}{-0.000}{10.9091}{mi}{104}\RG{79.355}{-0.000}{10.9091}{rm}{40}\RG{83.597}{-0.000}{10.9091}{mi}{39}\RG{90.734}{-0.000}{10.9091}{rm}{41}\RG{94.967}{-0.000}{10.9091}{rm}{40}\RG{99.219}{-0.000}{10.9091}{mi}{11}\RG{106.238}{-0.000}{10.9091}{rm}{41}\RG{113.492}{-0.000}{10.9091}{rm}{61}\RG{125.002}{-0.000}{10.9091}{rm}{48}\RG{130.480}{-0.000}{10.9091}{mi}{58}%
}
\]
由\allowbreak{}\hskip.1em \RB{7.067}{5.322}{0.620}{%
\RG{0.000}{-0.000}{10.9091}{mi}{11}%
}\hskip.1em\allowbreak{} 的任意性，得到\allowbreak{}\hskip.1em \RB{42.176}{8.682}{3.227}{%
\RG{0.000}{-0.000}{10.9091}{mi}{102}\RG{6.515}{-0.000}{10.9091}{rm}{40}\RG{10.757}{-0.000}{10.9091}{mi}{39}\RG{17.894}{-0.000}{10.9091}{rm}{41}\RG{25.149}{-0.000}{10.9091}{rm}{61}\RG{36.658}{-0.000}{10.9091}{rm}{48}%
}\hskip.1em\allowbreak{} 。

(3)只有平凡不变子空间的循环刻画

若\allowbreak{}\hskip.1em \RB{8.889}{7.947}{0.744}{%
\RG{0.000}{-0.004}{10.9091}{mi}{86}%
}\hskip.1em\allowbreak{} 只有平凡不变子空间，则任意\allowbreak{}\hskip.1em \RB{28.499}{8.307}{2.860}{%
\RG{0.000}{-0.004}{10.9091}{mi}{11}\RG{10.762}{-0.004}{10.9091}{sy}{54}\RG{10.762}{-0.004}{10.9091}{rm}{61}\RG{22.980}{-0.004}{10.9091}{rm}{48}%
}\hskip.1em\allowbreak{} 的循环子空间是非零不变子空间，因而只能等于

\[
\RB{20.196}{9.222}{2.687}{%
\RG{0.000}{-0.005}{10.9091}{mi}{86}\RT{8.787}{-0.005}{10.9091}{。}%
}
\]
反之，设每个非零向量都生成整个\allowbreak{}\hskip.1em \RB{8.889}{7.956}{0.735}{%
\RG{0.000}{0.005}{10.9091}{mi}{86}%
}\hskip.1em\allowbreak{} 。任取非零不变子空间\allowbreak{}\hskip.1em \RB{11.933}{7.956}{0.735}{%
\RG{0.000}{0.005}{10.9091}{mi}{87}%
}\hskip.1em\allowbreak{} 与\allowbreak{}\hskip.1em \RB{52.285}{8.316}{2.851}{%
\RG{0.000}{0.005}{10.9091}{rm}{48}\RG{8.485}{0.005}{10.9091}{sy}{54}\RG{8.485}{0.005}{10.9091}{rm}{61}\RG{20.000}{0.005}{10.9091}{mi}{11}\RG{30.049}{0.005}{10.9091}{sy}{50}\RG{40.352}{0.005}{10.9091}{mi}{87}%
}\hskip.1em\allowbreak{} 。由(1)，\allowbreak{}\hskip.1em \RB{7.067}{5.327}{0.615}{%
\RG{0.000}{0.005}{10.9091}{mi}{11}%
}\hskip.1em\allowbreak{} 的循环子空间包含于\allowbreak{}\hskip.1em \RB{11.933}{7.955}{0.736}{%
\RG{0.000}{0.004}{10.9091}{mi}{87}%
}\hskip.1em\allowbreak{} ；而该循环子空间等于\allowbreak{}\hskip.1em \RB{8.889}{7.955}{0.736}{%
\RG{0.000}{0.004}{10.9091}{mi}{86}%
}\hskip.1em\allowbreak{} ，故\allowbreak{}\hskip.1em \RB{35.252}{7.955}{0.736}{%
\RG{0.000}{0.004}{10.9091}{mi}{87}\RG{14.848}{0.004}{10.9091}{rm}{61}\RG{26.363}{0.004}{10.9091}{mi}{86}%
}\hskip.1em\allowbreak{} 。所以只有平凡不变子空间。

(4)与特征多项式不可约的等价性

先设\allowbreak{}\hskip.1em \RB{20.645}{8.682}{3.227}{%
\RG{0.000}{-0.000}{10.9091}{mi}{102}\RG{6.515}{-0.000}{10.9091}{rm}{40}\RG{10.757}{-0.000}{10.9091}{mi}{120}\RG{16.992}{-0.000}{10.9091}{rm}{41}%
}\hskip.1em\allowbreak{} 不可约。若存在非平凡不变子空间\allowbreak{}\hskip.1em \RB{11.933}{7.951}{0.740}{%
\RG{0.000}{-0.000}{10.9091}{mi}{87}%
}\hskip.1em\allowbreak{} ，取其一组基并扩张为\allowbreak{}\hskip.1em \RB{8.889}{7.951}{0.740}{%
\RG{0.000}{-0.000}{10.9091}{mi}{86}%
}\hskip.1em\allowbreak{} 的基，得到表示矩

\noindent \RB{304.063}{31.017}{14.125}{%
\RT{0.000}{21.789}{10.9091}{阵}\RG{154.696}{18.655}{10.9092}{ex}{32}\RG{191.452}{18.655}{10.9092}{ex}{33}\RG{163.332}{7.321}{10.9091}{mi}{66}\RG{184.058}{7.321}{10.9091}{sy}{3}\RG{201.907}{-0.000}{10.9091}{mi}{59}\RG{228.574}{-0.000}{10.9091}{rm}{48}\RG{237.058}{-0.000}{10.9091}{mi}{60}\RT{248.574}{-0.000}{10.9091}{d}\RT{254.640}{-0.000}{10.9091}{i}\RT{257.672}{-0.000}{10.9091}{m}\RG{268.577}{-0.000}{10.9091}{mi}{87}\RG{283.414}{-0.000}{10.9091}{mi}{60}\RG{294.923}{-0.000}{10.9091}{mi}{110}\RG{301.468}{-0.000}{10.9091}{mi}{58}\RG{165.016}{-7.990}{10.9091}{rm}{48}\RG{182.118}{-7.990}{10.9091}{mi}{68}%
}

于是

\[
\RB{213.259}{9.289}{3.824}{%
\RG{0.000}{-0.004}{10.9091}{mi}{102}\RG{6.515}{-0.004}{10.9091}{rm}{40}\RG{10.758}{-0.004}{10.9091}{mi}{120}\RG{16.993}{-0.004}{10.9091}{rm}{41}\RG{24.248}{-0.004}{10.9091}{rm}{61}\RT{35.781}{-0.004}{10.9091}{d}\RT{41.847}{-0.004}{10.9091}{e}\RT{46.690}{-0.004}{10.9091}{t}\RG{50.933}{-0.004}{10.9091}{rm}{40}\RG{55.176}{-0.004}{10.9091}{mi}{120}\RG{61.406}{-0.004}{10.9091}{mi}{73}\RG{93.814}{-0.004}{10.9091}{sy}{0}\RG{104.723}{-0.004}{10.9091}{mi}{66}\RG{113.545}{-0.004}{10.9091}{rm}{41}\RT{119.605}{-0.004}{10.9091}{d}\RT{125.670}{-0.004}{10.9091}{e}\RT{130.514}{-0.004}{10.9091}{t}\RG{134.758}{-0.004}{10.9091}{rm}{40}\RG{139.000}{-0.004}{10.9091}{mi}{120}\RG{145.229}{-0.004}{10.9091}{mi}{73}\RG{189.363}{-0.004}{10.9091}{sy}{0}\RG{200.272}{-0.004}{10.9091}{mi}{68}\RG{209.607}{-0.004}{10.9091}{rm}{41}\RT{66.207}{-1.778}{7.9701}{d}\RT{70.909}{-1.778}{7.9701}{i}\RT{73.260}{-1.778}{7.9701}{m}\RG{81.726}{-1.778}{7.9701}{mi}{87}\RG{150.031}{-1.778}{7.9701}{mi}{110}\RG{155.169}{-1.778}{7.9701}{sy}{0}\RT{161.755}{-1.778}{7.9701}{d}\RT{166.458}{-1.778}{7.9701}{i}\RT{168.809}{-1.778}{7.9701}{m}\RG{177.274}{-1.778}{7.9701}{mi}{87}%
}
\]
分解为两个正次数多项式的乘积，矛盾。因此只有平凡不变子空间。

反之，设只有平凡不变子空间。由(3)及(1)，任取\allowbreak{}\hskip.1em \RB{27.349}{8.310}{2.857}{%
\RG{0.000}{-0.001}{10.9091}{mi}{11}\RG{10.191}{-0.001}{10.9091}{sy}{54}\RG{10.191}{-0.001}{10.9091}{rm}{61}\RG{21.831}{-0.001}{10.9091}{rm}{48}%
}\hskip.1em\allowbreak{} ，都有\allowbreak{}\hskip.1em \RB{97.259}{9.757}{3.228}{%
\RG{0.000}{-0.001}{10.9091}{mi}{11}\RG{7.004}{-0.001}{10.9091}{mi}{59}\RG{11.847}{-0.001}{10.9091}{mi}{39}\RG{19.004}{-0.001}{10.9091}{rm}{40}\RG{23.247}{-0.001}{10.9091}{mi}{11}\RG{30.265}{-0.001}{10.9091}{rm}{41}\RG{34.508}{-0.001}{10.9091}{mi}{59}\RG{39.341}{-0.001}{10.9091}{mi}{58}\RG{44.185}{-0.001}{10.9091}{mi}{58}\RG{49.018}{-0.001}{10.9091}{mi}{58}\RG{53.861}{-0.001}{10.9091}{mi}{59}\RG{58.705}{-0.001}{10.9091}{mi}{39}\RG{65.888}{3.957}{7.9701}{mi}{110}\RG{71.026}{3.957}{7.9701}{sy}{0}\RG{77.612}{3.957}{7.9701}{rm}{49}\RG{82.345}{-0.001}{10.9091}{rm}{40}\RG{86.587}{-0.001}{10.9091}{mi}{11}\RG{93.606}{-0.001}{10.9091}{rm}{41}%
}\hskip.1em\allowbreak{} 构成\allowbreak{}\hskip.1em \RB{8.889}{7.949}{0.742}{%
\RG{0.000}{-0.002}{10.9091}{mi}{86}%
}\hskip.1em\allowbreak{} 的一组基。所以对任一非零多项式\allowbreak{}\hskip.1em \RB{5.846}{5.320}{2.618}{%
\RG{0.000}{-0.002}{10.9091}{mi}{112}%
}\hskip.1em\allowbreak{} ，只要deg\allowbreak{}\hskip.1em \RB{26.748}{6.389}{2.618}{%
\RG{0.000}{-0.002}{10.9091}{mi}{112}\RG{8.509}{-0.002}{10.9091}{mi}{60}\RG{20.018}{-0.002}{10.9091}{mi}{110}%
}\hskip.1em\allowbreak{} ，就有\allowbreak{}\hskip.1em \RB{56.661}{8.680}{3.229}{%
\RG{0.000}{-0.002}{10.9091}{mi}{112}\RG{5.488}{-0.002}{10.9091}{rm}{40}\RG{9.731}{-0.002}{10.9091}{mi}{39}\RG{16.867}{-0.002}{10.9091}{rm}{41}\RG{21.100}{-0.002}{10.9091}{rm}{40}\RG{25.352}{-0.002}{10.9091}{mi}{11}\RG{32.371}{-0.002}{10.9091}{rm}{41}\RG{39.644}{-0.002}{10.9091}{sy}{54}\RG{39.644}{-0.002}{10.9091}{rm}{61}\RG{51.143}{-0.002}{10.9091}{rm}{48}%
}\hskip.1em\allowbreak{} 。

若\allowbreak{}\hskip.1em \RB{6.522}{8.188}{2.739}{%
\RG{0.000}{-0.003}{10.9091}{mi}{102}%
}\hskip.1em\allowbreak{} 可约，写成\allowbreak{}\hskip.1em \RB{33.091}{8.188}{2.739}{%
\RG{0.000}{-0.003}{10.9091}{mi}{102}\RG{9.545}{-0.003}{10.9091}{rm}{61}\RG{21.060}{-0.003}{10.9091}{mi}{103}\RG{26.635}{-0.003}{10.9091}{mi}{104}%
}\hskip.1em\allowbreak{} ，其中\allowbreak{}\hskip.1em \RB{16.556}{7.762}{1.998}{%
\RG{0.000}{-0.003}{10.9091}{rm}{49}\RG{8.485}{-0.003}{10.9091}{sy}{20}%
}\hskip.1em\allowbreak{} deg\allowbreak{}\hskip.1em \RB{8.289}{5.319}{2.739}{%
\RG{0.000}{-0.003}{10.9091}{mi}{103}\RG{5.575}{-0.003}{10.9091}{mi}{59}%
}\hskip.1em\allowbreak{} deg\allowbreak{}\hskip.1em \RB{27.544}{8.068}{0.939}{%
\RG{0.000}{-0.003}{10.9091}{mi}{104}\RG{9.306}{-0.003}{10.9091}{mi}{60}\RG{20.815}{-0.003}{10.9091}{mi}{110}%
}\hskip.1em\allowbreak{} 。令

\[
\RB{152.242}{9.291}{3.229}{%
\RG{0.000}{-0.002}{10.9091}{mi}{12}\RG{9.776}{-0.002}{10.9091}{rm}{61}\RG{21.291}{-0.002}{10.9091}{mi}{104}\RG{27.576}{-0.002}{10.9091}{rm}{40}\RG{31.819}{-0.002}{10.9091}{mi}{39}\RG{38.955}{-0.002}{10.9091}{rm}{41}\RG{43.188}{-0.002}{10.9091}{rm}{40}\RG{47.440}{-0.002}{10.9091}{mi}{11}\RG{54.459}{-0.002}{10.9091}{rm}{41}\RG{58.702}{-0.002}{10.9091}{mi}{59}\RG{85.354}{-0.002}{10.9091}{mi}{87}\RG{100.218}{-0.002}{10.9091}{rm}{61}\RT{111.733}{-0.002}{10.9091}{k}\RT{117.493}{-0.002}{10.9091}{e}\RT{122.337}{-0.002}{10.9091}{r}\RG{128.431}{-0.002}{10.9091}{mi}{103}\RG{134.026}{-0.002}{10.9091}{rm}{40}\RG{138.268}{-0.002}{10.9091}{mi}{39}\RG{145.404}{-0.002}{10.9091}{rm}{41}\RG{149.647}{-0.002}{10.9091}{mi}{58}%
}
\]
由于deg\allowbreak{}\hskip.1em \RB{27.544}{8.069}{0.938}{%
\RG{0.000}{-0.002}{10.9091}{mi}{104}\RG{9.306}{-0.002}{10.9091}{mi}{60}\RG{20.815}{-0.002}{10.9091}{mi}{110}%
}\hskip.1em\allowbreak{} ，有\allowbreak{}\hskip.1em \RB{26.793}{8.309}{2.858}{%
\RG{0.000}{-0.002}{10.9091}{mi}{12}\RG{9.776}{-0.002}{10.9091}{sy}{54}\RG{9.776}{-0.002}{10.9091}{rm}{61}\RG{21.274}{-0.002}{10.9091}{rm}{48}%
}\hskip.1em\allowbreak{} ；由(2)，

\[
\RB{111.348}{8.681}{3.228}{%
\RG{0.000}{-0.001}{10.9091}{mi}{103}\RG{5.595}{-0.001}{10.9091}{rm}{40}\RG{9.838}{-0.001}{10.9091}{mi}{39}\RG{16.974}{-0.001}{10.9091}{rm}{41}\RG{21.207}{-0.001}{10.9091}{rm}{40}\RG{25.459}{-0.001}{10.9091}{mi}{12}\RG{32.205}{-0.001}{10.9091}{rm}{41}\RG{39.460}{-0.001}{10.9091}{rm}{61}\RG{50.993}{-0.001}{10.9091}{mi}{102}\RG{57.508}{-0.001}{10.9091}{rm}{40}\RG{61.750}{-0.001}{10.9091}{mi}{39}\RG{68.888}{-0.001}{10.9091}{rm}{41}\RG{73.120}{-0.001}{10.9091}{rm}{40}\RG{77.373}{-0.001}{10.9091}{mi}{11}\RG{84.392}{-0.001}{10.9091}{rm}{41}\RG{91.646}{-0.001}{10.9091}{rm}{61}\RG{103.155}{-0.001}{10.9091}{rm}{48}\RG{108.634}{-0.001}{10.9091}{mi}{59}%
}
\]
故\allowbreak{}\hskip.1em \RB{42.432}{8.681}{3.228}{%
\RG{0.000}{-0.001}{10.9091}{mi}{87}\RG{14.849}{-0.001}{10.9091}{sy}{54}\RG{14.849}{-0.001}{10.9091}{rm}{61}\RG{26.364}{-0.001}{10.9091}{sy}{102}\RG{31.818}{-0.001}{10.9091}{rm}{48}\RG{37.273}{-0.001}{10.9091}{sy}{103}%
}\hskip.1em\allowbreak{} 。又因\allowbreak{}\hskip.1em \RB{56.767}{8.681}{3.228}{%
\RG{0.000}{-0.001}{10.9091}{mi}{103}\RG{5.594}{-0.001}{10.9091}{rm}{40}\RG{9.837}{-0.001}{10.9091}{mi}{39}\RG{16.973}{-0.001}{10.9091}{rm}{41}\RG{21.206}{-0.001}{10.9091}{rm}{40}\RG{25.458}{-0.001}{10.9091}{mi}{11}\RG{32.477}{-0.001}{10.9091}{rm}{41}\RG{39.750}{-0.001}{10.9091}{sy}{54}\RG{39.750}{-0.001}{10.9091}{rm}{61}\RG{51.249}{-0.001}{10.9091}{rm}{48}%
}\hskip.1em\allowbreak{} ，有\allowbreak{}\hskip.1em \RB{35.253}{8.310}{2.857}{%
\RG{0.000}{-0.001}{10.9091}{mi}{87}\RG{14.849}{-0.001}{10.9091}{sy}{54}\RG{14.849}{-0.001}{10.9091}{rm}{61}\RG{26.364}{-0.001}{10.9091}{mi}{86}%
}\hskip.1em\allowbreak{} 。最后，\allowbreak{}\hskip.1em \RB{20.626}{8.681}{3.228}{%
\RG{0.000}{-0.001}{10.9091}{mi}{103}\RG{5.595}{-0.001}{10.9091}{rm}{40}\RG{9.837}{-0.001}{10.9091}{mi}{39}\RG{16.973}{-0.001}{10.9091}{rm}{41}%
}\hskip.1em\allowbreak{} 与\allowbreak{}\hskip.1em \RB{7.253}{5.321}{2.879}{%
\RG{0.000}{-0.001}{10.9091}{mi}{39}%
}\hskip.1em\allowbreak{} 可交换，所以对\allowbreak{}\hskip.1em \RB{33.370}{7.950}{0.937}{%
\RG{0.000}{-0.001}{10.9091}{mi}{119}\RG{11.134}{-0.001}{10.9091}{sy}{50}\RG{21.437}{-0.001}{10.9091}{mi}{87}%
}\hskip.1em\allowbreak{} ，

\[
\RB{144.113}{8.681}{3.228}{%
\RG{0.000}{-0.001}{10.9091}{mi}{103}\RG{5.595}{-0.001}{10.9091}{rm}{40}\RG{9.838}{-0.001}{10.9091}{mi}{39}\RG{16.974}{-0.001}{10.9091}{rm}{41}\RG{21.207}{-0.001}{10.9091}{rm}{40}\RG{25.459}{-0.001}{10.9091}{mi}{39}\RG{32.595}{-0.001}{10.9091}{rm}{40}\RG{36.838}{-0.001}{10.9091}{mi}{119}\RG{44.941}{-0.001}{10.9091}{rm}{41}\RG{49.174}{-0.001}{10.9091}{rm}{41}\RG{56.429}{-0.001}{10.9091}{rm}{61}\RG{67.972}{-0.001}{10.9091}{mi}{39}\RG{75.109}{-0.001}{10.9091}{rm}{40}\RG{79.351}{-0.001}{10.9091}{mi}{103}\RG{84.946}{-0.001}{10.9091}{rm}{40}\RG{89.188}{-0.001}{10.9091}{mi}{39}\RG{96.325}{-0.001}{10.9091}{rm}{41}\RG{100.558}{-0.001}{10.9091}{rm}{40}\RG{104.810}{-0.001}{10.9091}{mi}{119}\RG{112.913}{-0.001}{10.9091}{rm}{41}\RG{117.146}{-0.001}{10.9091}{rm}{41}\RG{124.400}{-0.001}{10.9091}{rm}{61}\RG{135.909}{-0.001}{10.9091}{rm}{48}\RG{141.398}{-0.001}{10.9091}{mi}{59}%
}
\]
即\allowbreak{}\hskip.1em \RB{11.933}{7.951}{0.740}{%
\RG{0.000}{-0.000}{10.9091}{mi}{87}%
}\hskip.1em\allowbreak{} 是非平凡不变子空间，矛盾。因此\allowbreak{}\hskip.1em \RB{6.522}{8.191}{2.736}{%
\RG{0.000}{-0.000}{10.9091}{mi}{102}%
}\hskip.1em\allowbreak{} 不可约。\allowbreak{}\hskip.1em \RB{8.366}{7.995}{0.500}{%
\RG{0.000}{-0.000}{10.9091}{am}{3}%
}\hskip.1em\allowbreak{}

小结：第一次出现的迭代关系给出循环基与友阵；不变子空间使特征多项式分解；最后将局部关系\allowbreak{}\hskip.1em \RB{44.767}{7.975}{2.988}{%
\RG{0.000}{0.003}{10.9091}{mi}{104}\RG{5.740}{0.003}{10.9091}{rm}{40}\RG{9.614}{0.003}{10.9091}{mi}{39}\RG{16.131}{0.003}{10.9091}{rm}{41}\RG{20.006}{0.003}{10.9091}{mi}{11}\RG{29.183}{0.003}{10.9091}{rm}{61}\RG{39.684}{0.003}{10.9091}{rm}{48}%
}\hskip.1em\allowbreak{} 提升为整体关系\allowbreak{}\hskip.1em \RB{38.560}{7.975}{2.988}{%
\RG{0.000}{0.003}{10.9091}{mi}{102}\RG{5.950}{0.003}{10.9091}{rm}{40}\RG{9.824}{0.003}{10.9091}{mi}{39}\RG{16.341}{0.003}{10.9091}{rm}{41}\RG{22.966}{0.003}{10.9091}{rm}{61}\RG{33.477}{0.003}{10.9091}{rm}{48}%
}\hskip.1em\allowbreak{} 。第(4)问的不可约性始终是在原数域\allowbreak{}\hskip.1em \RB{9.357}{7.308}{0.500}{%
\RG{0.000}{0.003}{10.9091}{mi}{75}%
}\hskip.1em\allowbreak{} 上讨论。

% Original page 8.
\section*{2015S07通过非异对称矩阵相似于转置}
题目设\allowbreak{}\hskip.1em \RB{8.366}{8.312}{0.500}{%
\RG{0.000}{0.001}{10.9091}{mi}{65}%
}\hskip.1em\allowbreak{} 为\allowbreak{}\hskip.1em \RB{6.729}{5.323}{0.619}{%
\RG{0.000}{0.001}{10.9091}{mi}{110}%
}\hskip.1em\allowbreak{} 阶复方阵，证明：存在\allowbreak{}\hskip.1em \RB{6.729}{5.323}{0.619}{%
\RG{0.000}{0.001}{10.9091}{mi}{110}%
}\hskip.1em\allowbreak{} 阶非异复对称阵\allowbreak{}\hskip.1em \RB{8.573}{8.192}{2.615}{%
\RG{0.000}{0.001}{10.9091}{mi}{81}%
}\hskip.1em\allowbreak{} ，使

\[
\RB{64.980}{10.306}{2.615}{%
\RG{8.182}{4.506}{7.9701}{sy}{48}\RG{34.147}{4.506}{7.9701}{sy}{0}\RG{40.733}{4.506}{7.9701}{rm}{49}\RG{0.000}{0.001}{10.9091}{mi}{65}\RG{14.007}{0.001}{10.9091}{rm}{61}\RG{25.522}{0.001}{10.9091}{mi}{81}\RG{45.465}{0.001}{10.9091}{mi}{65}\RG{53.647}{0.001}{10.9091}{mi}{81}\RG{62.266}{0.001}{10.9091}{mi}{59}%
}
\]
即\allowbreak{}\hskip.1em \RB{8.366}{8.313}{0.500}{%
\RG{0.000}{0.002}{10.9091}{mi}{65}%
}\hskip.1em\allowbreak{} 可通过非异复对称阵相似于其转置\allowbreak{}\hskip.1em \RB{10.866}{8.916}{0.500}{%
\RG{0.000}{0.002}{10.9091}{mi}{65}\RG{8.182}{3.961}{7.9701}{sy}{48}%
}\hskip.1em\allowbreak{} 。解：下文用\allowbreak{}\hskip.1em \RB{8.366}{8.312}{0.500}{%
\RG{0.000}{0.001}{10.9091}{mi}{65}%
}\hskip.1em\allowbreak{} T表示普通转置，即原题中的\allowbreak{}\hskip.1em \RB{10.866}{8.915}{0.500}{%
\RG{0.000}{0.001}{10.9091}{mi}{65}\RG{8.182}{3.960}{7.9701}{sy}{48}%
}\hskip.1em\allowbreak{} 。取Jordan分解

\[
\RB{238.522}{10.306}{4.306}{%
\RG{126.501}{8.838}{10.9092}{ex}{0}\RG{230.807}{8.838}{10.9092}{ex}{1}\RG{46.864}{4.506}{7.9701}{sy}{0}\RG{53.450}{4.506}{7.9701}{rm}{49}\RG{0.000}{0.002}{10.9091}{mi}{65}\RG{11.212}{0.002}{10.9091}{rm}{61}\RG{22.728}{0.002}{10.9091}{mi}{80}\RG{31.237}{0.002}{10.9091}{mi}{74}\RG{38.339}{0.002}{10.9091}{mi}{80}\RG{58.183}{0.002}{10.9091}{mi}{59}\RG{84.835}{0.002}{10.9091}{mi}{74}\RG{94.979}{0.002}{10.9091}{rm}{61}\RT{106.494}{0.002}{10.9091}{d}\RT{112.559}{0.002}{10.9091}{i}\RT{115.592}{0.002}{10.9091}{a}\RT{121.047}{0.002}{10.9091}{g}\RG{131.502}{0.002}{10.9091}{mi}{74}\RG{149.690}{0.002}{10.9091}{rm}{40}\RG{153.933}{0.002}{10.9091}{mi}{21}\RG{165.029}{0.002}{10.9091}{rm}{41}\RG{169.271}{0.002}{10.9091}{mi}{59}\RG{174.104}{0.002}{10.9091}{mi}{58}\RG{178.948}{0.002}{10.9091}{mi}{58}\RG{183.792}{0.002}{10.9091}{mi}{58}\RG{188.624}{0.002}{10.9091}{mi}{59}\RG{193.468}{0.002}{10.9091}{mi}{74}\RG{211.544}{0.002}{10.9091}{rm}{40}\RG{215.787}{0.002}{10.9091}{mi}{21}\RG{226.565}{0.002}{10.9091}{rm}{41}\RG{235.807}{0.002}{10.9091}{mi}{59}\RG{137.551}{-1.635}{7.9701}{mi}{109}\RG{160.297}{-1.635}{7.9701}{rm}{49}\RG{199.564}{-1.635}{7.9701}{mi}{109}\RG{222.151}{-1.635}{7.9701}{mi}{115}\RG{207.054}{-2.631}{5.9776}{mi}{115}\RG{145.041}{-2.742}{5.9776}{rm}{49}%
}
\]
其中\allowbreak{}\hskip.1em \RB{8.726}{7.953}{0.500}{%
\RG{0.000}{0.002}{10.9091}{mi}{80}%
}\hskip.1em\allowbreak{} 可逆，每个Jordan块的超对角线元素为\allowbreak{}\hskip.1em \RB{5.071}{7.768}{0.500}{%
\RG{0.000}{0.002}{10.9091}{rm}{49}%
}\hskip.1em\allowbreak{} 。

对每个正整数\allowbreak{}\hskip.1em \RB{9.751}{5.323}{0.619}{%
\RG{0.000}{0.001}{10.9091}{mi}{109}%
}\hskip.1em\allowbreak{} ，取反单位矩阵

\[
\RB{236.331}{24.490}{15.761}{%
\RG{139.404}{24.001}{10.9092}{ex}{56}\RG{139.404}{14.183}{10.9092}{ex}{60}\RG{149.101}{9.876}{10.9091}{rm}{49}\RG{154.556}{9.876}{10.9091}{mi}{59}\RG{168.476}{9.876}{10.9091}{mi}{105}\RG{174.678}{9.876}{10.9091}{rm}{43}\RG{185.588}{9.876}{10.9091}{mi}{106}\RG{193.735}{9.876}{10.9091}{rm}{61}\RG{205.250}{9.876}{10.9091}{mi}{109}\RG{217.253}{9.876}{10.9091}{rm}{43}\RG{228.151}{9.876}{10.9091}{rm}{49}\RG{233.616}{9.876}{10.9091}{mi}{59}\RG{0.000}{1.637}{10.9091}{mi}{82}\RG{19.302}{1.637}{10.9091}{rm}{61}\RG{30.800}{1.637}{10.9091}{rm}{40}\RG{35.059}{1.637}{10.9091}{mi}{114}\RG{47.246}{1.637}{10.9091}{rm}{41}\RG{86.005}{1.637}{10.9091}{mi}{59}\RG{112.656}{1.637}{10.9091}{mi}{114}\RG{127.889}{1.637}{10.9091}{rm}{61}\RG{8.283}{0.001}{7.9701}{mi}{109}\RG{39.981}{0.001}{7.9701}{mi}{105}\RG{42.859}{0.001}{7.9701}{mi}{106}\RG{51.489}{0.001}{7.9701}{rm}{49}\RG{55.723}{0.001}{7.9701}{sy}{20}\RG{62.311}{0.001}{7.9701}{mi}{105}\RG{65.188}{0.001}{7.9701}{mi}{59}\RG{67.539}{0.001}{7.9701}{mi}{106}\RG{71.430}{0.001}{7.9701}{sy}{20}\RG{78.017}{0.001}{7.9701}{mi}{109}\RG{117.594}{0.001}{7.9701}{mi}{105}\RG{120.471}{0.001}{7.9701}{mi}{106}\RG{139.404}{-5.454}{10.9092}{ex}{58}\RG{149.101}{-8.496}{10.9091}{rm}{48}\RG{154.556}{-8.497}{10.9091}{mi}{59}\RG{168.476}{-8.497}{10.9091}{mi}{105}\RG{174.678}{-8.496}{10.9091}{rm}{43}\RG{185.588}{-8.497}{10.9091}{mi}{106}\RG{193.735}{-8.496}{10.9091}{sy}{54}\RG{193.735}{-8.497}{10.9091}{rm}{61}\RG{205.250}{-8.496}{10.9091}{mi}{109}\RG{217.253}{-8.497}{10.9091}{rm}{43}\RG{228.151}{-8.497}{10.9091}{rm}{49}\RG{233.616}{-8.496}{10.9091}{mi}{58}%
}
\]
它将基向量的次序反转，因此

\[
\RB{257.767}{11.430}{3.276}{%
\RT{8.367}{4.506}{7.9701}{T}\RG{82.124}{4.506}{7.9701}{rm}{50}\RT{248.887}{4.506}{7.9701}{T}\RG{0.000}{0.001}{10.9091}{mi}{82}\RG{19.302}{0.001}{10.9091}{rm}{61}\RG{30.817}{0.001}{10.9091}{mi}{82}\RG{47.089}{0.001}{10.9091}{mi}{59}\RG{73.741}{0.001}{10.9091}{mi}{82}\RG{93.059}{0.001}{10.9091}{rm}{61}\RG{104.574}{0.001}{10.9091}{mi}{73}\RG{117.358}{0.001}{10.9091}{mi}{59}\RG{144.009}{0.001}{10.9091}{mi}{82}\RG{160.296}{0.001}{10.9091}{mi}{74}\RG{174.334}{0.001}{10.9091}{rm}{40}\RG{178.577}{0.001}{10.9091}{mi}{21}\RG{184.941}{0.001}{10.9091}{rm}{41}\RG{189.184}{0.001}{10.9091}{mi}{82}\RG{208.486}{0.001}{10.9091}{rm}{61}\RG{220.001}{0.001}{10.9091}{mi}{74}\RG{234.039}{0.001}{10.9091}{rm}{40}\RG{238.281}{0.001}{10.9091}{mi}{21}\RG{244.645}{0.001}{10.9091}{rm}{41}\RG{255.172}{0.001}{10.9091}{mi}{58}\RG{39.100}{-1.635}{7.9701}{mi}{109}\RG{109.369}{-1.635}{7.9701}{mi}{109}\RG{152.308}{-1.635}{7.9701}{mi}{109}\RG{166.345}{-1.635}{7.9701}{mi}{109}\RG{197.467}{-1.635}{7.9701}{mi}{109}\RG{226.050}{-1.635}{7.9701}{mi}{109}\RG{8.283}{-2.696}{7.9701}{mi}{109}\RG{82.040}{-2.696}{7.9701}{mi}{109}%
}
\]
最后一个等式也可直接看出：左右同时反转行、列，把Jordan块超对角线上的\allowbreak{}\hskip.1em \RB{5.071}{7.768}{0.500}{%
\RG{0.000}{0.002}{10.9091}{rm}{49}%
}\hskip.1em\allowbreak{} 移到次对角线上，而主对角线不变。

令\allowbreak{}\hskip.1em \RB{19.763}{7.951}{0.740}{%
\RG{0.000}{0.000}{10.9091}{mi}{82}\RG{11.397}{0.000}{10.9091}{rm}{61}%
}\hskip.1em\allowbreak{} diag\allowbreak{}\hskip.1em \RB{72.824}{8.682}{3.243}{%
\RG{0.000}{0.000}{10.9091}{rm}{40}\RG{4.242}{0.000}{10.9091}{mi}{82}\RG{12.525}{-1.636}{7.9701}{mi}{109}\RG{20.016}{-2.743}{5.9776}{rm}{49}\RG{24.665}{0.000}{10.9091}{mi}{59}\RG{29.498}{0.000}{10.9091}{mi}{58}\RG{34.342}{0.000}{10.9091}{mi}{58}\RG{39.174}{0.000}{10.9091}{mi}{58}\RG{44.018}{0.000}{10.9091}{mi}{59}\RG{48.862}{0.000}{10.9091}{mi}{82}\RG{57.190}{-1.636}{7.9701}{mi}{109}\RG{64.682}{-2.632}{5.9776}{mi}{115}\RG{69.172}{0.000}{10.9091}{rm}{41}%
}\hskip.1em\allowbreak{} ，则

\[
\RB{152.820}{11.429}{2.604}{%
\RT{8.367}{4.505}{7.9701}{T}\RG{60.478}{4.505}{7.9701}{sy}{0}\RG{67.064}{4.505}{7.9701}{rm}{49}\RT{143.941}{4.505}{7.9701}{T}\RG{0.000}{0.001}{10.9091}{mi}{82}\RG{17.682}{0.001}{10.9091}{rm}{61}\RG{29.197}{0.001}{10.9091}{mi}{82}\RG{40.594}{0.001}{10.9091}{rm}{61}\RG{52.110}{0.001}{10.9091}{mi}{82}\RG{71.797}{0.001}{10.9091}{mi}{59}\RG{98.448}{0.001}{10.9091}{mi}{82}\RG{106.815}{0.001}{10.9091}{mi}{74}\RG{113.906}{0.001}{10.9091}{mi}{82}\RG{125.327}{0.001}{10.9091}{rm}{61}\RG{136.842}{0.001}{10.9091}{mi}{74}\RG{150.225}{0.001}{10.9091}{mi}{58}%
}
\]
现在构造

\[
\RB{60.842}{11.429}{2.615}{%
\RT{48.575}{4.505}{7.9701}{T}\RG{0.000}{0.001}{10.9091}{mi}{81}\RG{11.655}{0.001}{10.9091}{rm}{61}\RG{23.170}{0.001}{10.9091}{mi}{80}\RG{31.679}{0.001}{10.9091}{mi}{82}\RG{40.047}{0.001}{10.9091}{mi}{80}\RG{58.247}{0.001}{10.9091}{mi}{58}%
}
\]
一方面，

\[
\RB{292.220}{11.432}{3.223}{%
\RT{8.624}{4.508}{7.9701}{T}\RT{46.340}{4.508}{7.9701}{T}\RT{61.145}{4.508}{7.9701}{T}\RT{107.380}{4.508}{7.9701}{T}\RG{237.617}{4.508}{7.9701}{rm}{50}\RG{0.000}{0.004}{10.9091}{mi}{81}\RG{17.939}{0.004}{10.9091}{rm}{61}\RG{29.454}{0.004}{10.9091}{mi}{80}\RG{37.963}{0.004}{10.9091}{mi}{82}\RG{52.625}{0.004}{10.9091}{mi}{80}\RG{70.459}{0.004}{10.9091}{rm}{61}\RG{81.975}{0.004}{10.9091}{mi}{80}\RG{90.484}{0.004}{10.9091}{mi}{82}\RG{98.851}{0.004}{10.9091}{mi}{80}\RG{116.695}{0.004}{10.9091}{rm}{61}\RG{128.210}{0.004}{10.9091}{mi}{81}\RG{136.828}{0.004}{10.9091}{mi}{59}\RT{163.501}{0.004}{10.9091}{d}\RT{169.567}{0.004}{10.9091}{e}\RT{174.411}{0.004}{10.9091}{t}\RG{180.471}{0.004}{10.9091}{mi}{81}\RG{192.127}{0.004}{10.9091}{rm}{61}\RG{203.625}{0.004}{10.9091}{rm}{40}\RT{207.884}{0.004}{10.9091}{d}\RT{213.949}{0.004}{10.9091}{e}\RT{218.793}{0.004}{10.9091}{t}\RG{224.855}{0.004}{10.9091}{mi}{80}\RG{233.374}{0.004}{10.9091}{rm}{41}\RT{244.167}{0.004}{10.9091}{d}\RT{250.233}{0.004}{10.9091}{e}\RT{255.076}{0.004}{10.9091}{t}\RG{261.138}{0.004}{10.9091}{mi}{82}\RG{272.536}{0.004}{10.9091}{sy}{54}\RG{272.536}{0.004}{10.9091}{rm}{61}\RG{284.035}{0.004}{10.9091}{rm}{48}\RG{289.505}{0.004}{10.9091}{mi}{59}%
}
\]
所以\allowbreak{}\hskip.1em \RB{8.573}{8.188}{2.619}{%
\RG{0.000}{-0.003}{10.9091}{mi}{81}%
}\hskip.1em\allowbreak{} 是非异复对称阵。另一方面，记\allowbreak{}\hskip.1em \RB{14.893}{7.948}{0.503}{%
\RG{0.000}{-0.003}{10.9091}{mi}{80}\RG{8.519}{3.956}{7.9701}{sy}{0}%
}\hskip.1em\allowbreak{} T\allowbreak{}\hskip.1em \RB{39.249}{9.756}{3.230}{%
\RG{0.000}{-0.003}{10.9091}{rm}{61}\RG{11.498}{-0.003}{10.9091}{rm}{40}\RG{15.758}{-0.003}{10.9091}{mi}{80}\RG{24.277}{3.956}{7.9701}{sy}{0}\RG{30.863}{3.956}{7.9701}{rm}{49}\RG{35.596}{-0.003}{10.9091}{rm}{41}%
}\hskip.1em\allowbreak{} T，有

\[
\RB{176.523}{11.426}{58.129}{%
\RG{8.624}{4.502}{7.9701}{sy}{0}\RG{15.211}{4.502}{7.9701}{rm}{49}\RG{59.814}{4.502}{7.9701}{sy}{0}\RT{66.401}{4.502}{7.9701}{T}\RG{89.572}{4.502}{7.9701}{sy}{0}\RG{96.158}{4.502}{7.9701}{rm}{49}\RG{129.270}{4.502}{7.9701}{sy}{0}\RG{135.856}{4.502}{7.9701}{rm}{49}\RT{170.236}{4.502}{7.9701}{T}\RG{0.000}{-0.003}{10.9091}{mi}{81}\RG{19.943}{-0.003}{10.9091}{mi}{65}\RG{28.125}{-0.003}{10.9091}{mi}{81}\RG{39.780}{-0.003}{10.9091}{rm}{61}\RG{51.295}{-0.003}{10.9091}{mi}{80}\RG{72.686}{-0.003}{10.9091}{mi}{82}\RG{81.042}{-0.003}{10.9091}{mi}{80}\RG{100.891}{-0.003}{10.9091}{rm}{40}\RG{105.133}{-0.003}{10.9091}{mi}{80}\RG{113.642}{-0.003}{10.9091}{mi}{74}\RG{120.744}{-0.003}{10.9091}{mi}{80}\RG{140.588}{-0.003}{10.9091}{rm}{41}\RG{144.831}{-0.003}{10.9091}{mi}{80}\RG{153.340}{-0.003}{10.9091}{mi}{82}\RG{161.708}{-0.003}{10.9091}{mi}{80}\RG{59.814}{-13.798}{7.9701}{sy}{0}\RT{66.401}{-13.798}{7.9701}{T}\RT{105.038}{-13.798}{7.9701}{T}\RG{39.780}{-18.302}{10.9091}{rm}{61}\RG{51.295}{-18.302}{10.9091}{mi}{80}\RG{72.686}{-18.302}{10.9091}{mi}{82}\RG{81.042}{-18.302}{10.9091}{mi}{74}\RG{88.133}{-18.302}{10.9091}{mi}{82}\RG{96.500}{-18.302}{10.9091}{mi}{80}\RG{59.814}{-32.097}{7.9701}{sy}{0}\RT{66.401}{-32.097}{7.9701}{T}\RT{79.784}{-32.097}{7.9701}{T}\RT{94.587}{-32.097}{7.9701}{T}\RG{39.780}{-36.602}{10.9091}{rm}{61}\RG{51.295}{-36.602}{10.9091}{mi}{80}\RG{72.686}{-36.602}{10.9091}{mi}{74}\RG{86.068}{-36.602}{10.9091}{mi}{80}\RG{79.674}{-50.397}{7.9701}{sy}{0}\RG{86.261}{-50.397}{7.9701}{rm}{49}\RT{95.236}{-50.397}{7.9701}{T}\RT{124.247}{-50.397}{7.9701}{T}\RG{39.780}{-54.901}{10.9091}{rm}{61}\RG{51.279}{-54.901}{10.9091}{rm}{40}\RG{55.537}{-54.901}{10.9091}{mi}{80}\RG{64.047}{-54.901}{10.9091}{mi}{74}\RG{71.149}{-54.901}{10.9091}{mi}{80}\RG{90.993}{-54.901}{10.9091}{rm}{41}\RG{104.550}{-54.901}{10.9091}{rm}{61}\RG{116.066}{-54.901}{10.9091}{mi}{65}\RG{130.532}{-54.901}{10.9091}{mi}{58}%
}
\]
这就是所求的相似关系。\allowbreak{}\hskip.1em \RB{8.366}{7.999}{0.500}{%
\RG{0.000}{0.004}{10.9091}{am}{3}%
}\hskip.1em\allowbreak{}

小结：先在每个Jordan块上用反单位矩阵实现转置，再把这些矩阵拼接起来。回到原坐标时取\allowbreak{}\hskip.1em \RB{23.424}{7.303}{0.721}{%
\RG{0.000}{-0.002}{10.9091}{mi}{80}\RG{7.771}{-0.002}{10.9091}{mi}{82}\RG{15.412}{-0.002}{10.9091}{mi}{80}%
}\hskip.1em\allowbreak{} T；这个合同形式保证对称性，而\allowbreak{}\hskip.1em \RB{33.227}{8.745}{0.720}{%
\RG{0.000}{-0.001}{10.9091}{mi}{80}\RG{7.771}{-0.001}{10.9091}{mi}{82}\RG{15.412}{-0.001}{10.9091}{mi}{80}\RG{23.202}{3.614}{7.9701}{sy}{0}\RG{29.428}{3.614}{7.9701}{rm}{49}%
}\hskip.1em\allowbreak{} 一般不能保证。本题的“对称”不是Hermite对称。

本组题的联系2015S03、S04、S06都通过迭代向量寻找不变子空间；2015S02、S05通过特殊结构构造特征向量；2015S07则利用Jordan块逐块构造满足附加条件的相似矩阵。


\end{document}
